\chapter{The Pauli basis test}
\label{chap:qpbt-test}
% Paper label sec:pauli-verifier kept for cross-reference fidelity; it resolves to this chapter.
\label{sec:pauli-verifier}

This chapter defines the two testing primitives underlying the
question-reduction step: the \emph{classical low individual degree game}, in
its simultaneous form, and the \emph{Pauli basis test}, a nonlocal game built
from the classical game together with the Magic Square game. We state the
quantum soundness of the classical game (Theorem~\ref{lem:ld-soundness},
imported, with the facts required by the identification, which remain to be
proved, recorded in \cite{gap:qpbt_ld-dimension-divisibility}; see
Remark~\ref{rem:ld-soundness-provider}), the rigidity of the Magic Square game
(Theorem~\ref{thm:ms-rigidity}, proved with the consistency
hypothesis on the two distinguished variable measurements that the source
omits --- the source's unrestricted statement is false; the counterexamples
and the correction are recorded in
\cite{gap:qpbt_ms-rigidity-symmetric-strategies}; see
Remark~\ref{rem:ms-rigidity-strategy-class}), completeness of the Pauli basis
test (Lemma~\ref{lem:pauli-completeness}), and the main soundness theorem for
the Pauli basis test (Theorem~\ref{thm:pauli}). The latter is proved in
Chapter~\ref{chap:qpbt-extraction} from soundness of the directly indexed
low-degree game, the global polynomial-pair construction, extraction of the
EPR state and Pauli observables, Naimark dilation, and the transfer from
completed Pauli measurements to the prescribed-answer effects. The unresolved
seed-indexed argument printed in the source remains recorded in
\cite{gap:qpbt_ld-dimension-divisibility}; it supplies no hypothesis of
Theorem~\ref{thm:pauli}. The chapter closes with the qubit form of the soundness
statement (Corollary~\ref{cor:pauli-binary}) and the canonical parameter choice
with its soundness bound (Definition~\ref{def:introparams},
Lemma~\ref{lem:delta-bound}).

Throughout, the two players are denoted $\mathrm A$ and $\mathrm B$; for $w \in \{\mathrm A,\mathrm B\}$ we write $\overline w$ for the other player. Games, strategies, values, conditionally linear (CL) functions and distributions, typed distributions sampled from the edges of a graph, and the state-dependent approximation $\approx_\delta$ between families of measurement operators are as in Chapter~\ref{chap:qpbt-games}; the algebraic material on lines, canonical linear maps, the low-degree encoding, and generalized Pauli observables is from Chapter~\ref{chap:qpbt-algebra}.

\section{The classical low individual degree game}
% Paper labels sec:ld-verifier and sec:ld-game kept for cross-reference fidelity.
\label{sec:ld-verifier}
\label{sec:ld-game}

The classical low individual degree game is a two-player game based on the low individual degree variant of the multilinearity test of Babai, Fortnow, and Lund. In the basic version, one player (the ``points player'') receives a uniformly random point $u \in \F_q^m$ and answers with a value $a \in \F_q$, while the other player (the ``lines player'') receives a random line through $u$ --- either axis-parallel or diagonal --- and answers with a univariate polynomial describing the restriction of a global function to that line. If the players share a polynomial $g\colon \F_q^m \to \F_q$ of individual degree at most $d$ and answer with $g(u)$ and with the restriction $g|_\ell$, they win with probability $1$; the soundness theorem below states that the converse approximately holds. In the simultaneous version, the players' answers are $k$-tuples, checked coordinatewise against a $k$-tuple of low individual degree polynomials.

\begin{definition}[The classical low individual degree game: parameters and registers]\label{def:ld-game}
	\lean{MIPStarRE.QPBT.ldGame}
	\leanok
	\uses{def:game, def:admissible-size}
	% The source fixes \N to be the set of positive integers
	% (references/qpbt-paper/04_preliminaries.tex:6) and takes m, d, ldc in \N
	% (references/qpbt-paper/08_classical_and_quantum_low_degree_tests.tex:33-35),
	% so the positivity below is the source's implicit domain, stated explicitly.
	% It is load-bearing: at d = 0 or k = 0 the prefactor a(dmk)^a of delta_ld in
	% lem:ld-soundness vanishes, and at d = 0 the resulting exact conclusions fail
	% (e.g. at (q,m,d,k) = (2,1,0,1) with eps = 1). The case k = 0, though it would
	% hold trivially, is likewise outside the source's domain and is omitted.
	The game $\mathfrak{G}^{\mathrm{ld}}$ is parametrized by a tuple $\mathsf{ldparams} = (q,m,d,k)$, where $m, d, k \ge 1$ are integers, $q \in \N$ is an admissible field size (Definition~\ref{def:admissible-size}), and $m$ divides $q$. We write $\mathfrak{G}^{\mathrm{ld}}_{\mathsf{ldparams}}$ to emphasize the dependence on the parameter tuple. The type set of $\mathfrak{G}^{\mathrm{ld}}$ is
	\[
		\mathcal T^{\mathrm{ld}} = \{\mathsf{Point}, \mathsf{ALine}, \mathsf{DLine}\}\;,
	\]
	and the distribution over type pairs $(t_{\mathrm A}, t_{\mathrm B})$ is uniform over $\mathcal T^{\mathrm{ld}} \times \mathcal T^{\mathrm{ld}}$. The ambient space is the direct sum $V = V_X \oplus V_I \oplus V_V$ of three register subspaces: the \emph{point register} $V_X \cong \F_q^m$, the \emph{coordinate register} $V_I \cong \F_q$, and the \emph{direction register} $V_V \cong \F_q^m$. Elements of $V$ are identified with triples $(u,s,v) \in \F_q^m \times \F_q \times \F_q^m$. The question distribution and the win predicate of $\mathfrak{G}^{\mathrm{ld}}$ are specified in Definitions~\ref{def:ld-question-distribution} and~\ref{def:ld-win-predicate} below.
\end{definition}

\begin{definition}[Question distribution of the classical low individual degree game]\label{def:ld-question-distribution}
	\uses{def:ld-game, def:cl-func, def:cl-dist, lem:cl-concat, def:cl-canonical, def:line-representative}
	Let $\mathsf{ldparams} = (q,m,d,k)$ be as in Definition~\ref{def:ld-game}. Define $L_{\mathsf{Point}}$ to be the $1$-level CL function on $V$ that projects onto $V_X$: for every $(u,s,v) \in V$,
	\begin{equation}
		\label{eq:cl-ptf}
		L_{\mathsf{Point}}(u,s,v) = (u,0,0)\;.
	\end{equation}
	For $s \in \F_q$, let $\chi(s)$ be the unique integer such that
	\begin{equation}
		\label{eq:chi-func}
		s = (\chi(s)-1)\,\frac{q}{m} + r
	\end{equation}
	for some $0 \le r < q/m$, where $s$ is interpreted as an integer in $\{0,1,\ldots,q-1\}$ under the identification of $\F_q$ with binary strings fixed in Chapter~\ref{chap:qpbt-algebra}. Since $m$ divides $q$, when $s$ is uniformly distributed on $\F_q$ the index $\chi(s)$ is uniformly distributed on $\{1,2,\ldots,m\}$. Define $L_{\mathsf{ALine}}$ by
	\begin{equation}
		\label{eq:cl-alnf}
		L_{\mathsf{ALine}}(u,s,v) = \bigl(L^{\mathrm{Ln}}_{e_i}(u),\, s,\, 0\bigr)\quad\text{for all } (u,s,v) \in V\;,
	\end{equation}
	where $i = \chi(s)$ and $L^{\mathrm{Ln}}_v$ denotes the canonical linear map with one-dimensional kernel spanned by $v$, used to compute canonical representatives of lines (Definition~\ref{def:line-representative}). The function $L_{\mathsf{ALine}}$ is the concatenation (in the sense of Lemma~\ref{lem:cl-concat}) of the $1$-level CL function projecting $V_I \oplus V_V$ onto $V_I$ with the family of $1$-level CL functions $\{L^{\mathrm{Ln}}_{e_i}\}$ on $V_X$ indexed by $i \in \{1,\ldots,m\}$; thus $L_{\mathsf{ALine}}$ is a $2$-level CL function. Define $L_{\mathsf{DLine}}$ by
	\begin{equation}
		\label{eq:cl-dlnf}
		L_{\mathsf{DLine}}(u,s,v) = \bigl(L^{\mathrm{Ln}}_{\pi_{i-1}(v)}(u),\, s,\, \pi_{i-1}(v)\bigr)\;,
	\end{equation}
	where $i = \chi(s)$ and $\pi_{i-1}$ zeroes out the first $i-1$ coordinates of its input. The function $L_{\mathsf{DLine}}$ is the concatenation of (i) the identity on $V_I$, (ii) for each $i$ the $1$-level CL function $v \mapsto \pi_{i-1}(v)$ on $V_V$, and (iii) for each pair $(s,v') \in V_I \oplus V_V$ the $1$-level CL function $L^{\mathrm{Ln}}_{v'}$ on $V_X$; thus $L_{\mathsf{DLine}}$ is a $3$-level CL function.

	The question distribution of $\mathfrak{G}^{\mathrm{ld}}$ is the typed CL distribution obtained as follows: sample a type pair $(t_{\mathrm A}, t_{\mathrm B})$ uniformly from $\mathcal T^{\mathrm{ld}} \times \mathcal T^{\mathrm{ld}}$, sample $z = (u,s,v)$ uniformly from $V$, and send each player $w \in \{\mathrm A,\mathrm B\}$ the question $(t_w, L_{t_w}(z))$.
\end{definition}

\begin{lemma}[Level of the point map]\label{lem:ld-point-level}
	\lean{MIPStarRE.QPBT.isCondLinear_ldPointCL}
	\leanok
	\uses{def:ld-question-distribution, def:cl-func}
	Fix $\mathsf{ldparams} = (q,m,d,k)$ as in Definition~\ref{def:ld-game}. The map $L_{\mathsf{Point}}$ of Equation~\eqref{eq:cl-ptf} is a $1$-level conditionally linear function on $V$.
\end{lemma}
\begin{proof}
	\leanok
	\uses{def:ld-question-distribution, def:cl-func}
	The map $L_{\mathsf{Point}}$ is the linear projection onto the point register $V_X$, hence is a $1$-level CL function.
\end{proof}

\begin{lemma}[Level of the axis-parallel line map]\label{lem:ld-aline-level}
	\lean{MIPStarRE.QPBT.isCondLinear_ldALineCL}
	\leanok
	\uses{def:ld-question-distribution, def:cl-func}
	Fix $\mathsf{ldparams} = (q,m,d,k)$ as in Definition~\ref{def:ld-game}. The map $L_{\mathsf{ALine}}$ of Equation~\eqref{eq:cl-alnf} is a $2$-level conditionally linear function on $V$.
\end{lemma}
\begin{proof}
	\leanok
	\uses{def:cl-func, def:register-subspace, def:ld-question-distribution, def:line-representative}
	The following data exhibit $L_{\mathsf{ALine}}$ as a $2$-level conditionally linear function. Its first level is the projection of $V$ onto the coordinate register $V_I$, which is linear. For a first-level value with coordinate component $s$, its second level is the linear map $L^{\mathrm{Ln}}_{e_i}$ with $i = \chi(s)$, acting on the point register $V_X$. The registers $V_I$ and $V_X$ are disjoint, so the second level is supported on a register subspace disjoint from the first, and summing the two contributions gives $\bigl(L^{\mathrm{Ln}}_{e_i}(u),\, s,\, 0\bigr)$, which is $L_{\mathsf{ALine}}(u,s,v)$ by Equation~\eqref{eq:cl-alnf}.
\end{proof}

\begin{lemma}[Level of the diagonal line map]\label{lem:ld-dline-level}
	\lean{MIPStarRE.QPBT.isCondLinear_ldDLineCL}
	\leanok
	\uses{def:ld-question-distribution, def:cl-func}
	Fix $\mathsf{ldparams} = (q,m,d,k)$ as in Definition~\ref{def:ld-game}. The map $L_{\mathsf{DLine}}$ of Equation~\eqref{eq:cl-dlnf} is a $3$-level conditionally linear function on $V$.
\end{lemma}
\begin{proof}
	\leanok
	\uses{def:cl-func, def:register-subspace, def:ld-question-distribution, def:line-representative}
	The following data exhibit $L_{\mathsf{DLine}}$ as a $3$-level conditionally linear function. Its first level is the projection of $V$ onto the coordinate register $V_I$. For a first-level value with coordinate component $s$, its second level is the linear map $v \mapsto \pi_{i-1}(v)$ with $i = \chi(s)$ on the direction register $V_V$. For a second-level value with direction component $v'$, its third level is the linear map $L^{\mathrm{Ln}}_{v'}$ on the point register $V_X$. The registers $V_I$, $V_V$ and $V_X$ are pairwise disjoint, so each level is supported on a register subspace disjoint from the preceding ones, and summing the three contributions gives $\bigl(L^{\mathrm{Ln}}_{\pi_{i-1}(v)}(u),\, s,\, \pi_{i-1}(v)\bigr)$, which is $L_{\mathsf{DLine}}(u,s,v)$ by Equation~\eqref{eq:cl-dlnf}.
\end{proof}

\begin{lemma}[Formalization support for the low-degree typed CL family]\label{lem:ld-question-cl-family}
	\lean{MIPStarRE.QPBT.isTypedCondLinearFamily_ldCL}
	\leanok
	\uses{def:ld-question-distribution, def:typed-cl-functions}
	This is not a separately named statement of the source article; it isolates the family condition in Lemma~\ref{lem:ld-question-typed-cl}. For every low-degree type $t \in \mathcal T^{\mathrm{ld}}$, the map $L_t$ of Definition~\ref{def:ld-question-distribution} is a $3$-level CL function. Thus $L = \{L_t\}_{t \in \mathcal T^{\mathrm{ld}}}$ is a finite typed family of $3$-level CL functions.
\end{lemma}
\begin{proof}
	\leanok
	\uses{lem:cl-level-mono, lem:ld-point-level, lem:ld-aline-level, lem:ld-dline-level}
	The point, axis-line, and diagonal-line maps have levels $1$, $2$, and $3$, respectively. Raise the first two representations to level $3$ by Lemma~\ref{lem:cl-level-mono}.
\end{proof}

\begin{lemma}[Formalization support for the low-degree typed distribution equality]\label{lem:ld-question-typed-equality}
	\lean{MIPStarRE.QPBT.ldQuestionDistribution_eq_typedCL}
	\leanok
	\uses{def:ld-question-distribution, def:typed-cl-distributions, def:graph-distribution}
	This is not a separately named statement of the source article; it isolates the distribution equality in Lemma~\ref{lem:ld-question-typed-cl}. Fix $\mathsf{ldparams} = (q,m,d,k)$ as in Definition~\ref{def:ld-game}, let $L = \{L_t\}_{t \in \mathcal T^{\mathrm{ld}}}$ be the family of Definition~\ref{def:ld-question-distribution}, and let $G^{\mathrm{ld}}$ be the complete type graph. The question distribution of $\mathfrak G^{\mathrm{ld}}$ equals the distribution $\mu^{G^{\mathrm{ld}}}_{L,L}$ that the construction of Definition~\ref{def:typed-cl-distributions} produces from $(L,L)$ and $G^{\mathrm{ld}}$. This statement records only the distribution identity; Lemma~\ref{lem:ld-question-cl-family} supplies the common-level typed-family assertion.
\end{lemma}
\begin{proof}
	\leanok
	\uses{def:ld-question-distribution, def:typed-cl-distributions, def:graph-distribution, def:cl-dist, lem:distribution-map-comp, lem:uniform-bind-map}
	Every unordered pair of types is an edge of $G^{\mathrm{ld}}$, so the graph distribution $\mu_{G^{\mathrm{ld}}}$ of Definition~\ref{def:graph-distribution} is uniform on $\mathcal T^{\mathrm{ld}} \times \mathcal T^{\mathrm{ld}}$. For each ordered type pair $(t_{\mathrm A}, t_{\mathrm B})$, the CL distribution $\mu_{L_{t_{\mathrm A}},\, L_{t_{\mathrm B}}}$ is the push-forward of the uniform distribution on $V$ along $z \mapsto \bigl(L_{t_{\mathrm A}}(z), L_{t_{\mathrm B}}(z)\bigr)$. Binding the uniform type pair to these push-forwards and tagging each content with its type is exactly the two-stage sampling procedure of Definition~\ref{def:ld-question-distribution}.
\end{proof}

\begin{lemma}[The low-degree question distribution is a typed CL distribution]\label{lem:ld-question-typed-cl}
	\lean{MIPStarRE.QPBT.ldQuestionDistribution_isTypedCL}
	\leanok
	\uses{def:ld-question-distribution, def:typed-cl-distributions}
	Fix $\mathsf{ldparams} = (q,m,d,k)$ as in Definition~\ref{def:ld-game}, let $L = \{L_t\}_{t \in \mathcal T^{\mathrm{ld}}}$ be the family of Definition~\ref{def:ld-question-distribution}, and let $G^{\mathrm{ld}}$ be the graph on $\mathcal T^{\mathrm{ld}}$ all of whose unordered pairs, self-loops included, are edges. The family $L$ is a finite $\mathcal T^{\mathrm{ld}}$-typed family of $3$-level conditionally linear functions, and the question distribution of $\mathfrak{G}^{\mathrm{ld}}$ equals the $(\mathcal T^{\mathrm{ld}}, G^{\mathrm{ld}})$-typed conditionally linear distribution $\mu^{G^{\mathrm{ld}}}_{L,\,L}$ of Definition~\ref{def:typed-cl-distributions}.
\end{lemma}
\begin{proof}
	\leanok
	\uses{lem:ld-question-cl-family, lem:ld-question-typed-equality}
	Combine Lemma~\ref{lem:ld-question-cl-family}, which supplies the typed-family hypothesis, with Lemma~\ref{lem:ld-question-typed-equality}, which identifies the two distributions.
\end{proof}

The next two lemmas show that the CL functions $L_{\mathsf{Point}}$, $L_{\mathsf{ALine}}$, $L_{\mathsf{DLine}}$ give rise to the axis-parallel and diagonal line-versus-point distributions.
We label the coordinates by $r\in\{0,\ldots,m-1\}$, with label $r$
representing the paper's coordinate $r+1$.

\begin{lemma}\label{lem:alnf}
	\lean{MIPStarRE.QPBT.aLinePointDist_point_marginal_uniform, MIPStarRE.QPBT.aLinePointDist_mem_line}
	\leanok
	\uses{def:line-point-dist, def:line}
	Consider the pair $(\ell, u)$ sampled in the following way: first, sample the tuple $\bigl((u_0,s,0),(u,0,0)\bigr)$ from the CL distribution $\mu_{L_{\mathsf{ALine}}, L_{\mathsf{Point}}}$, and let $\ell = \ell(u_0, e_i)$ where $i = \chi(s)$. Then $u$ is a uniformly random point in $\F_q^m$, and $\ell$ is a line that goes through $u$ and is parallel to a uniformly random axis $i \in \{1,2,\ldots,m\}$.
\end{lemma}
\begin{proof}
	\leanok
	\uses{def:cl-dist, def:ld-question-distribution, def:ld-line-description, def:line-representative, lem:distribution-map-comp, lem:uniform-map-equivalence, lem:uniform-product-first-marginal, lem:uniform-product-second-marginal, lem:chi-index-uniform, lem:line-representative-idempotent, lem:line-rep-incidence}
	By Definition~\ref{def:cl-dist} the distribution $\mu_{L_{\mathsf{ALine}}, L_{\mathsf{Point}}}$ is the push-forward of the uniform distribution on $V$ along $z \mapsto \bigl(L_{\mathsf{ALine}}(z), L_{\mathsf{Point}}(z)\bigr)$, so each marginal below is the push-forward of the uniform distribution on $V$ along the composite of that map with the corresponding coordinate. The point marginal is the image of the uniform distribution on $V$ under $(u,s,v) \mapsto u$, hence uniform on $\F_q^m$. The axis marginal is its image under $(u,s,v) \mapsto \chi(s)$; since $m$ divides $q$, every fibre of $\chi$ has exactly $q/m$ elements, so $i = \chi(s)$ is uniform on $\{1,2,\ldots,m\}$ and $(u_0,s)$ specifies a uniformly random axis-parallel line $\ell(u_0,e_i)$. Finally the base point recorded in the sample is $u_0 = L^{\mathrm{Ln}}_{e_i}(u)$, which is fixed by $L^{\mathrm{Ln}}_{e_i}$ (Definition~\ref{def:line-representative}) and satisfies $u \in \ell(u_0,e_i)$ by Lemma~\ref{lem:line-rep-incidence}.
\end{proof}

\begin{lemma}\label{lem:dlnf}
	\lean{MIPStarRE.QPBT.dLinePointDist_point_marginal_uniform, MIPStarRE.QPBT.dLinePointDist_mem_line, MIPStarRE.QPBT.dLinePointDist_prefix_zero}
	\leanok
	\uses{def:line-point-dist, def:line}
	Consider the pair $(\ell, u)$ sampled in the following way: first, sample the tuple $\bigl((u_0,s,v),(u,0,0)\bigr)$ from the CL distribution $\mu_{L_{\mathsf{DLine}}, L_{\mathsf{Point}}}$, and let $\ell = \ell(u_0, v)$ where $i = \chi(s)$. Then $u$ is a uniformly random point in $\F_q^m$, and $\ell$ is a random diagonal line that goes through $u$ and shares its first $i-1$ coordinates with $u$, for uniformly random $i \in \{1,2,\ldots,m\}$.
\end{lemma}
\begin{proof}
	\leanok
	\uses{def:cl-dist, def:ld-question-distribution, def:ld-line-description, def:line-representative, lem:distribution-map-comp, lem:uniform-map-equivalence, lem:uniform-product-first-marginal, lem:uniform-product-second-marginal, lem:chi-index-uniform, lem:prefix-projection-idempotent, lem:line-representative-idempotent, lem:line-description-prefix, lem:line-rep-incidence}
	The three steps are those of the proof of Lemma~\ref{lem:alnf}, and none of them depends on that lemma. The two marginals are again push-forwards of the uniform distribution on $V$ along coordinates of $z \mapsto \bigl(L_{\mathsf{DLine}}(z), L_{\mathsf{Point}}(z)\bigr)$, by Definition~\ref{def:cl-dist}: the point marginal is uniform on $\F_q^m$, and $i = \chi(s)$ is uniform on $\{1,2,\ldots,m\}$ because every fibre of $\chi$ has $q/m$ elements. The direction component of the sampled tuple is $v' = \pi_{i-1}(v)$, whose first $i-1$ coordinates vanish and which is unchanged by a further application of $\pi_{i-1}$. Finally $u_0 = L^{\mathrm{Ln}}_{v'}(u)$ is fixed by $L^{\mathrm{Ln}}_{v'}$ (Definition~\ref{def:line-representative}) and satisfies $u \in \ell(u_0,v')$ by Lemma~\ref{lem:line-rep-incidence}.
\end{proof}

\begin{definition}[Line-point distributions]\label{def:line-point-dist}
	% A line description retains the seed used to select its coordinate.
	\lean{MIPStarRE.QPBT.aLinePointDist, MIPStarRE.QPBT.dLinePointDist}
	\lean{MIPStarRE.QPBT.linePointDist}
	\lean{MIPStarRE.QPBT.aLinePointDist_isProbability}
	\lean{MIPStarRE.QPBT.dLinePointDist_isProbability}
	\lean{MIPStarRE.QPBT.linePointDist_isProbability}
	\leanok
	% The statement-level \leanok records that the six declarations tagged
	% above exist with matching statements, and that the three total-mass
	% assertions among them are proved in Lean with no sorryAx in their
	% axiom closure.
	% The refined laws retain data not present in the axis source law.
	\uses{def:ld-question-distribution, def:ld-line-description}
	The source \emph{axis line-point distribution} $D_{\mathsf{ALine}}$ is
	the law $\mu_{L_{\mathsf{ALine}}, L_{\mathsf{Point}}}$ on pairs
	$(\ell=(u_0,e_i),u)$, and the source \emph{diagonal line-point
	distribution} $D_{\mathsf{DLine}}$ is the corresponding law on pairs
	$(\ell=(u_0,s,v),u)$ from
	$\mu_{L_{\mathsf{DLine}}, L_{\mathsf{Point}}}$. Their equal mixture is
	$D_{\mathsf{Line}}$. We write
	$\widetilde D_{\mathsf{ALine}}$, $\widetilde D_{\mathsf{DLine}}$, and
	$\widetilde D_{\mathsf{Line}}$ for the seed-bearing refinements and their
	equal mixture represented by the linked distribution declarations. An axis
	description records $u_0$ and a seed $s$ selecting $i=\chi(s)$, subject to
	$L^{\mathrm{Ln}}_{e_i}(u_0)=u_0$. A diagonal description records
	$(u_0,s,v)$, subject to $L^{\mathrm{Ln}}_v(u_0)=u_0$ and to the prefix
	condition $v_j=0$ for $j<i$. Define forgetful maps on line-point samples by
	\[
		\rho_{\mathsf{ALine}}((u_0,s),u)
		=((u_0,e_{\chi(s)}),u),
		\qquad
		\rho_{\mathsf{DLine}}((u_0,s,v),u)=((u_0,s,v),u).
	\]
	Then $(\rho_{\mathsf{ALine}})_*\widetilde D_{\mathsf{ALine}}
	=D_{\mathsf{ALine}}$ and
	$(\rho_{\mathsf{DLine}})_*\widetilde D_{\mathsf{DLine}}
	=D_{\mathsf{DLine}}$. The map $\rho_{\mathsf{Line}}$ acting by the
	corresponding map in each component likewise sends
	$\widetilde D_{\mathsf{Line}}$ to $D_{\mathsf{Line}}$.
\end{definition}

\begin{definition}[Win predicate of the classical low individual degree game]\label{def:ld-win-predicate}
	\uses{def:ld-game, def:ld-question-distribution, def:line, def:polyfunc}
	Fix $\mathsf{ldparams} = (q,m,d,k)$. Answers in the game $\mathfrak{G}^{\mathrm{ld}}_{\mathsf{ldparams}}$ take the following forms, depending on the question type; answers not of the prescribed form are rejected.
	\begin{itemize}
		\item Type $\mathsf{Point}$, question content $u \in \F_q^m$: the answer is an element of $\F_q^k$.
		\item Type $\mathsf{ALine}$, question content $(u_0,s) \in \F_q^m \times \F_q$, specifying the axis-parallel line $\ell(u_0,e_i)$ with $i = \chi(s)$: the answer is a $k$-tuple $(f_1,\ldots,f_k)$ of univariate polynomials $f_j\colon \F_q \to \F_q$ of degree at most $d$, each given by its $d+1$ coefficients.
		\item Type $\mathsf{DLine}$, question content $(u_0,s,v) \in \F_q^m \times \F_q \times \F_q^m$, specifying the diagonal line $\ell(u_0,v')$ with $v' = \pi_{i-1}(v)$ and $i = \chi(s)$: the answer is a $k$-tuple of univariate polynomials of degree at most $md$.
	\end{itemize}
	The win predicate $D^{\mathrm{ld}} = D^{\mathrm{ld}}_{\mathsf{ldparams}}$ is the function of the tuple $(t_{\mathrm A}, x_{\mathrm A}, t_{\mathrm B}, x_{\mathrm B}, a_{\mathrm A}, a_{\mathrm B})$ determined by the following clauses, evaluated for $w \in \{\mathrm A,\mathrm B\}$; in all cases where no action is indicated, the predicate accepts.
	\begin{enumerate}
		\item (\textbf{Consistency test}) If $t_{\mathrm A} = t_{\mathrm B}$, accept if and only if $a_{\mathrm A} = a_{\mathrm B}$.
		% The source phrases both clauses with "the" t such that x = u_0 + tv; that phrase
		% does not determine t on diagonal questions with direction v' = 0, which the
		% question distribution samples with positive probability. The universal
		% quantification over t used here is the local correction; see the remark below.
		\item (\textbf{Axis-parallel line-versus-point test}) If $t_w = \mathsf{ALine}$ and $t_{\overline w} = \mathsf{Point}$, accept if and only if $f_j(t) = (a_{\overline w})_j$ for all $j \in \{1,2,\ldots,k\}$ and all $t \in \F_q$ such that $x_{\overline w} = u_0 + t e_i$, where $i = \chi(s)$.
		\item (\textbf{Diagonal line-versus-point test}) If $t_w = \mathsf{DLine}$ and $t_{\overline w} = \mathsf{Point}$, accept if and only if $f_j(t) = (a_{\overline w})_j$ for all $j \in \{1,2,\ldots,k\}$ and all $t \in \F_q$ such that $x_{\overline w} = u_0 + t v'$, where $v' = \pi_{i-1}(v)$ and $i = \chi(s)$.
	\end{enumerate}
	The game $\mathfrak{G}^{\mathrm{ld}}_{\mathsf{ldparams}}$ is the typed game with type set $\mathcal T^{\mathrm{ld}}$, the question distribution of Definition~\ref{def:ld-question-distribution}, and win predicate $D^{\mathrm{ld}}_{\mathsf{ldparams}}$.
\end{definition}

\begin{remark}[Diagonal questions with zero direction]\label{rem:ld-win-zero-direction}
	The source states the two line-versus-point clauses with ``the'' $t \in \F_q$ such that $x_{\overline w} = u_0 + t v$. When the direction is nonzero and $x_{\overline w}$ lies on the line, exactly one such $t$ exists and the formulation above agrees with the source. But zero directions occur: Definition~\ref{def:line} explicitly permits them, and the question distribution of Definition~\ref{def:ld-question-distribution} produces $\mathsf{DLine}$ questions with $v' = \pi_{i-1}(v) = 0$ with positive probability, namely whenever the last $m - i + 1$ coordinates of $v$ vanish. In that case $\ell(u_0, 0) = \{u_0\}$, every $t \in \F_q$ satisfies $x_{\overline w} = u_0 + t v'$ when $x_{\overline w} = u_0$, and the source clause does not determine acceptance for a nonconstant answer polynomial. The universal quantification over $t$ adopted in Definition~\ref{def:ld-win-predicate} makes the predicate a well-defined function of its arguments; it accepts vacuously when no such $t$ exists, a case the question distribution never produces. Completeness is unaffected: the restriction of a polynomial $g$ to the singleton line $\ell(u_0,0)$, expressed through the parameterization $t \mapsto u_0 + t \cdot 0$, is the constant polynomial with value $g(u_0)$, which passes the check for every $t$. Soundness is likewise unaffected: zero directions occur with probability at most $2/(mq)$, so any acceptance convention on this event changes success probabilities by at most that amount; the game-correspondence obligation enumerated in the proof of Theorem~\ref{lem:ld-soundness} accounts for this event explicitly, and an error of this shape is absorbed into the universal constants of $\delta_{\mathrm{ld}}$.
\end{remark}

The soundness properties of the low individual degree game are expressed in terms of the following class of measurements.

\begin{definition}[Low-degree polynomial measurements]\label{def:ld-meas}
	\lean{MIPStarRE.QPBT.PolyIndex, MIPStarRE.QPBT.PolyMeas, MIPStarRE.QPBT.PolyMeasFamily, MIPStarRE.QPBT.PolyMeasTuple}
	\leanok
	\uses{def:polyfunc, def:polymeasurements, def:polynomials-degree}
	Let $\polymeas{m}{q}{d}$ denote the set of POVM measurements indexed by $\polyfunc{m}{q}{d}$, the polynomial representatives of individual degree at most $d$ in each variable; this is the notion of Definition~\ref{def:polymeasurements}, and the argument order $(m,q,d)$ follows that definition (the paper orders the arguments as $(m,d,q)$). By Definition~\ref{def:polynomials-degree}, evaluation identifies these representatives with the corresponding polynomial functions $g\colon \F_q^m \to \F_q$, bijectively whenever $d \le q-1$; at the canonical parameters (Definition~\ref{def:introparams}, where $d = 1 \le q-1$) the identification is bijective; the general statements of this chapter quantify over all admissible tuples and are read through the representative convention. More generally, for an integer $k$ and tuples $m = (m_1,\ldots,m_k)$, $q = (q_1,\ldots,q_k)$, $d = (d_1,\ldots,d_k)$, let $\mathrm{PolyMeas}(m,q,d,k)$ denote the set of measurements $G = \{G_{g_1,\ldots,g_k}\}$ such that for each $i \in \{1,\ldots,k\}$, the outcome index $g_i$ is a polynomial $g_i\colon \F_{q_i}^{m_i} \to \F_{q_i}$ of individual degree at most $d_i$.
	The tuple components are labelled by $r\in\{0,\ldots,k-1\}$, with
	label $r$ representing the paper's component $r+1$.
\end{definition}

Quantum soundness of a version of the classical low-degree test for polynomials of low \emph{total} degree was claimed by Natarajan and Vidick, building on an analysis of a three-prover version of the test by Vidick; there is a known gap in the soundness analysis of that test, and it is not used here. Instead, the construction relies on the low \emph{individual} degree test, whose quantum soundness is proved in~\cite{Ji2020LowIndividualDegree} and generalized in~\cite{Ji2021Quantum}. The statement below is the form most directly useful in what follows.

% The paper states this result in a theorem environment with the label lem:ld-soundness.
% We keep both the environment and the label verbatim for cross-reference fidelity.
\begin{theorem}[Quantum soundness of the simultaneous classical low individual degree test]\label{lem:ld-soundness}
	\lean{MIPStarRE.QPBT.deltaLd, MIPStarRE.QPBT.exists_ld_soundness}
	\leanok
	\uses{def:ld-game, def:ld-question-distribution, def:ld-win-predicate, def:ld-meas, def:projective-strategy-general, def:tensor-product-value, def:consistency, def:bracket}
	There exists a function
	\[
		\delta_{\mathrm{ld}}(\eps, q, m, d, k) = a\,(dmk)^a \bigl(\eps^b + q^{-b} + 2^{-bmd}\bigr)
	\]
	for universal constants $a \ge 1$ and $0 < b \le 1$ such that the following holds. For all $\eps > 0$ and every parameter tuple $\mathsf{ldparams} = (q,m,d,k)$ as in Definition~\ref{def:ld-game} (so $m, d, k \ge 1$, $q$ is an admissible field size, and $m$ divides $q$), and for all projective strategies $(\psi, A, B)$ that succeed with probability at least $1-\eps$ in the game $\mathfrak{G}^{\mathrm{ld}}_{\mathsf{ldparams}}$, there exist measurements
	\[
		G^w \in \mathrm{PolyMeas}(m,q,d,k)
	\]
	on $\hilb_w$ for $w \in \{\mathrm A,\mathrm B\}$, where by slight abuse of notation $m$, $q$, $d$ denote the constant tuples $(m,\ldots,m)$, $(q,\ldots,q)$, $(d,\ldots,d)$ of length $k$, such that
	\begin{align*}
		A^{\mathsf{Point},\, u}_{a_1,\ldots,a_k} \ot I_{\mathrm B}
		&\simeq_{\delta_{\mathrm{ld}}}
		I_{\mathrm A} \ot G^{\mathrm B}_{[(g_1(u),\ldots,g_k(u)) = (a_1,\ldots,a_k)]}\;,\\
		G^{\mathrm A}_{[(g_1(u),\ldots,g_k(u)) = (a_1,\ldots,a_k)]} \ot I_{\mathrm B}
		&\simeq_{\delta_{\mathrm{ld}}}
		I_{\mathrm A} \ot B^{\mathsf{Point},\, u}_{a_1,\ldots,a_k}\;,\\
		G^{\mathrm A}_{g_1,\ldots,g_k} \ot I_{\mathrm B}
		&\simeq_{\delta_{\mathrm{ld}}}
		I_{\mathrm A} \ot G^{\mathrm B}_{g_1,\ldots,g_k}\;,
	\end{align*}
	where $\delta_{\mathrm{ld}} = \delta_{\mathrm{ld}}(\eps,q,m,d,k)$, all three relations are the state-dependent consistency $\simeq$ of Definition~\ref{def:consistency} with respect to the state $\ket{\psi}$, and the bracket notation for post-processed outcomes is that of Definition~\ref{def:bracket}. The first two relations hold on average over the uniform distribution on $u \in \F_q^m$; no question distribution is associated with the third.
\end{theorem}
\begin{proof}
	\uses{def:ld-question-distribution, def:ld-win-predicate, def:ld-meas, def:povm-distance}
	The route transcribed below is the one the paper prints, and it is not the
	route of the formalization. The Lean proof of this statement follows instead
	the seed-indexed direct reduction and seed compression of
	Remark~\ref{rem:ld-soundness-provider}, recorded as
	Theorem~\ref{lem:ld-soundness-formalized}; it establishes all three
	conclusions for every $k$ and uses neither step of the printed route. Those
	two steps --- the identification of $\mathfrak{G}^{\mathrm{ld}}$ with the
	two-prover tensor-code test, and the admissibility of the printed parameter
	choice $m^3 d$ against the bound $12m(d+1)$ required by Theorem~4.7
	of~\cite{Ji2021Quantum} --- remain open source gaps, tracked as issue 527 of
	this repository and enumerated below. This node accordingly carries a
	statement-level formalization mark and no proof-level one, since a
	proof-level mark here would certify the printed route, which nothing in the
	formalization discharges. The following account records that printed route
	and its separate verification obligations.
	Let $\Sigma = \F_q$ and let $\code \subseteq \Sigma^q$ be the Reed--Solomon code of degree $d$: the set of all $(p(x_1),\ldots,p(x_q))$ where $p\colon \F_q \to \F_q$ is a univariate polynomial of degree at most $d$ and $x_1,\ldots,x_q$ enumerate $\F_q$. We may assume $d < q$. In the opposite regime the statement is immediate: for $d \ge q$ the prefactor gives $\delta_{\mathrm{ld}} \ge a\,(dmk)^a\, q^{-b} \ge q^{a-b} \ge 1$, using $m, k \ge 1$, while the quantity bounded in Definition~\ref{def:consistency} never exceeds $4$ for families of POVM elements: $\sum_a \lVert (A_a \ot I - I \ot B_a)\ket{\psi}\rVert^2 \le 2\sum_a \bigl(\bra{\psi} A_a^2 \ot I \ket{\psi} + \bra{\psi} I \ot B_a^2 \ket{\psi}\bigr) \le 4$, using $0 \le M_a^2 \le M_a$ and $\sum_a M_a = \Id$ for POVM elements, so the three relations hold for an arbitrary choice of $G^w \in \mathrm{PolyMeas}(m,q,d,k)$. For $d < q$, evaluation is injective on polynomials of degree at most $d$ and a nonzero polynomial of degree at most $d$ vanishes at no more than $d$ of the $q$ points, so $\code$ is a linear $[q,\, d+1,\, q-d]_\Sigma$ code.
	% The paper prints the parameters as [q, d+1, n-d+1] with an undefined n. With the
	% intended value n = q the printed distance q-d+1 exceeds the true minimum distance
	% q-d of the degree-d evaluation code, and the dimension claim d+1 requires d < q.
	% Both points are corrected above; the regime d >= q, which the parameter tuple of
	% Definition def:ld-game permits, is handled by the triviality argument. The
	% boundary values d = 0 and k = 0 are excluded by the parameter domain of
	% Definition def:ld-game, following the source's convention that \N denotes
	% the positive integers.
	% The source asserts a flat identity between its game and the two-prover
	% tensor-code test and applies Theorem 4.7 of ji2021quantum with the
	% parameter choice m^3 d
	% (references/qpbt-paper/08_classical_and_quantum_low_degree_tests.tex:446-454),
	% without checking that theorem's hypothesis k >= 12mt on the parameter.
	% Inspection of the cited test (arXiv:2111.08131, Figures 1 and 2,
	% Definitions 4.5-4.6, Theorem 4.7) shows the asserted identity is not one:
	% the question formats, the branch weights, and the zero-direction
	% convention all differ, and the parameter bound fails at small m. Per the
	% import policy, this proof states the source's route and enumerates
	% the facts required by the identification, which remain to be proved, as open obligations; the
	% detailed comparison and the case arithmetic are in the import-obligations
	% section of docs/paper-gaps/qpbt_ld-dimension-divisibility.tex.
	In this regime the statement is imported, and the source's route is as follows: identify the game $\mathfrak{G}^{\mathrm{ld}}_{\mathsf{ldparams}}$ with $k = 1$ with the two-prover tensor-code test for the code $\code^{\ot m}$ of~\cite{Ji2021Quantum} (Figures~1 and~2 there, the roles of the two players being assigned uniformly at random); apply Theorem~4.7 of~\cite{Ji2021Quantum}, instantiated with its auxiliary integer parameter --- denoted $k$ there and unrelated to the simultaneity parameter $k$ of $\mathsf{ldparams}$ --- chosen as $m^3 d$, so that its error function $\mathrm{poly}(m, d+1, m^3d) \cdot \mathrm{poly}(\eps, q^{-1}, e^{-md})$ takes the displayed form of $\delta_{\mathrm{ld}}$; and extend from $k = 1$ to general $k$ by a standard reduction, following exactly the derivation of Theorem~4.43 from Theorem~4.40 in~\cite{Natarajan2018NEEXP}. The import carries two verification obligations, which remain open; they are enumerated here as obligations and detailed, together with candidate repairs, in the import-obligations section of~\cite{gap:qpbt_ld-dimension-divisibility}.
	\begin{enumerate}
		\item (\emph{Game correspondence.}) The source asserts that the two games are identical; they are not. In the cited test, the subcube commutation and pair-synchronicity branches send an ordered pair of points of a \emph{hidden} subcube and receive a pair of field values --- a question type that $\mathfrak{G}^{\mathrm{ld}}$ does not have --- its line questions carry no analogue of the seed $s$, its nontrivial cross-tests have mass $1/3$ where the present ones have $2/9$, and it fixes no acceptance convention on the zero-direction event of Remark~\ref{rem:ld-win-zero-direction}, whose probability is at most $\frac1m \sum_{j=1}^m q^{-j} \le \frac{1}{m(q-1)} \le \frac{2}{mq}$ over uniform $(s,v)$. What must be proved is a reduction converting any projective strategy for $\mathfrak{G}^{\mathrm{ld}}$ with success probability $1-\eps$ into a two-prover strategy for the tensor-code test with rejection probability at most $\tfrac32\,\eps + O(1/(mq))$ --- the factor $\tfrac32$ being the target bookkeeping constant from the branch weights, and the $O(1/(mq))$ the zero-direction term. A deficit of this shape is then absorbed into $\delta_{\mathrm{ld}}$ by enlarging the universal constant $a$, since $(x+y)^b \le x^b + y^b$ for $0 < b \le 1$ and $(2/(mq))^b \le 2\,q^{-b}$ for $m \ge 1$, using $dmk \ge 1$.
		\item (\emph{Parameter bound.}) Theorem~4.7 requires its auxiliary parameter to be at least $12mt$, where $t$ is the interpolation parameter of the base code, equal to $d+1$ for the degree-$d$ Reed--Solomon code $\code$ (blocklength $q$, distance $q-d$). The choice $m^3 d$ satisfies $m^3 d \ge 12m(d+1)$ if and only if $m^2 d \ge 12(d+1)$, which holds for all $d \ge 1$ when $m \ge 5$ and, at $m = 4$, exactly when $d \ge 3$, but fails for every $d$ when $m \le 3$. Over the parameter domain of Definition~\ref{def:ld-game}, where $m$ divides the admissible field size $q$ and is therefore a power of $2$, the bound thus fails exactly at $m \in \{1,2\}$ for every $d$ and at $m = 4$ for $d \in \{1,2\}$; the case check and the candidate repair --- instantiating the parameter as $\max(m^3 d,\, 12m(d+1))$ and re-absorbing the error terms --- are in the note.
	\end{enumerate}
	Quantum soundness of the low individual degree test itself is the result of~\cite{Ji2020LowIndividualDegree}; Remark~\ref{rem:ld-soundness-provider} relates the present statement to Theorem~\ref{thm:main-formal}, which proves that result.
\end{proof}

% This node is not a statement from the paper: it records the route by which
% Theorem lem:ld-soundness is formalized, so that a proof-level formalization
% mark certifies the formalized route and not the printed tensor-code one.
\begin{theorem}[Formalization support for Theorem~\ref{lem:ld-soundness}: the seed-indexed route]
	\label{lem:ld-soundness-formalized}
	\lean{MIPStarRE.QPBT.exists_ld_soundness}
	\lean{MIPStarRE.QPBT.polyTupleAgreement_avg_le_mdq}
	\lean{MIPStarRE.QPBT.ten_sqrt_deltaLd_le}
	\lean{MIPStarRE.QPBT.error_and_collision_le_deltaLd}
	\leanok
	\uses{def:ld-game, def:ld-question-distribution, def:ld-win-predicate, def:ld-meas, def:projective-strategy-general, def:tensor-product-value, def:consistency, def:bracket}
	This formalization-only node is not a named statement of the paper; it carries the conclusion of Theorem~\ref{lem:ld-soundness} verbatim, and exists so that the route the formalization actually takes is the route a proof-level formalization mark certifies. There exists a function
	\[
		\delta_{\mathrm{ld}}(\eps, q, m, d, k) = a\,(dmk)^a \bigl(\eps^b + q^{-b} + 2^{-bmd}\bigr)
	\]
	for universal constants $a \ge 1$ and $0 < b \le 1$ such that all three consistency relations displayed in Theorem~\ref{lem:ld-soundness} hold with that bound, with the same question laws and the same bracket convention, for every parameter tuple $\mathsf{ldparams} = (q,m,d,k)$ as in Definition~\ref{def:ld-game}, every $\eps > 0$, and every projective strategy $(\psi, A, B)$ succeeding with probability at least $1-\eps$ in $\mathfrak{G}^{\mathrm{ld}}_{\mathsf{ldparams}}$. The statement is that of Theorem~\ref{lem:ld-soundness}; only the proof recorded below differs from the one printed there. The auxiliary declarations tagged on this node are the tuple collision bound and the two scalar absorption estimates that proof uses.
\end{theorem}
\begin{proof}
	\leanok
	\uses{prop:ld-simultaneous-general-k, lem:direct-ld-residue-compression, lem:ld-point-agreement, lem:direct-ld-polynomial-agreement, fact:triangle-for-simeq, lem:ld-soundness-k-one}
	Neither obligation enumerated in the proof of Theorem~\ref{lem:ld-soundness} is used or assumed here: the two-prover tensor-code test does not occur in this argument at all.

	Dilate the given seed-indexed strategy by a maximally correlated register indexed by the seed fibre, so that each player answers the reconstructed directly indexed question in the block selected by their copy. The dilation preserves projectivity and preserves the value exactly, because the residue is uniform and shared. Proposition~\ref{prop:ld-simultaneous-general-k}, applied to the dilated strategy, supplies polynomial-tuple measurements for every $k$ on the extended Hilbert spaces. Compress both of them back to the original Hilbert spaces by the normalized diagonal-block average over the correlated register. The two mixed point-versus-polynomial relations transfer through that compression exactly, by the two compression-transfer lemmas of Lemma~\ref{lem:direct-ld-residue-compression}: in each of them the opposite player's operator is the block-diagonal amplification of the original point measurement and therefore leaves the correlated register untouched.

	The third relation, in which both operators act on the correlated register, is not preserved by compression and is reproved on the original spaces. The point-versus-point branch of Definition~\ref{def:ld-win-predicate} accepts only on equal answers and carries one ninth of the question law, so the two point measurements are consistent up to $9\eps$ (Lemma~\ref{lem:ld-point-agreement}), and the triangle inequality of Lemma~\ref{fact:triangle-for-simeq} chains the two compressed relations through that agreement. The defect this leaves is the probability that two \emph{distinct} polynomial tuples take the same value at a uniform point, which is at most $md/q$: distinct tuples differ in some coordinate, and the Schwartz--Zippel estimate of Lemma~\ref{lem:direct-ld-polynomial-agreement} is applied at that one coordinate, so no union bound over the $k$ coordinates is taken and the estimate does not degrade in $k$. This bounds collisions between tuples that have already been constructed; it does not construct simultaneous measurements coordinatewise, and so it does not touch the coordinatewise route refuted in Remark~\ref{rem:ld-soundness-simultaneity}.

	The accumulated error $E + 2\sqrt{9\eps + E} + md/q$, where $E$ is the error supplied by Proposition~\ref{prop:ld-simultaneous-general-k}, is absorbed into the displayed shape of $\delta_{\mathrm{ld}}$ by replacing the direct constants $(a,b)$ by $(10a, b/2)$: the square root halves the exponent, and the factor ten is taken up by the prefactor. The case $k = 1$ is recorded separately as Theorem~\ref{lem:ld-soundness-k-one}.
\end{proof}

\begin{remark}[Relation to the low individual degree test of Chapter~\ref{chap:test}]\label{rem:ld-soundness-provider}
	The printed proof of Theorem~\ref{lem:ld-soundness} uses an import: the proof above records the source's reduction to Theorem~4.7 of~\cite{Ji2021Quantum}, instantiated with that theorem's auxiliary parameter chosen as $m^3 d$ and followed by the extension to simultaneous $k$-tuples of~\cite{Natarajan2018NEEXP}, together with the two open facts required by the identification, which remain to be proved, --- the game correspondence and the parameter bound, detailed in~\cite{gap:qpbt_ld-dimension-divisibility}. The result of~\cite{Ji2020LowIndividualDegree} that underlies it is proved here as Theorem~\ref{thm:main-formal}: for every projective strategy passing the $(m,q,d)$ low individual degree test of Definition~\ref{def:low-individual-degree-test} with probability at least $1-\eps$, that theorem produces measurements in $\polymeas{m}{q}{d}$ satisfying consistency relations of the same shape as the three approximations of Theorem~\ref{lem:ld-soundness}. Theorem~\ref{lem:ld-soundness} is not a consequence of Theorem~\ref{thm:main-formal} alone: a derivation from it would require, in addition, an equivalence between the typed game $\mathfrak{G}^{\mathrm{ld}}_{(q,m,d,1)}$ and the $(m,q,d)$ test of Chapter~\ref{chap:test} --- matching the two question distributions, the two win conditions, and the two error functions --- together with the extension from $k = 1$ to general $k$ by the reduction of~\cite{Natarajan2018NEEXP} cited above. Proposition~\ref{prop:ld-simultaneous-general-k} establishes the directly indexed combining reduction for every $k$. Its correlated seed dilation preserves projectivity and value, and the two point/polynomial relations compress exactly. Point agreement and the Schwartz--Zippel bound for distinct polynomial tuples recover the global relation on the original spaces. The error is absorbed by replacing the direct constants $(a,b)$ by $(10a,b/2)$, proving Theorem~\ref{lem:ld-soundness} for every $k$ without invoking either open tensor-code assertion; the case $k=1$ is recorded separately in Theorem~\ref{lem:ld-soundness-k-one}. The printed tensor-code route above does not invoke Theorem~\ref{thm:main-formal}.
\end{remark}

% This remark records the formal status of the extension to general k;
% it is not a statement from the paper.
\begin{remark}[The extension to general $k$ is not coordinatewise]\label{rem:ld-soundness-simultaneity}
	The extension from $k = 1$ to general $k$ cited in the proof of Theorem~\ref{lem:ld-soundness} is the reduction deriving Theorem~4.43 from Theorem~4.40 in \cite{Natarajan2018NEEXP}: the $k$ answer polynomials are combined into the single polynomial $(\alpha, u) \mapsto \sum_i \alpha_i g_i(u)$ in $m + k$ variables, the case $k = 1$ is applied once to the combined strategy, and the $k$ polynomials are recovered from the resulting measurement by exact linearity in $\alpha$ and the Schwartz--Zippel lemma. It cannot be replaced by applying the case $k = 1$ to each coordinate separately and combining the coordinate measurements by the sandwich of Lemma~\ref{lem:ld-sandwich}: over $\F_{2^s}$ with $d = 1$ and $k = 2$ there are a joint point measurement and coordinate polynomial measurements, on the maximally entangled state, for which every coordinate relation of the theorem holds with error $0$ while every POVM with polynomial-pair outcomes has consistency defect at least $1 - 2/q$ against the joint point measurement; the computation is in \cite{gap:qpbt_ld-simultaneous-sandwich}. The formalization accordingly proves the polynomial-tuple conclusion of Theorem~\ref{lem:ld-soundness} for $k = 1$ from Theorem~\ref{thm:main-formal}, which is the only value instantiated in the proof of Lemma~\ref{lem:qld-4-7}; the general case is the combining reduction of Proposition~\ref{prop:ld-simultaneous-general-k}, carried out for the directly indexed game. The correlated seed-fibre dilation and compression yield the seed-indexed theorem for every $k$: distinct tuples differ in some coordinate, so their collision probability at a uniform point is at most $md/q$. This recovers global consistency from the two compressed point relations and the point-agreement branch. The construction of simultaneous measurements still comes from the combining reduction; the coordinate argument here only bounds collisions of already constructed tuples. The case $k = 1$ is recorded separately in Theorem~\ref{lem:ld-soundness-k-one}. The two games and the transport are compared in \cite{gap:qpbt_ld-dimension-divisibility}.
	% This remark carries no Lean metadata, since remarks are not blueprint
	% dependency nodes; the Lean declarations are linked from the
	% formalization support nodes cited above.
\end{remark}

% The nodes of this subsection are formalization-support nodes recording the
% reduction cited in the proof of Theorem lem:ld-soundness; the paper states
% none of them separately.
\subsection{The combining reduction to simultaneity parameter one}
\label{sec:ld-combining-reduction}

The extension from $k = 1$ to general $k$ cited in the proof of
Theorem~\ref{lem:ld-soundness} is recorded here in the form in which it is
formalized.  The nodes below are formalization-support nodes: the source states
none of them separately, and together they constitute its derivation of
Theorem~4.43 from Theorem~4.40 in~\cite{Natarajan2018NEEXP}.  The game in
question is the directly indexed low-degree game of
\cite{gap:qpbt_ld-dimension-divisibility}, which samples the coordinate index of
a line directly instead of through a fiber of $\chi$; its point questions, line
questions, answer formats and win predicate are those of
Definition~\ref{def:ld-game}, and it is used here because the reduction applies
the case $k = 1$ at dimension $m + k$, where the divisibility condition of
Definition~\ref{def:ld-game} is not available.
% A node of this subsection is marked as formalized only when the Lean
% declarations it names are complete and carry no unproved placeholder.

\begin{definition}[Combined parameters of the reduction]
	\label{def:ld-combining-parameters}
	\lean{MIPStarRE.QPBT.DirectLdParams, MIPStarRE.QPBT.DirectLdParams.combined}
	\leanok
	\uses{def:admissible-size}
	For directly indexed parameters $\mathsf{ldparams} = (q,m,d,k)$, where
	$m,d,k \ge 1$ are integers and $q$ is an admissible field size, no condition
	$m \mid q$ is imposed.  The \emph{combined parameter tuple} is
	\[
		\mathsf{ldparams}^{\mathrm c} = (q,\; m+k,\; d,\; 1)\;:
	\]
	the dimension grows by the simultaneity parameter, the simultaneity
	parameter becomes $1$, and the field size, the degree and the admissibility
	of the field size are unchanged.  The degree is unchanged because the
	combining map of Definition~\ref{def:ld-combining-map} is linear in the
	fresh coordinates and therefore, in the individual-degree class of
	Definition~\ref{def:polyfunc}, raises no individual degree above $d$ as soon
	as $d \ge 1$.  The source works with total degree and passes from $d$ to
	$d+1$ at this step; that increment is an artifact of the total-degree
	formulation and does not occur here.
\end{definition}

\begin{definition}[The combining map on polynomial tuples]
	\label{def:ld-combining-map}
	\lean{MIPStarRE.QPBT.combinePolyTuple, MIPStarRE.QPBT.combinePolyTuple_eval}
	\lean{MIPStarRE.QPBT.combinePolyTuple_mem_polyFunc}
	\lean{MIPStarRE.QPBT.directCombinedPolynomial}
	\lean{MIPStarRE.QPBT.directCombinedPolynomial_eval}
	\leanok
	\uses{def:polyfunc, def:ld-meas, def:ld-combining-parameters}
	Coordinates of $\F_q^{m+k}$ are written $(u,\alpha)$, with $u \in \F_q^m$
	occupying the first $m$ coordinates and $\alpha \in \F_q^k$ the last $k$;
	the coordinates carrying $\alpha$ are called the \emph{combining
	coordinates}.  For a tuple $g = (g_0,\ldots,g_{k-1})$ of polynomials
	$g_r \in \polyfunc{m}{q}{d}$, the \emph{combined polynomial} is
	\[
		\mathrm{comb}_g(u,\alpha) \;=\; \sum_{r=0}^{k-1} \alpha_r\, g_r(u)\;.
	\]
	It lies in $\polyfunc{m+k}{q}{d}$ whenever $d \ge 1$: a component $g_r$
	involves no combining coordinate, so multiplying it by $\alpha_r$ leaves its
	individual degrees in the first $m$ coordinates unchanged and produces
	individual degree $1$ in the coordinate carrying $\alpha_r$.

	The order of the two coordinate blocks is part of the construction.  The
	point coordinates precede the combining coordinates, so that zeroing a
	prefix of a direction of $\F_q^{m+k}$, as the diagonal-line branch of
	Definition~\ref{def:ld-question-distribution} does, restricts on the first
	$m$ coordinates to zeroing the corresponding prefix of the point part of
	that direction; this is what makes the diagonal-line questions of the
	combined game project to diagonal-line questions of the original game in
	Definition~\ref{def:ld-combined-strategy}.  Definition~\ref{def:combine-map}
	is a variant of the same construction with two components in \emph{disjoint}
	blocks of point coordinates; the map above keeps one common block of point
	coordinates, as the reduction requires.
\end{definition}

\begin{lemma}[Recovery of the components of a combined polynomial]
	\label{lem:ld-combining-split}
	\lean{MIPStarRE.QPBT.splitCombinedPoly}
	\lean{MIPStarRE.QPBT.splitCombinedPoly_combinePolyTuple}
	\lean{MIPStarRE.QPBT.combinePolyTuple_injective}
	\lean{MIPStarRE.QPBT.IsDirectCombined, MIPStarRE.QPBT.directTupleOfCombined}
	\lean{MIPStarRE.QPBT.directTupleOfCombined_directCombinedPolynomial}
	\leanok
	\uses{def:ld-combining-map}
	Let $\mathrm{split}_r$ denote the substitution of the $r$-th standard basis
	vector of $\F_q^k$ for the combining coordinates, a ring homomorphism
	$\F_q[u,\alpha] \to \F_q[u]$ fixing the point coordinates.  Then
	$\mathrm{split}_r(\mathrm{comb}_g) = g_r$ for every tuple $g$ and every
	$r$.  Consequently $g \mapsto \mathrm{comb}_g$ is injective, and the map
	sending $p \in \polyfunc{m+k}{q}{d}$ to the unique tuple $g$ with
	$\mathrm{comb}_g = p$ when $p$ is of that form, and to the zero tuple
	otherwise, is a well-defined total map $\mathrm{split}$ with
	$\mathrm{split}(\mathrm{comb}_g) = g$.
\end{lemma}
\begin{proof}
	\leanok
	\uses{def:ld-combining-map}
	Applying the homomorphism $\mathrm{split}_r$ to the sum defining
	$\mathrm{comb}_g$ sends the summand of index $s$ to $g_s$ when $s = r$ and to
	$0$ otherwise, since the combining coordinate carrying $\alpha_s$ is sent to
	$1$ for $s = r$ and to $0$ otherwise, while $g_s$, which involves no
	combining coordinate, is fixed; the sum is therefore $g_r$.
	Injectivity follows, and it makes the case distinction defining
	$\mathrm{split}$ unambiguous.  The value on polynomials that are not
	combined is immaterial: the source leaves it arbitrary, and the estimates
	below use only the value on combined polynomials.
\end{proof}

\begin{lemma}[Collision of two linear forms in the combining variables]
	\label{lem:ld-combining-collision}
	\lean{MIPStarRE.QPBT.combiningLinearForm}
	\lean{MIPStarRE.QPBT.combiningLinearForm_injective}
	\lean{MIPStarRE.QPBT.combiningLinearForm_agreementProbability_le}
	\lean{MIPStarRE.QPBT.directCombiningCollision_avg_le}
	\leanok
	\uses{lem:schwartz-zippel-total-degree}
	Let $c, b \in \F_q^k$ with $c \neq b$.  Then
	\[
		\Pr_{\alpha \in \F_q^k}\Bigl[\; \sum_{r} c_r \alpha_r
			= \sum_{r} b_r \alpha_r \;\Bigr] \;\le\; \frac1q\;,
	\]
	the probability being over a uniformly random $\alpha$.
\end{lemma}
\begin{proof}
	\leanok
	\uses{lem:schwartz-zippel-total-degree}
	The two sides are the evaluations at $\alpha$ of the linear forms
	$\sum_r c_r X_r$ and $\sum_r b_r X_r$ in $k$ variables.  The coefficient of
	$X_r$ in such a form is $c_r$, so the two forms are distinct polynomials,
	and each has total degree at most $1$.  Lemma~\ref{lem:schwartz-zippel-total-degree}
	bounds the probability that two distinct polynomials of total degree at most
	$1$ agree at a uniformly random point of $\F_q^k$ by $1/q$.
\end{proof}

\begin{definition}[The combined strategy]
	\label{def:ld-combined-strategy}
	\lean{MIPStarRE.QPBT.directCombinedMeasuredQuestion}
	\lean{MIPStarRE.QPBT.directCombinedMeasuredDirection}
	\lean{MIPStarRE.QPBT.directCombinedRebaseParameter}
	\lean{MIPStarRE.QPBT.directCombinedAnswerMap}
	\lean{MIPStarRE.QPBT.directCombinedStrategy}
	\leanok
	\uses{def:ld-game, def:ld-question-distribution, def:ld-win-predicate,
		def:tensor-product-strategy, def:bracket,
		def:ld-combining-parameters, def:ld-combining-map}
	Let $\mathcal S = (\psi, A, B)$ be a strategy for the directly indexed
	low-degree game with parameters $\mathsf{ldparams}$.  The
	\emph{combined strategy} $\mathcal S^{\mathrm c}$ for the directly indexed
	low-degree game with parameters $\mathsf{ldparams}^{\mathrm c}$ has the same
	state and the same two Hilbert spaces, and answers each question of the
	combined game by measuring one question of $\mathcal S$ and post-processing
	its outcome.  Writing questions of the combined game in the coordinates of
	Definition~\ref{def:ld-combining-map}, the five cases are the following.
	\begin{enumerate}
		\item On the point question $(u,\alpha) \in \F_q^{m+k}$: measure the
		point question $u$, obtaining $b \in \F_q^k$, and answer
		$\sum_r \alpha_r b_r$.
		\item On the axis-parallel line question with stored index $i < m$ and
		canonical base $(u_0,\alpha_0)$: measure the axis-parallel line question
		with index $i$ and base $u_0$, obtaining univariate polynomials
		$f_0,\ldots,f_{k-1}$ of degree at most $d$, and answer
		$\sum_r (\alpha_0)_r f_r$, again of degree at most $d$.
		\item On the axis-parallel line question with stored index $m + r$ and
		canonical base $(u_0,\alpha_0)$, a line along which $u$ is constant and
		only the coordinate $\alpha_r$ varies: measure the point question $u_0$,
		obtaining $b$, and answer the degree-one polynomial
		$t \mapsto \sum_s (\alpha_0)_s b_s + t\, b_r$.
		\item On the diagonal line question with stored index $i < m$, canonical
		base $(u_0,\alpha_0)$ and direction $(v,w)$: measure the diagonal line
		question with index $i$, direction $v$ and the canonical base of the
		line $\ell(u_0, v)$, rebase the answer along the affine change of
		parameter relating the two presentations of that line, obtaining
		univariate polynomials $f_0,\ldots,f_{k-1}$ of degree at most $md$ in the
		parameter of the combined line, and answer
		$t \mapsto \sum_r \bigl((\alpha_0)_r + t\, w_r\bigr) f_r(t)$, of degree
		at most $md+1 \le (m+k)d$.
		\item On the diagonal line question with stored index $i \ge m$,
		canonical base $(u_0,\alpha_0)$ and direction $(0,w)$, a line along which
		$u$ is constant: measure the point question $u_0$, obtaining $b$, and
		answer $t \mapsto \sum_r (\alpha_0)_r b_r + t \sum_r w_r b_r$.
	\end{enumerate}
	For any of these five question forms, if the outcome of the measured
	question has a different answer form from the one specified in the
	corresponding case, the post-processing returns the zero answer of the form
	required by the combined question.  This convention makes the
	post-processing a total map on the common answer alphabet.
	For the point measurements of the original strategy, let $r$ be the map
	that sends a point answer to its tuple in $\F_q^k$ and sends either kind of
	line answer to the zero tuple.  Write
	\[
		\widehat A^{\mathsf{Point},u}_a
		= \sum_{x:\,r(x)=a} A^{\mathsf{Point},u}_x,
		\qquad
		\widehat B^{\mathsf{Point},u}_a
		= \sum_{x:\,r(x)=a} B^{\mathsf{Point},u}_x.
	\]
	These are the point measurements post-processed by
	\texttt{directLdPointValuesOrZero}.  In particular, each zero-tuple effect
	is the original zero-tuple point effect plus all wrong-form answer effects;
	those additional effects need not vanish.  Use the same convention for the
	combined strategy, with its own answer alphabet.

	Cases~3 and~5 use that a line of $\F_q^{m+k}$ whose direction
	has vanishing point part is contained in a fiber of the projection to
	$\F_q^m$, so that the combined polynomial restricted to it is determined by
	the $k$ values of the tuple at the single point $u_0$; the rebasing in
	case~4 is the affine change of parameter of the previous paragraph, needed
	because the
	canonical base of a line of $\F_q^{m+k}$ need not project to the canonical
	base of its image line in $\F_q^m$.
\end{definition}

\begin{theorem}[Projectivity of the combined strategy]
	\label{thm:ld-combined-strategy-projective}
	\lean{MIPStarRE.QPBT.directCombinedStrategy_isProjective}
	\leanok
	\uses{def:ld-combined-strategy, def:projective-strategy-general}
	For every strategy $\mathcal S$ for the directly indexed low-degree game, if
	$\mathcal S$ is projective, then the combined strategy
	$\mathcal S^{\mathrm c}$ is projective.
\end{theorem}
\begin{proof}
	\leanok
	\uses{def:ld-combined-strategy, lem:sandwich-postprocess-projective}
	For either player and every question of the combined game, the corresponding
	measurement of $\mathcal S^{\mathrm c}$ is obtained by relabelling the outcomes
	of one measurement of $\mathcal S$.  Every such post-processing of a
	projective measurement is projective.
\end{proof}

\begin{lemma}[The questions measured by the combined strategy]
	\label{lem:ld-combined-question-law}
	\lean{MIPStarRE.QPBT.directCombined_measuredQuestion_of_pointVar}
	\lean{MIPStarRE.QPBT.directCombined_measuredQuestion_of_coefficientVar}
	\lean{MIPStarRE.QPBT.directCombinedSample_pointParts_uniform}
	\lean{MIPStarRE.QPBT.directCombinedSample_index_uniform}
	\lean{MIPStarRE.QPBT.directCombinedSample_indexPointParts_uniform}
	\leanok
	\uses{def:ld-question-distribution, def:ld-combined-strategy}
	Sample a question pair of the directly indexed game with parameters
	$\mathsf{ldparams}^{\mathrm c}$, and let $\mathbf{i}$ be the coordinate index
	of its common sample, uniform on $\{0,1,\ldots,m+k-1\}$.  Conditioned on
	$\mathbf{i} < m$, an event of probability $m/(m+k)$, the pair of questions
	measured by $\mathcal S^{\mathrm c}$ in Definition~\ref{def:ld-combined-strategy}
	is distributed as the question pair of the directly indexed game with
	parameters $\mathsf{ldparams}$ conditioned on the same ordered type pair.
	Conditioned on $\mathbf{i} \ge m$, both measured questions are the point
	question at the point part of the sampled point, and that point part is
	uniform on $\F_q^m$; that is, the measured pair is a point-versus-point pair
	of the game with parameters $\mathsf{ldparams}$.
\end{lemma}
\begin{proof}
	\leanok
	\uses{def:ld-question-distribution, def:ld-combining-map}
	The common sample of the combined game consists of a uniform point
	$(u,\alpha) \in \F_q^{m+k}$, a uniform index $\mathbf{i}$, and a uniform
	direction $(v,w)$, the three independent.  Conditioned on $\mathbf{i} < m$,
	the triple $(u, \mathbf{i}, v)$ is uniform on
	$\F_q^m \times \{0,\ldots,m-1\} \times \F_q^m$, which is the common sample of
	the game with parameters $\mathsf{ldparams}$, and $\alpha$ and $w$ are
	independent of it.  Under this conditioning the coordinate direction of an
	axis-parallel question and the prefix projection of a diagonal direction have
	point parts equal to the coordinate direction and the prefix projection
	computed in $\F_q^m$, by the coordinate order fixed in
	Definition~\ref{def:ld-combining-map}, and cases~2 and~4 of
	Definition~\ref{def:ld-combined-strategy} measure exactly the canonical
	question of the game with parameters $\mathsf{ldparams}$ attached to that
	sample.  Conditioned on $\mathbf{i} \ge m$, every direction produced by the
	question distribution has vanishing point part, so cases~1, 3 and~5 apply and
	all measure the point question at $u$, which is uniform.

	The two conditionings are legitimate because the coordinate index and
	the two point parts of the common sample are jointly uniform on
	$\{0,\ldots,m+k-1\} \times \F_q^m \times \F_q^m$, so the index is
	independent of the point parts.
\end{proof}

\begin{lemma}[Value of the combined strategy]
	\label{lem:ld-combined-value}
	\lean{MIPStarRE.QPBT.directLdRejectionProbability_directCombinedStrategy_le}
	\lean{MIPStarRE.QPBT.directCombinedStrategy_value_ge}
	\leanok
	\uses{def:ld-combined-strategy, lem:ld-combined-question-law,
		def:ld-win-predicate, def:tensor-product-value}
	The statement holds with the universal constant $C = 10$: if $\mathcal S$
	succeeds with probability at least $1 - \eps$ in the directly indexed
	low-degree game with parameters $\mathsf{ldparams}$, then
	$\mathcal S^{\mathrm c}$ succeeds with probability at least $1 - 10\eps$ in
	the directly indexed low-degree game with parameters
	$\mathsf{ldparams}^{\mathrm c}$.  The value $10$ is used in the absorption of
	Proposition~\ref{prop:ld-simultaneous-general-k}.
\end{lemma}
\begin{proof}
	\leanok
	\uses{def:ld-combined-strategy, lem:ld-combined-question-law,
		def:ld-win-predicate}
	Both games weight their nine ordered type pairs uniformly, so it suffices to
	bound the rejection probability of each branch of the combined game by a
	bounded number of branch rejection probabilities of the original game and to
	sum.  Fix a branch and split according to the event of
	Lemma~\ref{lem:ld-combined-question-law}.

	On the event $\mathbf{i} < m$ the two combined answers are the images of the
	two original answers under the relabelings of
	Definition~\ref{def:ld-combined-strategy}, and each clause of
	Definition~\ref{def:ld-win-predicate} for the combined game follows from the
	corresponding clause for the original game.  For the point-versus-point
	clause this is immediate.  For a line-versus-point clause, let $t$ present
	the sampled combined point on the sampled combined line; then $t$ presents
	the point part of the sampled point on the image line, the combining part of
	the sampled point is $\alpha_0 + t w$ in the diagonal case and $\alpha_0$ in
	the axis-parallel case, and the combined line answer evaluated at $t$ is
	$\sum_r (\alpha_0 + t w)_r f_r(t)$, respectively $\sum_r (\alpha_0)_r f_r(t)$;
	substituting the original clause $f_r(t) = b_r$ for every $r$ turns this into
	the combined point answer $\sum_r \alpha_r b_r$.  For a line-versus-line
	clause, equality of the two original tuples of line answers gives equality of
	their combinations.

	On the event $\mathbf{i} \ge m$ all measured questions are the point question
	at the same point $u$, and every clause of the combined game follows from the
	agreement $b = b'$ of the two original point answers, which is the
	point-versus-point clause of the original game: in cases~3 and~5 the combined
	line answer is a fixed affine function of the measured tuple, and evaluating
	it at the parameter presenting the sampled combined point returns
	$\sum_r \alpha_r b_r$, the combined point answer computed from the same
	tuple.

	Each branch of the combined game is therefore rejected with probability at
	most the sum of two branch rejection probabilities of the original game, and
	summing over the nine branches bounds the rejection probability of
	$\mathcal S^{\mathrm c}$ by an absolute multiple of that of $\mathcal S$.
\end{proof}

\begin{lemma}[Exact linearity in the combining variables]
	\label{lem:ld-combining-exact-linearity}
	\lean{MIPStarRE.QPBT.combinedRestrict}
	\lean{MIPStarRE.QPBT.eval_combinedRestrict}
	\lean{MIPStarRE.QPBT.combinedRestrict_combinePolyTuple}
	\lean{MIPStarRE.QPBT.combinedRestrict_eq_of_eq_combiningLinearForm}
	\lean{MIPStarRE.QPBT.combinedCoef}
	\lean{MIPStarRE.QPBT.combinePolyTuple_combinedCoef_iff}
	\lean{MIPStarRE.QPBT.exists_eq_combiningLinearForm_combinedRestrict_iff}
	\lean{MIPStarRE.QPBT.totalDegree_combinedCoef_le}
	\lean{MIPStarRE.QPBT.combinedRestrict_combiningLinearForm_avg_le}
	\lean{MIPStarRE.QPBT.directCombinedExactLinearity_avg_le}
	\leanok
	\uses{def:ld-combining-map, lem:ld-combining-split,
		lem:schwartz-zippel-total-degree}
	Let $p \in \polyfunc{m+k}{q}{d}$ be a polynomial which is not the combined
	polynomial of any tuple.  Then, for a uniformly random $u \in \F_q^m$, the
	polynomial $p(u,\cdot) \in \F_q[\alpha_0,\ldots,\alpha_{k-1}]$ obtained by
	substituting $u$ for the point coordinates is a linear form
	$\sum_r c_r \alpha_r$ with probability at most $md/q$.
\end{lemma}
\begin{proof}
	\leanok
	\uses{lem:schwartz-zippel-total-degree, def:ld-combining-map,
		lem:ld-combining-split}
	Write $p = \sum_{\mu} c_\mu \, \alpha^\mu$, the sum over exponent vectors
	$\mu \in \N^k$ of the combining coordinates, with coefficients
	$c_\mu \in \polyfunc{m}{q}{d}$.  By
	Definition~\ref{def:ld-combining-map} and
	Lemma~\ref{lem:ld-combining-split}, $p$ is a combined polynomial precisely
	when $c_\mu = 0$ for every $\mu$ outside the $k$ standard basis vectors, and
	$p(u,\cdot)$ is a linear form precisely when $c_\mu(u) = 0$ for every such
	$\mu$.  If $p$ is not combined, fix one such $\mu$ with $c_\mu \neq 0$; the
	total degree of $c_\mu$ is at most $md$, so
	Lemma~\ref{lem:schwartz-zippel-total-degree} bounds the probability of
	$c_\mu(u) = 0$ by $md/q$.
\end{proof}

\begin{lemma}[Recovery of the polynomial-tuple measurements]
	\label{lem:ld-combining-recovery}
	\lean{MIPStarRE.QPBT.totalDegree_combinedRestrict_le}
	\lean{MIPStarRE.QPBT.combinedRestrict_collision_avg_le}
	\lean{MIPStarRE.QPBT.directCombinedRecoveryLocalBound}
	\lean{MIPStarRE.QPBT.directCombinedRecoveryEvent_avg_le}
	\lean{MIPStarRE.QPBT.directCombinedRecoveryLocalBound_avg_le}
	\lean{MIPStarRE.QPBT.directCombinedRecoveryBadEvent_avg_le}
	\lean{MIPStarRE.QPBT.consistencyDefect_le_of_recoveryBound}
	\lean{MIPStarRE.QPBT.consistencyDefect_le_of_recoveryBound_right}
	\lean{MIPStarRE.QPBT.directCombinedPointEquiv}
	\lean{MIPStarRE.QPBT.directCombinedRecovery_relation_one}
	\lean{MIPStarRE.QPBT.directCombinedRecovery_relation_two}
	\lean{MIPStarRE.QPBT.directCombinedRecovery_relation_three}
	\leanok
	\uses{def:ld-meas, def:consistency, def:bracket, def:ld-combining-map,
		lem:ld-combining-split, lem:ld-combining-collision,
		lem:ld-combining-exact-linearity, def:ld-combined-strategy}
	Let $\mathcal S$ be a projective strategy for the directly indexed game with
	parameters $\mathsf{ldparams}$, and let $G^{\mathrm A}$, $G^{\mathrm B}$ be
	measurements with outcomes in $\polyfunc{m+k}{q}{d}$ satisfying, for the
	combined strategy $\mathcal S^{\mathrm c}$ and with error $\delta$, the three
	conclusions of Theorem~\ref{lem:ld-soundness} at parameters
	$\mathsf{ldparams}^{\mathrm c}$, using the relabeled point measurements of
	Definition~\ref{def:ld-combined-strategy}.  Let $H^w$ be $G^w$ post-processed by the
	map $\mathrm{split}$ of Lemma~\ref{lem:ld-combining-split}, so that
	$H^w \in \mathrm{PolyMeas}(m,q,d,k)$.  Then the three conclusions of
	Theorem~\ref{lem:ld-soundness} hold for $\mathcal S$ at parameters
	$\mathsf{ldparams}$, with the measurements $H^{\mathrm A}$, $H^{\mathrm B}$
	and with error $O(\delta + (m+k)d/q)$, where the point effects in the first
	two conclusions are $\widehat A^{\mathsf{Point},u}$ and
	$\widehat B^{\mathsf{Point},u}$, respectively.
\end{lemma}
\begin{proof}
	\leanok
	\uses{lem:ld-combining-collision, lem:ld-combining-exact-linearity,
		lem:ld-combining-split, def:bracket, lem:schwartz-zippel-total-degree}
	The first conclusion for $\mathcal S^{\mathrm c}$ says that, on average over
	a uniform $(u,\alpha) \in \F_q^{m+k}$, the value $p(u,\alpha)$ of the outcome
	$p$ of $G^{\mathrm B}$ agrees with the combined point answer
	$\sum_r \alpha_r b_r$ computed from the outcome $b$ of the point measurement
	$\widehat A^{\mathsf{Point},u}$ of $\mathcal S$ at $u$, up to $\delta$.

	The outcome is a combined polynomial with probability at least
	$1 - O(\delta + (m+k)d/q)$.  Indeed, condition on an outcome $p$ that is not
	combined.  By Lemma~\ref{lem:ld-combining-exact-linearity}, $p(u,\cdot)$ fails
	to be a linear form except with probability $md/q$ over $u$; the combined
	point answer is a linear form in $\alpha$ by construction; and two distinct
	polynomials in $\alpha$ of total degree at most $kd$ agree at a uniform
	$\alpha$ with probability at most $kd/q$ by
	Lemma~\ref{lem:schwartz-zippel-total-degree}.  Hence, conditioned on a
	non-combined outcome, the agreement above holds with probability at most
	$(m+k)d/q$, and comparing with the bound $\delta$ gives the claim.

	Replacing $G^w$ by $H^w$ therefore changes each of the first two relations by
	at most $O(\delta + (m+k)d/q)$, since the two measurements differ only on the
	event that the outcome is not combined, and
	$\mathrm{split}(\mathrm{comb}_g) = g$ on its complement; the third relation
	for $H^{\mathrm A}$, $H^{\mathrm B}$ follows from the third relation for
	$G^{\mathrm A}$, $G^{\mathrm B}$ by data processing, $H^w$ being a
	post-processing of $G^w$.

	It remains to upgrade the first relation from the combined value to the tuple
	of values.  Let $g$ be the outcome of $H^{\mathrm B}$ and $b$ the outcome of
	the relabeled point measurement $\widehat A^{\mathsf{Point},u}$.
	If $(g_0(u),\ldots,g_{k-1}(u)) \neq b$, then by
	Lemma~\ref{lem:ld-combining-collision} the combined values
	$\sum_r \alpha_r g_r(u)$ and $\sum_r \alpha_r b_r$ agree with probability at
	most $1/q$ over a uniform $\alpha$.  Averaging over $\alpha$ therefore
	converts the relation for the combined values into the tuple relation
	\[
		\widehat A^{\mathsf{Point},\,u}_{a_0,\ldots,a_{k-1}} \ot \Id_{\mathrm B}
		\simeq \Id_{\mathrm A} \ot
			H^{\mathrm B}_{[(g_0(u),\ldots,g_{k-1}(u)) =
				(a_0,\ldots,a_{k-1})]}
	\]
	with error $O(\delta + (m+k)d/q)$, and symmetrically for the second relation.
\end{proof}

\begin{proposition}[Simultaneous polynomial measurements for general $k$]
	\label{prop:ld-simultaneous-general-k}
	\lean{MIPStarRE.QPBT.exists_directSimultaneousPolynomialMeasurements_combinedError}
	\lean{MIPStarRE.QPBT.exists_directCombinedTransportConstants}
	\lean{MIPStarRE.QPBT.exists_direct_ld_soundness}
	\leanok
	\uses{def:ld-meas, def:consistency, def:bracket,
		def:projective-strategy-general, def:ld-combining-parameters,
		lem:ld-combined-value, lem:ld-combining-recovery, lem:ld-soundness,
		thm:main-formal}
	There exist universal constants $a \ge 1$ and $0 < b \le 1$ such that the
	following holds.  For every directly indexed parameter tuple
	$\mathsf{ldparams} = (q,m,d,k)$, every $\eps > 0$, and every projective
	strategy $\mathcal S = (\psi, A, B)$ succeeding with probability at least
	$1 - \eps$ in the directly indexed low-degree game with those parameters,
	there exist measurements $G^{\mathrm A}, G^{\mathrm B} \in
	\mathrm{PolyMeas}(m,q,d,k)$ satisfying the three consistency relations of
	Theorem~\ref{lem:ld-soundness} with error
	$\delta_{\mathrm{ld}}(\eps,q,m,d,k)$ evaluated at these constants, using
	$\widehat A^{\mathsf{Point},u}$ and $\widehat B^{\mathsf{Point},u}$ from
	Definition~\ref{def:ld-combined-strategy} in the first two relations.

	Before scalar absorption, the first two relations hold with the error of
	Theorem~\ref{thm:main-formal} at the combined parameters
	$\mathsf{ldparams}^{\mathrm c}$ of
	Definition~\ref{def:ld-combining-parameters}, sampling parameter
	$2560000\,(m+k)^{3}d$ and pass bound $30\eps$ --- increased by the recovery
	loss $(m+k)d/q$; the polynomial/polynomial relation carries the same error
	without the recovery loss.  For $0 < \eps \le 1$ the former error is at most
	$\delta_{\mathrm{ld}}(\eps,q,m,d,k)$, and for $\eps \ge 1$ this error
	function is at least one while a consistency defect on a unit state is at
	most one.

	The source records this error in the abbreviated form
	$\delta_{\mathrm{ld}}(\eps, q, m+k, d, 1) + O((m+k)d/q)$; the first summand
	is the error function of Theorem~\ref{lem:ld-soundness} at the combined
	parameters, which is not identified here with the error of
	Theorem~\ref{thm:main-formal} at those parameters.
\end{proposition}
\begin{proof}
	\leanok
	\uses{lem:ld-combined-value, lem:ld-combining-recovery,
		def:ld-combining-parameters, lem:direct-ld-simultaneous-single,
		thm:ld-combined-strategy-projective, thm:main-formal, lem:ld-soundness}
	By Theorem~\ref{thm:ld-combined-strategy-projective},
	$\mathcal S^{\mathrm c}$ is projective, and
	Lemma~\ref{lem:ld-combined-value} shows that it succeeds with probability at
	least $1 - 10\eps$ in the directly indexed game with parameters
	$\mathsf{ldparams}^{\mathrm c}$.  That game has simultaneity parameter $1$, so
	Lemma~\ref{lem:direct-ld-simultaneous-single} --- the polynomial-tuple
	conclusion for simultaneity parameter $1$, obtained for the directly indexed
	game from Theorem~\ref{thm:main-formal} --- applies and produces measurements
	$G^{\mathrm A}$, $G^{\mathrm B}$ with outcomes in $\polyfunc{m+k}{q}{d}$
	satisfying the three relations with the error of
	Theorem~\ref{thm:main-formal} at the combined parameters, sampling parameter
	$2560000\,(m+k)^{3}d$ and pass bound $3 \cdot 10\eps = 30\eps$.
	Lemma~\ref{lem:ld-combining-recovery} converts them into measurements in
	$\mathrm{PolyMeas}(m,q,d,k)$ satisfying the three relations for
	$\mathcal S$ with the stated error.

	The absorption is carried out with the universal constants $a = 10^{23}$ and
	$b = 1/80000$, the exponent being half the exponent $1/40000$ carried by
	Theorem~\ref{thm:main-formal}.  For $0 < \eps \le 1$ the error of
	Theorem~\ref{thm:main-formal} at the combined parameters and the sampling
	parameter $2560000\,(m+k)^{3}d$ carries the prefactor
	$655360000000000000\,(m+k)^{10}d^{2}$, and its three terms are bounded by
	$30\,\eps^{b}$, $d\,q^{-b}$ and $2^{-bmd}$, while the recovery loss
	$(m+k)d/q$ is bounded by $(m+k)d\,q^{-b}$.  Since $m + k \le 2dmk$ and
	$d \le dmk$, the total is at most
	$2.048 \cdot 10^{22}\,(dmk)^{13}\bigl(\eps^{b} + q^{-b} + 2^{-bmd}\bigr)$,
	hence at most $\delta_{\mathrm{ld}}(\eps,q,m,d,k)$ at those constants.  For
	$\eps \ge 1$ the error function is at least one, while the consistency defect
	of a pair of measurement families on a unit bipartite state is at most one, so
	the three relations hold for any witnesses.
\end{proof}

% This node is not a statement from the paper: it is the formalization-support
% specialization of Theorem lem:ld-soundness to the simultaneity parameter
% k = 1, which the proof of Lemma lem:qld-4-7 is the only consumer of.
\begin{theorem}[Formalization support for Theorem~\ref{lem:ld-soundness}: one simultaneous coordinate]
	\label{lem:ld-soundness-k-one}
	\lean{MIPStarRE.QPBT.exists_ld_soundness_of_k_eq_one}
	\leanok
	\uses{def:ld-game, def:ld-question-distribution, def:ld-win-predicate, def:ld-meas, def:projective-strategy-general, def:tensor-product-value, def:consistency, def:bracket}
	This formalization-only auxiliary is not a named statement of the paper; it is the case $k = 1$ of Theorem~\ref{lem:ld-soundness}, the only value instantiated in the proof of Lemma~\ref{lem:qld-4-7} and, by Remark~\ref{rem:ld-soundness-simultaneity}, the only value the coordinatewise route can supply. There exists a function $\delta_{\mathrm{ld}}$ of the shape displayed in Theorem~\ref{lem:ld-soundness}, for universal constants $a \ge 1$ and $0 < b \le 1$, such that all three consistency conclusions of that theorem hold verbatim for every parameter tuple $\mathsf{ldparams} = (q,m,d,k)$ with $k = 1$, every $\eps > 0$, and every projective strategy succeeding with probability at least $1-\eps$ in $\mathfrak{G}^{\mathrm{ld}}_{\mathsf{ldparams}}$.
\end{theorem}
\begin{proof}
	\leanok
	\uses{lem:direct-ld-simultaneous-single, lem:direct-ld-residue-compression, fact:triangle-for-simeq, lem:sandwich-codeword-defect, lem:direct-ld-polynomial-agreement, lem:ld-point-agreement, lem:ld-transport-absorption, lem:ld-scalar-square-root}
	The proof does not follow the source's route through the tensor-code test; it applies Theorem~\ref{thm:main-formal} directly, and therefore neither of the two obligations enumerated in the proof of Theorem~\ref{lem:ld-soundness} is used or assumed.

	The game $\mathfrak{G}^{\mathrm{ld}}_{\mathsf{ldparams}}$ carries a seed $s$ whose fibre $\chi^{-1}(i)$ is not recorded by the $(m,q,d)$ low individual degree test of Definition~\ref{def:low-individual-degree-test}. Dilate the given strategy by a maximally correlated register indexed by that fibre and let each player answer the reconstructed seed-indexed question in the block selected by their copy; the dilation is projective and has exactly the original value, because the residue is uniform and shared. The dilated strategy plays the directly indexed variant of the game, in which the coordinate index is sampled directly instead of through a seed, and its coordinate reading passes the low individual degree test with failure probability at most $3\eps$, the factor three being the ratio of the branch weights of the two tests. Theorem~\ref{thm:main-formal}, instantiated at the auxiliary sampling parameter $2560000\,m^3 d$ --- which satisfies that theorem's large-sampling hypothesis for every admissible tuple --- produces polynomial measurements on the extended Hilbert spaces satisfying the three consistency relations with the error of that theorem. For $k = 1$ these are already the polynomial-tuple measurements required here.

	Compressing the correlated register by its normalized diagonal-block average returns measurements on the original Hilbert spaces. The compression is exact for the two point-versus-polynomial relations, because in each of them the opposite player's operator is the block-diagonal amplification of the original point measurement and therefore leaves the correlated register untouched. It is not exact for the third relation, in which both operators act on that register: on a maximally correlated pair the compressed measurements need not remain close while the uncompressed ones agree exactly. The third relation is therefore reproved on the original Hilbert spaces. The consistency test of Definition~\ref{def:ld-win-predicate} accepts the point-versus-point branch only on equal answers, and that branch carries one ninth of the question law, so the two point measurements are consistent up to $9\eps$.
	The triangle inequality of Lemma~\ref{fact:triangle-for-simeq} combines the three point relations into consistency of the evaluated polynomial measurements. Lemma~\ref{lem:sandwich-codeword-defect}, with the collision bound of Lemma~\ref{lem:direct-ld-polynomial-agreement}, then gives consistency of the polynomial measurements with an additional error $md/q$.

	The resulting error is $C\bigl(\mathrm{poly}(m,d) \cdot \mathrm{poly}(\eps, q^{-1}, e^{-md}) + md/q + \eps\bigr)^{1/2}$ with $C=10$, which is at most $\delta_{\mathrm{ld}}(\eps,q,m,d,k)$ for the universal constants $a = 2.5 \cdot 10^{10}$ and $b = 1/80000$: the square root halves the exponent $1/40000$ carried by the error of Theorem~\ref{thm:main-formal}, and the polynomial prefactor is absorbed by $(dmk)^a$. Outside the regime where both $\eps$ and that argument are at most one the conclusion is vacuous, since $\delta_{\mathrm{ld}} \ge 1$ there while a consistency defect of a pair of measurement families on a unit state never exceeds $1$.
\end{proof}

\begin{theorem}[Formalization support: a scalar square-root bound]
\label{lem:ld-scalar-square-root}
\lean{MIPStarRE.QPBT.le_sqrt_of_le_of_le_one}
\leanok
For real numbers $x,T$ with $x\le T\le1$, one has $x\le\sqrt T$,
where the real square root is zero on negative arguments.
This is a scalar auxiliary, not a statement of the source paper.
\end{theorem}
\begin{proof}
\leanok
If $x\le0$, use nonnegativity of the square root. Otherwise
$0<x\le T\le1$ implies $x^2\le x\le T$.
\end{proof}

\begin{theorem}[Formalization support: absorption of transport bounds]
\label{lem:ld-transport-absorption}
\lean{MIPStarRE.QPBT.consistencyDefect_le_deltaLd_of_transportBound}
\leanok
\uses{def:consistency, lem:direct-ld-coordinate-soundness}
For directly indexed parameters $D=(q,m,d,k)$, write
\[
 T_D(\eps)=E_D(\eps)+md/q+\eps,
\]
where $E_D(\eps)$ is the error of Theorem~\ref{thm:main-formal} at
sampling parameter $2560000m^3d$ and pass bound $3\eps$.
Fix $a\ge1$, $b>0$, and $C_0\ge1$, and assume that for every such $D$
and $0<\eps\le1$,
\[
 C_0 k\sqrt{T_D(\eps)}\le
 \delta_{\mathrm{ld}}(\eps,q,m,d,k).
\]
Here $\delta_{\mathrm{ld}}$ has the displayed formula of
Theorem~\ref{lem:ld-soundness} with constants $a,b$.
For any fixed $D$ and $\eps>0$, let $\mu$ be a probability distribution
on a finite question set, let $A,B$ be families of POVMs with a common
finite outcome set on finite-dimensional Hilbert spaces, and let $\psi$
be a unit bipartite state. Denote their consistency defect by $\Delta$.
If $\Delta\le C_0\sqrt{T_D(\eps)}$ whenever both $\eps\le1$ and
$T_D(\eps)\le1$, then
$\Delta\le\delta_{\mathrm{ld}}(\eps,q,m,d,k)$.
This is an auxiliary estimate with its scalar absorption hypothesis
displayed explicitly, not an additional hypothesis of the paper theorem.
\end{theorem}
\begin{proof}
\leanok
A POVM consistency defect on a unit state is at most one. If $\eps>1$,
the error formula is at least one. If $\eps\le1<T_D(\eps)$, use
$1\le C_0\sqrt{T_D(\eps)}$; otherwise use the assumed defect bound.
In the latter two cases, $k\ge1$ and the scalar absorption hypothesis
give the conclusion.
\end{proof}

\begin{theorem}[Formalization support: directly indexed soundness for $k=1$]
\label{lem:direct-ld-soundness-k-one}
\lean{MIPStarRE.QPBT.exists_direct_ld_soundness_of_k_eq_one}
\leanok
\uses{lem:direct-ld-simultaneous-single, def:ld-meas, def:consistency}
There exist universal constants $a\ge1$ and $0<b\le1$ such that for
every directly indexed parameter tuple $D=(q,m,d,k)$ with $k=1$, every
$\eps>0$, and every projective strategy of value at least $1-\eps$
for the directly indexed game, there are polynomial-tuple POVMs
$G^{\mathrm A},G^{\mathrm B}$ satisfying all three conclusions of
Theorem~\ref{lem:ld-soundness} with error
$\delta_{\mathrm{ld}}(\eps,q,m,d,k)$: the point/polynomial and
polynomial/point defects are averaged over uniform $u\in\F_q^m$,
and the polynomial/polynomial defect has no question parameter.
The directly indexed domain requires $m,d,k\ge1$ and admissible $q$,
but not $m\mid q$. This is a formalization-only specialization for that
game, not the unrestricted source theorem.
\end{theorem}
\begin{proof}
\leanok
\uses{lem:direct-ld-simultaneous-single, lem:ld-scalar-square-root,
lem:ld-transport-absorption}
Apply Lemma~\ref{lem:direct-ld-simultaneous-single}. Its error $E_D(\eps)$
is at most $T_D(\eps)$, hence at most $\sqrt{T_D(\eps)}$ when
$T_D(\eps)\le1$. The scalar absorption estimate holds with
$C_0=1$, $a=2.5\cdot10^9$, and $b=1/80000$.
Theorem~\ref{lem:ld-transport-absorption} applies to each defect.
\end{proof}

\begin{theorem}[Formalization support: point agreement and residue compression]
\label{lem:ld-point-agreement}
\lean{MIPStarRE.QPBT.directLdPointPair_consistencyDefect_le}
\lean{MIPStarRE.QPBT.averageDiagonalBlock_blockDiagonal_const}
\lean{MIPStarRE.QPBT.ldStrategyToDirect_pointPair_compression}
\lean{MIPStarRE.QPBT.ldPointPair_consistencyDefect_le}
\leanok
\uses{def:ld-game, def:consistency, lem:direct-ld-residue-compression}
These are formalization auxiliaries, not named assertions of the paper.
For directly indexed parameters $D=(q,m,d,k)$, a projective strategy of
value at least $1-\eps$, with $\eps\in\R$, has point-readout consistency
defect at most $9\eps$, averaged over uniform $u\in\F_q^m$.
The readouts send malformed point answers to the zero tuple.
The same assertion holds for seed-indexed parameters $L=(q,m,d,k)$
(which additionally require $m\mid q$).

More precisely, for any seed-indexed strategy, without projectivity or
value assumptions, its point-readout defect equals that of its
correlated-residue extension from Lemma~\ref{lem:direct-ld-residue-compression}.
The matrix identity used here is valid for every operator $M$ on a finite
Hilbert space and every nonempty finite register $R$:
\[
 \frac1{|R|}\sum_{r\in R}
 (M\ot I_R)_{r,r}=M,
\]
where the subscripts denote diagonal blocks indexed by the register.
\end{theorem}
\begin{proof}
\leanok
\uses{lem:direct-ld-residue-compression}
In the directly indexed game, unequal point readouts imply rejection of
the point-versus-point branch. That branch has weight $1/9$, giving the
bound $9\eps$. Every diagonal block of $M\ot I_R$ equals $M$, so its
normalized average is $M$. Point measurements in the extended strategy
are these constant block-diagonal amplifications; compression on one side
and this identity on the other give equality of the point defects.
The extension preserves projectivity and value, so the seed-indexed bound
follows from the directly indexed bound.
\end{proof}

\section{The Magic Square game}
% Paper label sec:ms kept for cross-reference fidelity.
\label{sec:ms}

We recall the Magic Square game of Mermin and Peres. It is a simple self-test for two EPR pairs, and it allows one to test that a pair of binary observables \emph{anticommutes}; both properties are used in the Pauli basis test. Here it is presented as a binary constraint satisfaction game. Throughout this section and the next, $\sigma^X$ and $\sigma^Z$ denote the single-qubit Pauli observables, with associated two-outcome projective measurements $\{\sigma^W_b\}_{b \in \F_2}$ for $W \in \{X,Z\}$.

\begin{definition}[The Magic Square game]\label{def:ms-game}
	\uses{def:game, def:graph-distribution}
	The Magic Square game $\mathfrak{G}^{\mathrm{MS}}$ is the following two-player game. There are $9$ variables $x_1,\ldots,x_9$ taking values in $\F_2$, arranged in a $3 \times 3$ grid: the cell in row $r$ and column $c$ contains the variable $x_{3(r-1)+c}$. There are $6$ linear equations over these variables: for each of the three rows and each of the first two columns, the sum of the three variables in it equals $0$; for the third column, the sum of the three variables in it equals $1$.

	The question types are
	\begin{align*}
		\mathcal T^{\mathrm{mc}} &= \{\mathsf{Constraint}_i : i = 1,2,\ldots,6\}\;,\\
		\mathcal T^{\mathrm{mv}} &= \{\mathsf{Variable}_j : j = 1,2,\ldots,9\}\;,\\
		\mathcal T^{\mathrm{ms}} &= \mathcal T^{\mathrm{mc}} \cup \mathcal T^{\mathrm{mv}}\;,
	\end{align*}
	where $\mathsf{Constraint}_i$ for $i \in \{1,2,3\}$ denotes the constraint of row $i$ (with variables $x_{3(i-1)+1}, x_{3(i-1)+2}, x_{3(i-1)+3}$), $\mathsf{Constraint}_4$ and $\mathsf{Constraint}_5$ denote the constraints of the first two columns, and $\mathsf{Constraint}_6$ denotes the constraint of the third column (with variables $x_{i-3}, x_{i}, x_{i+3}$ for $\mathsf{Constraint}_i$, $i \in \{4,5,6\}$).

	The question distribution is as follows: the verifier samples $\mathsf{Constraint}_i$ uniformly from $\mathcal T^{\mathrm{mc}}$, then samples $\mathsf{Variable}_j$ uniformly among the three variables appearing in the row or column of $\mathsf{Constraint}_i$; one player, chosen uniformly at random, receives $\mathsf{Constraint}_i$ and the other receives $\mathsf{Variable}_j$. Equivalently, the verifier samples a uniformly random edge of the bipartite \emph{type graph} $G^{\mathrm{ms}}$ on vertex sets $\mathcal T^{\mathrm{mc}}$ and $\mathcal T^{\mathrm{mv}}$, whose edges are $\mathsf{Constraint}_i \,\text{---}\, \mathsf{Variable}_{3(i-1)+j}$ for $i,j \in \{1,2,3\}$ and $\mathsf{Constraint}_i \,\text{---}\, \mathsf{Variable}_{(i-3)+3(j-1)}$ for $i \in \{4,5,6\}$, $j \in \{1,2,3\}$ ($18$ edges in total, and no self-loops), and sends one endpoint to a uniformly random player and the other endpoint to the other player.

	The $\mathsf{Constraint}$ player answers with three bits $(\beta_{v_1}, \beta_{v_2}, \beta_{v_3}) \in \F_2^3$, where $(v_1,v_2,v_3)$ are the indices of the three variables of $\mathsf{Constraint}_i$; the $\mathsf{Variable}$ player answers with a single bit $\gamma \in \F_2$. The players win if and only if the $\mathsf{Constraint}$ player's answer satisfies the equation associated with $\mathsf{Constraint}_i$ and $\gamma = \beta_j$.
\end{definition}

The following theorem records the self-testing (rigidity) properties of the Magic Square game. It is used in the soundness analysis of the Pauli basis test to enforce anticommutation relations between certain pairs of observables.
With zero-based labels, the paper's variables $1$ and $5$ have labels $0$ and
$4$, respectively.

\begin{definition}[Variable-measurement agreement distance]\label{def:ms-variable-agreement}
	\lean{MIPStarRE.QPBT.msVariableConsistencyDefect}
	\leanok
	\uses{def:ms-game, def:povm-distance, def:consistency}
	% Formalization-support node (issue #172, owner decision B5 on issue #26): it names the
	% quantity bounded by the corrected hypothesis of Theorem thm:ms-rigidity, and fixes the
	% binary relabeling of the answer alphabet that the Lean definition applies.
	Let $\mathcal A^{\mathrm{ms}}$ be the answer alphabet of $\mathfrak{G}^{\mathrm{MS}}$ (Definition~\ref{def:ms-game}), whose elements are the parity triples $\F_2^3$ and the single cell bits $\F_2$, and let $\beta \colon \mathcal A^{\mathrm{ms}} \to \F_2$ be the relabeling that sends a single bit $\gamma$ to $\gamma$ and every parity triple --- an answer of the wrong form for a $\mathsf{Variable}$ question --- to $0$. For a strategy $\mathcal S = (\ket{\psi}, A, B)$ for $\mathfrak{G}^{\mathrm{MS}}$ and a variable $\mathsf{Variable}_j$, the associated \emph{binary effects} are
	\[
		\hat A^{\mathsf{Variable}_j}_b = \sum_{a \,:\, \beta(a) = b} A^{\mathsf{Variable}_j}_a \;, \qquad \hat B^{\mathsf{Variable}_j}_b = \sum_{a \,:\, \beta(a) = b} B^{\mathsf{Variable}_j}_a \qquad (b \in \F_2)\;,
	\]
	two-outcome POVMs that agree with the reported bit on answers of the form prescribed for a $\mathsf{Variable}$ question. The \emph{variable-measurement agreement distance} of $\mathcal S$ at $\mathsf{Variable}_j$ is
	\[
		\Delta_j(\mathcal S) = \sum_{b \in \F_2} \bigl\| \bigl( \hat A^{\mathsf{Variable}_j}_b \ot I_{\mathrm B} - I_{\mathrm A} \ot \hat B^{\mathsf{Variable}_j}_b \bigr) \ket{\psi} \bigr\|^2 \;,
	\]
	the squared state-dependent distance of Definition~\ref{def:povm-distance} between the two players' binary $\mathsf{Variable}_j$ effects on $\ket{\psi}$, under the one-point distribution and with the constant inside the $O(\cdot)$ taken to be $1$. It is not the consistency defect of Definition~\ref{def:consistency}, which is an off-diagonal outcome weight.
\end{definition}

\begin{remark}[The name of the agreement distance in the formalization]\label{rem:ms-variable-agreement-name}
	The Lean declaration for Definition~\ref{def:ms-variable-agreement} is \texttt{msVariableConsistencyDefect}. The word \emph{consistency} there names the failure of the two players' $\mathsf{Variable}_j$ measurements to be consistent with one another, which is what $\Delta_j(\mathcal S)$ measures; it does not name the consistency defect of Definition~\ref{def:consistency}, the off-diagonal outcome weight of a pair of measurements. The two quantities are different, and only the second is the source's consistency defect.
\end{remark}

\begin{theorem}[Rigidity of the Magic Square game]\label{thm:ms-rigidity}
	\lean{MIPStarRE.QPBT.exists_ms_rigidity}
	\leanok
	\uses{def:ms-game, def:tensor-product-value, def:EPR, def:povm-distance, def:consistent-measurement, def:ms-variable-agreement}
	% Corrected statement (issue #172, owner decision B5 on issue #26). The source, at
	% references/qpbt-paper/08_classical_and_quantum_low_degree_tests.tex:612-652, quantifies over
	% every strategy of value 1 - eps and carries no consistency hypothesis; that statement is
	% false (see the proof, Remark~\ref{rem:ms-rigidity-strategy-class}, and the cited note). The
	% case delta = 0 below recovers the source's error scale for completed binary effects;
	% comparison with the prescribed-answer effects remains a separate construction.
	Let $\eps, \delta \ge 0$, and let $\mathcal S = (\ket{\psi}, A, B)$ be a strategy that succeeds in the Magic Square game $\mathfrak{G}^{\mathrm{MS}}$ with probability at least $1-\eps$, where $\ket{\psi} \in \hilb_{\mathrm A} \ot \hilb_{\mathrm B}$, and whose $\mathsf{Variable}_1$ and $\mathsf{Variable}_5$ measurements are consistent between the two players up to $\delta$ on $\ket{\psi}$, in the binary effects $\hat A$, $\hat B$ of Definition~\ref{def:ms-variable-agreement}:
	\begin{equation}\label{eq:ms-rigidity-consistency}
		\Delta_j(\mathcal S) = \sum_{b \in \F_2} \bigl\| \bigl( \hat A^{\mathsf{Variable}_j}_b \ot I_{\mathrm B} - I_{\mathrm A} \ot \hat B^{\mathsf{Variable}_j}_b \bigr) \ket{\psi} \bigr\|^2 \le \delta \qquad \text{for } j \in \{1, 5\}\,.
	\end{equation}
	This is the relation $\hat A^{\mathsf{Variable}_j}_b \ot I_{\mathrm B} \approx_\delta I_{\mathrm A} \ot \hat B^{\mathsf{Variable}_j}_b$ of Definition~\ref{def:povm-distance} with respect to $\ket{\psi}$, with constant $1$; for a symmetric strategy and $\delta = 0$ it says that the two binary measurements are consistent on $\ket{\psi}$ in the sense of Definition~\ref{def:consistent-measurement}, so every symmetric consistent strategy satisfies it with $\delta = 0$. Write $\eta = \sqrt\eps + \sqrt\delta$. Then there exist local isometries $\phi_{\mathrm A}\colon \hilb_{\mathrm A} \to \hilb_{\mathrm A'} \ot \hilb_{\mathrm A''}$ and $\phi_{\mathrm B}\colon \hilb_{\mathrm B} \to \hilb_{\mathrm B'} \ot \hilb_{\mathrm B''}$, where $\hilb_{\mathrm A'}, \hilb_{\mathrm B'} \cong (\C^2)^{\ot 2}$ and $\hilb_{\mathrm A''}, \hilb_{\mathrm B''}$ are finite dimensional, and a state $\ket{\mathrm{aux}} \in \hilb_{\mathrm A''} \ot \hilb_{\mathrm B''}$, such that:
	\begin{enumerate}
		\item $\bigl\| \phi_{\mathrm A} \ot \phi_{\mathrm B} \ket{\psi} - \ket{\mathrm{EPR}_2}^{\ot 2} \ot \ket{\mathrm{aux}} \bigr\| \le O(\eta)$;
		\item letting $\tilde A^{\mathsf{Variable}_j}_b = \phi_{\mathrm A} \hat A^{\mathsf{Variable}_j}_b \phi_{\mathrm A}^\dagger$ and $\tilde B^{\mathsf{Variable}_j}_b = \phi_{\mathrm B} \hat B^{\mathsf{Variable}_j}_b \phi_{\mathrm B}^\dagger$ denote the conjugated binary effects of Definition~\ref{def:ms-variable-agreement},
		\begin{align*}
			\tilde A^{\mathsf{Variable}_1}_b \ot I_{\mathrm{B'B''}} &\approx_{\eta} (\sigma^X_b)_{\mathrm A'} \ot I_{\mathrm{A''B'B''}}\;,\\
			\tilde A^{\mathsf{Variable}_5}_b \ot I_{\mathrm{B'B''}} &\approx_{\eta} (\sigma^Z_b)_{\mathrm A'} \ot I_{\mathrm{A''B'B''}}\;,\\
			I_{\mathrm{A'A''}} \ot \tilde B^{\mathsf{Variable}_1}_b &\approx_{\eta} I_{\mathrm{A'A''B''}} \ot (\sigma^X_b)_{\mathrm B'}\;,\\
			I_{\mathrm{A'A''}} \ot \tilde B^{\mathsf{Variable}_5}_b &\approx_{\eta} I_{\mathrm{A'A''B''}} \ot (\sigma^Z_b)_{\mathrm B'}\;,
		\end{align*}
		where each $\approx_{\eta}$ holds with respect to the extracted ideal state $\ket{\Theta} = \ket{\mathrm{EPR}_2}^{\ot 2}_{\mathrm{A'B'}} \ot \ket{\mathrm{aux}}_{\mathrm{A''B''}}$ and the answer summation is over $b \in \{0,1\}$, and the approximation bounds carry implicit universal constants.
	\end{enumerate}
	As a consequence, letting $\tilde A^{\mathsf{Variable}_1} = \tilde A^{\mathsf{Variable}_1}_0 - \tilde A^{\mathsf{Variable}_1}_1$ and $\tilde A^{\mathsf{Variable}_5} = \tilde A^{\mathsf{Variable}_5}_0 - \tilde A^{\mathsf{Variable}_5}_1$ denote the associated observables, it holds that
	\begin{align*}
		\tilde A^{\mathsf{Variable}_1} \tilde A^{\mathsf{Variable}_5} \ot I_{\mathrm{B'B''}} &\approx_{\eta} -\, \tilde A^{\mathsf{Variable}_5} \tilde A^{\mathsf{Variable}_1} \ot I_{\mathrm{B'B''}}\;,\\
		I_{\mathrm{A'A''}} \ot \tilde B^{\mathsf{Variable}_1} \tilde B^{\mathsf{Variable}_5} &\approx_{\eta} -\, I_{\mathrm{A'A''}} \ot \tilde B^{\mathsf{Variable}_5} \tilde B^{\mathsf{Variable}_1}\;,
	\end{align*}
	where, for operators $M$, $N$ that are not POVM elements, the notation $M \approx_\zeta N$ means $\bra{\Theta} (M-N)^\dagger (M-N) \ket{\Theta} \le O(\zeta)$, evaluated on the extracted ideal state $\ket{\Theta}$ of item~2 --- the state on which the transported operators act --- and not on $\ket{\psi}$. For $\delta = 0$, in particular for every symmetric consistent strategy, $\eta = \sqrt\eps$.
	This recovers the source's error scale for the completed binary effects of
	Definition~\ref{def:ms-variable-agreement}; comparison with the
	prescribed-answer effects remains a separate construction. Zero agreement
	does not eliminate answers of the wrong form.
\end{theorem}
\begin{proof}
	\leanok
	\uses{thm:naimark, lem:povm-to-obs, lem:ms-varobs-cross-player, def:ms-variable-agreement, def:povm-distance, lem:ms-rigidity-support-swap-algebra, lem:ms-rigidity-support-constants, lem:ms-rigidity-support-ideal-target, lem:ms-rigidity-support-two-qubit-swap, lem:ms-rigidity-support-second-pair, lem:ms-rigidity-support-swap-pair, lem:ms-rigidity-support-swap-transport, lem:ms-rigidity-support-joint-state, lem:ms-rigidity-support-intertwine, lem:ms-rigidity-support-marginals, lem:ms-rigidity-support-assembly, lem:ms-rigidity-support-transfer}
	The source imports this result, for every strategy of value $1-\eps$, from Theorem~6.9 of~\cite{Coladangelo2017Robust}, noting two minor differences between the statements. First, that theorem states a robustness of $O(\eps)$, but measures the closeness of the actual and ideal states in trace norm; expressed as the Euclidean distance between $\phi_{\mathrm A} \ot \phi_{\mathrm B}\ket{\psi}$ and $\ket{\mathrm{EPR}_2}^{\ot 2} \ot \ket{\mathrm{aux}}$, this becomes $O(\sqrt\eps)$. Second, its ideal observables for the $\mathsf{Variable}_1$ and $\mathsf{Variable}_5$ questions are $I \ot \sigma^Z$ and $\sigma^Z\sigma^X \ot \sigma^Z\sigma^X$; under a local change of basis these are equivalent to $\sigma^X \ot I$ and $\sigma^Z \ot I$, respectively.

	Stated for every strategy, as in the source, the result is false, and the hypothesis~\eqref{eq:ms-rigidity-consistency} is the correction adopted for this formalization (\cite{gap:qpbt_ms-rigidity-symmetric-strategies}; owner decision B5 on issue~\#26). Theorem~6.9 of~\cite{Coladangelo2017Robust} is a self-test for the \emph{one-way} linear constraint system game, in which one player always receives a constraint and the other always receives a variable; its conclusion constrains the constraint player's observables and the variable player's observables, and says nothing about the constraint player's variable observables, which that game never asks for. In the game of Definition~\ref{def:ms-game} both players receive both kinds of question, and the conclusion above forces the two players' $\mathsf{Variable}_1$ and $\mathsf{Variable}_5$ measurements to agree on $\ket{\psi}$ --- a pair of questions that the question distribution never samples, being supported on constraint--variable pairs. Remark~\ref{rem:ms-rigidity-strategy-class} gives a strategy of value $1$ violating the source's conclusion at $\eps = 0$, and a symmetric one, so neither the value nor symmetry supplies that agreement; both strategies have $\delta = 1$ in~\eqref{eq:ms-rigidity-consistency}, which is precisely the agreement that the conclusion needs.

	Under the hypothesis, the derivation from the cited self-test is the following. Write $p_{\mathrm A}$ and $p_{\mathrm B}$ for the success probabilities of the two orientations, so that the hypothesis reads $\tfrac12(p_{\mathrm A} + p_{\mathrm B}) \ge 1-\eps$. Conditioning on the orientation in which player~$\mathrm A$ receives the constraint gives a strategy for the one-way game of value $p_{\mathrm A} \ge 1-2\eps$, since $p_{\mathrm B} \le 1$; when $\mathcal S$ is symmetric the two orientations have the same success probability and the bound improves to $p_{\mathrm A} \ge 1-\eps$. It is a strategy to which Theorem~6.9 of~\cite{Coladangelo2017Robust} applies after a Naimark dilation (Theorem~\ref{thm:naimark}); it yields $\phi_{\mathrm A}$, $\phi_{\mathrm B}$ and $\ket{\mathrm{aux}}$ with item~1 and with the two relations of item~2 for player~$\mathrm B$ at scale $\sqrt\eps$. For player~$\mathrm A$, the hypothesis gives that $(\tilde A^{\mathsf{Variable}_j}_b \ot I)\,(\phi_{\mathrm A} \ot \phi_{\mathrm B})\ket{\psi} = (\phi_{\mathrm A} \ot \phi_{\mathrm B})\,(A^{\mathsf{Variable}_j}_b \ot I)\ket{\psi}$ is within $\sqrt\delta$ of $(\phi_{\mathrm A} \ot \phi_{\mathrm B})\,(I \ot B^{\mathsf{Variable}_j}_b)\ket{\psi} = (I \ot \tilde B^{\mathsf{Variable}_j}_b)\,(\phi_{\mathrm A} \ot \phi_{\mathrm B})\ket{\psi}$; player~$\mathrm B$'s relation and item~1 move this to $(I \ot \sigma^W_b)\ket{\mathrm{EPR}_2}^{\ot 2} \ot \ket{\mathrm{aux}}$, and $(I \ot \sigma^W_b)\ket{\mathrm{EPR}_2}^{\ot 2} = (\sigma^W_b \ot I)\ket{\mathrm{EPR}_2}^{\ot 2}$ because $\sigma^X_b$ and $\sigma^Z_b$ are real symmetric. Since every operator involved is a contraction, player~$\mathrm A$'s relations hold at scale $\sqrt\eps + \sqrt\delta$, and the anticommutation consequences follow from the four relations as in the source, by the same transposition identity. A second derivation, which avoids importing the cited theorem, stays finite dimensional and extracts the two EPR pairs by swap isometries: each player carries two approximately anticommuting pairs of reflections, by the solution-group computation on the projective dilation of Theorem~\ref{thm:naimark}; the cross-player agreement of the first pair comes from~\eqref{eq:ms-rigidity-consistency} through the leakage bounds of that dilation (Lemma~\ref{lem:ms-varobs-cross-player}) and that of the second from the constraint--variable subtests at the two cells it uses; and item~1 follows from both together. In either derivation this is where the hypothesis enters, and it is also why symmetry is neither assumed nor used.

	The linked declaration states the theorem above, over an arbitrary strategy of the Magic Square game with the two consistency defects bounded by $\delta$, and it is proved. The formal proof splits at $\sqrt\eps + \sqrt\delta \ge 1$, where the seven conclusions hold for the trivial witness; in the complementary regime the witness is the tensor of the two players' two-qubit controlled-swap embeddings, composed with the dilation embedding. Each of those embeddings composes two one-qubit controlled swaps, one for each of two logical pairs of that player's $\pm 1$ reflections on the dilated state, and the two players' pairs are not chosen symmetrically. Bob's are his variable reflections at the cells of $\mathsf{Variable}_1$ and $\mathsf{Variable}_5$ and at those of $\mathsf{Variable}_2$ and $\mathsf{Variable}_4$. Alice's first pair is her variable reflections at the cells of $\mathsf{Variable}_1$ and $\mathsf{Variable}_5$, but her second is her \emph{constraint} reflections at the cells of $\mathsf{Variable}_2$ and $\mathsf{Variable}_4$, read from the first and from the second row. That asymmetry is forced by the hypothesis: \eqref{eq:ms-rigidity-consistency} compares the two players only at $\mathsf{Variable}_1$ and $\mathsf{Variable}_5$, so their $\mathsf{Variable}_2$ and $\mathsf{Variable}_4$ measurements cannot be compared, whereas Alice's constraint reflections at those two cells are compared with Bob's variable reflections there by a question pair that the distribution does sample. Item~1 is the joint state estimate on the dilation, which needs the approximate anticommutation of both of Alice's pairs, the cross-player agreement of all four pairs of corresponding reflections, and the approximate commutation of the two pairs of each player; the anticommutation of the second pair comes from a second closed path through the grid, which crosses the exceptional third column once and therefore returns the reversed product with the opposite sign. The bit-measurement relations come from the transport of the first logical pair through those embeddings, exactly for the phase observable and up to the anticommutator defect of that pair for the shift observable; the anticommutation consequences come from the solution-group computation on the dilation, products composing under conjugation by an isometry. Every conclusion is then moved from the dilated strategy to the original one by the leakage bounds of the dilation, with an explicit universal constant, which the formalization takes to be $2 \cdot 10^{12}$. The specialization to symmetric consistent strategies ($\delta = 0$) is stated as a separate declaration derived from it.
	% Formalization status is carried by the tags: the statement-level \leanok records that
	% the Lean statement matches the statement above, over an arbitrary strategy with both
	% agreement distances bounded by delta, and the proof-level \leanok records that the Lean
	% proof of that statement is complete (issue #105). The specialization to symmetric
	% consistent strategies (delta = 0) is Corollary cor:ms-rigidity-symmetric-consistent, and
	% the cross-player agreement of the dilated reflections is Lemma lem:ms-varobs-cross-player.
\end{proof}

\begin{remark}[The strategy class in Theorem~\ref{thm:ms-rigidity}]\label{rem:ms-rigidity-strategy-class}
	Fix the canonical Mermin--Peres operator solution on $\C^2 \ot \C^2$ --- the nine $\pm 1$ observables $\sigma^X \ot I$, $I \ot \sigma^X$, $\sigma^X \ot \sigma^X$ in the first row, $I \ot \sigma^Z$, $\sigma^Z \ot I$, $\sigma^Z \ot \sigma^Z$ in the second, and $\sigma^X \ot \sigma^Z$, $\sigma^Z \ot \sigma^X$, $\sigma^Y \ot \sigma^Y$ in the third --- and let $\ket{\Omega}$ be the maximally entangled state of two copies of $\C^2 \ot \C^2$. Measuring the three commuting observables of a row or column on a $\mathsf{Constraint}$ question, and the single observable of a cell on a $\mathsf{Variable}$ question, is a symmetric, projective, consistent strategy of value $1$ (Theorem~\ref{thm:ms-from-ac}). Now give each player two copies of $\C^2 \ot \C^2$ paired into two copies of $\ket{\Omega}$, and split the two roles across them: on a constraint question the first player measures copy~$1$ and the second copy~$2$, and on a variable question the first player measures copy~$2$ and the second copy~$1$. Every sampled question pair consists of one constraint and one variable, so the two measurements actually performed are the canonical perfect pair acting on a single copy of $\ket{\Omega}$, and the value is $1$. The two players' $\mathsf{Variable}_1$ observables, however, act on halves of different copies, whose joint reduced state is a product of maximally mixed states, so their correlation is $0$, while the second item of Theorem~\ref{thm:ms-rigidity} forces it to be $1$ at $\eps = \delta = 0$; the same computation applies to $\mathsf{Variable}_5$. Giving each player a flag qubit, putting the two flags in the state $(\ket{01}+\ket{10})/\sqrt2$, and letting each player's flag select which of the two roles it plays turns this into a \emph{symmetric} strategy with the same value and the same vanishing correlations, so a restriction to symmetric strategies (Definition~\ref{def:symmetric-game}) would not repair the source's statement. Agreement of the two players' $\mathsf{Variable}_1$ and $\mathsf{Variable}_5$ measurements on the state, that is, the hypothesis~\eqref{eq:ms-rigidity-consistency} --- exact consistency in the sense of Definition~\ref{def:consistent-measurement} when the strategy is symmetric and $\delta = 0$ --- does force the correlation to be $1$ up to $O(\delta)$, and both strategies above have agreement distance $\Delta_1 = \Delta_5 = 1$ (Definition~\ref{def:ms-variable-agreement}). It is also what the analysis of the Pauli basis test has available when it applies Theorem~\ref{thm:ms-rigidity}: the Magic Square strategy induced at an anticommuting tuple has, on average over the tuple, agreement distances of order $\eps$ by items~1 and~7 of Lemma~\ref{lem:qld-win-implications} (the point measurements are consistent between the players, and each player's point measurement agrees with the other player's variable measurement), once item~1 is conditioned on the anticommutation event and both are passed through the binary relabeling of Definition~\ref{def:ms-variable-agreement}; that analysis uses only the anticommutation consequence for one player. The theorem is therefore stated with that agreement as its hypothesis, for arbitrary and not necessarily symmetric strategies; the class of symmetric consistent strategies named in the owner decision is the case $\delta = 0$. The details are in \cite{gap:qpbt_ms-rigidity-symmetric-strategies}.
\end{remark}

\begin{definition}[The unasserted printed rigidity claim]
  \label{def:ms-rigidity-printed-claim}
  \lean{MIPStarRE.QPBT.PrintedMagicSquareRigidityClaim}
  \leanok
  \uses{def:ms-game, def:tensor-product-value, def:povm-distance}
  Define $\mathrm{PrintedMS}$ to be the following proposition. There is a
  universal constant $C\geq1$ such that, for every $\eps\geq0$ and every
  strategy of the Magic Square game of value exactly $1-\eps$, there are
  local isometries and a unit auxiliary state satisfying all seven bounds
  in the conclusion of Theorem~\ref{thm:ms-rigidity} at scale $C\sqrt\eps$.
  Here the bit effects are the prescribed single-bit answers themselves,
  without assigning wrong-form answers to zero, and the observables are
  the differences of those same effects. Both anticommutator distances
  are read on the ideal extracted state, as are the four measurement
  distances. No agreement or symmetry hypothesis is included.
  This defines a proposition and does not assert it. The independent-copy
  examples in Remark~\ref{rem:ms-rigidity-strategy-class} refute the
  unrestricted assertion. The Lean declaration does not assert its
  refutation either; see \cite{gap:qpbt_ms-rigidity-symmetric-strategies}.
\end{definition}

% Formalization-support and specialization nodes for Theorem~\ref{thm:ms-rigidity}. None of
% the four statements below is a named result of the source article: the first two isolate
% the facts about Definition~\ref{def:ms-variable-agreement} that the formalization proves,
% the third is the estimate that carries the hypothesis to the projective dilation, and the
% last is the specialization to the strategy class of owner decision B5 on issue #26.
\begin{lemma}[Vanishing of the agreement distance on consistent measurements]\label{lem:ms-variable-agreement-zero}
	\lean{MIPStarRE.QPBT.msVariableConsistencyDefect_eq_zero_of_forall, MIPStarRE.QPBT.msVariableConsistencyDefect_eq_zero_of_isConsistentOn, MIPStarRE.QPBT.msVariableConsistencyDefect_eq_zero_of_isConsistent}
	\leanok
	\uses{def:ms-variable-agreement, def:consistent-measurement, def:consistent-strategy, def:symmetric-game}
	Let $\mathcal S = (\ket{\psi}, A, B)$ be a strategy for $\mathfrak{G}^{\mathrm{MS}}$ and let $j$ be a cell. If $(\hat A^{\mathsf{Variable}_j}_b \ot I)\ket{\psi} = (I \ot \hat B^{\mathsf{Variable}_j}_b)\ket{\psi}$ for both $b \in \F_2$, then $\Delta_j(\mathcal S) = 0$. In particular, if $\mathcal S = (\ket{\psi}, M)$ is symmetric and its $\mathsf{Variable}_j$ measurement is consistent on $\ket{\psi}$ in the sense of Definition~\ref{def:consistent-measurement}, then $\Delta_j(\mathcal S) = 0$; and a symmetric strategy all of whose measurements are consistent on $\ket{\psi}$ --- in particular a symmetric consistent strategy in the sense of Definition~\ref{def:consistent-strategy}, projectivity not being used --- has $\Delta_j(\mathcal S) = 0$ at every cell $j$.
\end{lemma}
\begin{proof}
	\leanok
	\uses{def:ms-variable-agreement, def:consistent-measurement, lem:ms-variable-agreement-support}
	Each summand of $\Delta_j(\mathcal S)$ is the squared norm of a vector that the hypothesis makes zero. For the second assertion, consistency of $\{M^{\mathsf{Variable}_j}_a\}$ on $\ket{\psi}$ gives $(M^{\mathsf{Variable}_j}_a \ot I)\ket{\psi} = (I \ot M^{\mathsf{Variable}_j}_a)\ket{\psi}$ for every answer $a$, and the binary effects of Definition~\ref{def:ms-variable-agreement} are finite sums of such effects, so the same equality holds for $\hat M^{\mathsf{Variable}_j}_b$ for both $b$. The third assertion is the second applied at every cell.
\end{proof}

\begin{lemma}[Cross-player agreement of the dilated variable reflections]\label{lem:ms-varobs-cross-player}
	\lean{MIPStarRE.QPBT.MagicSquareRigidity.norm_msVarObsA_sub_msVarObsB_sq_le, MIPStarRE.QPBT.MagicSquareRigidity.msVarObsA_close_msVarObsB_opDistSq, MIPStarRE.QPBT.MagicSquareRigidity.msVarObsA_close_msVarObsB}
	\leanok
	\uses{def:ms-variable-agreement, def:ms-game, def:povm-distance, thm:naimark, def:tensor-product-value}
	Let $\mathcal S = (\ket{\psi}, A, B)$ be a strategy for $\mathfrak{G}^{\mathrm{MS}}$ of value at least $1-\eps$, let $j$ be a cell, and let $\Delta_j(\mathcal S) \le \delta$. Apply the questionwise Naimark dilation of Theorem~\ref{thm:naimark} to both players, and let $\ket{\psi'}$ be the dilated state and $X_{\mathrm A}$, $X_{\mathrm B}$ the two players' $\pm 1$-valued reflections associated with the dilated, projective $\mathsf{Variable}_j$ measurements. Then
	\[
		\bigl\| \bigl( X_{\mathrm A} \ot I - I \ot X_{\mathrm B} \bigr) \ket{\psi'} \bigr\|^2 \le 864\,\eps + 6\,\delta \;,
	\]
	that is, $X_{\mathrm A} \ot I \approx_{864\eps + 6\delta} I \ot X_{\mathrm B}$ on $\ket{\psi'}$ in the convention of Definition~\ref{def:povm-distance}, and the two reflections are within $\sqrt{864\eps + 6\delta}$ of each other in norm on $\ket{\psi'}$.
\end{lemma}
\begin{proof}
	\leanok
	\uses{def:ms-variable-agreement, lem:povm-to-obs, thm:naimark, lem:ms-variable-agreement-support}
	Compress each dilated reflection to the ground slice of the dilation. The compressed difference acts on the dilated state exactly as the difference of the original binary observables $\hat{\mathcal O}_{\mathrm A} = \hat A^{\mathsf{Variable}_j}_0 - \hat A^{\mathsf{Variable}_j}_1$ and $\hat{\mathcal O}_{\mathrm B}$ acts on $\ket{\psi}$, and passing from the two binary effects to their observable costs the answer-alphabet cardinality $2$ by Lemma~\ref{lem:povm-to-obs}, so that squared difference is at most $2\,\Delta_j(\mathcal S) \le 2\delta$. The two complementary terms are the leakages of the dilated reflections out of the ground slice; each has squared norm at most $144\,\eps$, because the cell of index $j$ is incident to a constraint and the corresponding cell-consistency mass of a strategy of value $1-\eps$ is $O(\eps)$. A triangle inequality on the three terms, followed by the inequality $(x+y+z)^2 \le 3(x^2+y^2+z^2)$, gives the displayed bound. The unsquared form follows by taking square roots.
\end{proof}

\begin{lemma}[Formalization support for the agreement distance and its transfer]\label{lem:ms-variable-agreement-support}
	\lean{MIPStarRE.QPBT.msVariableConsistencyDefect_nonneg, MIPStarRE.QPBT.msVariableConsistencyDefect_eq_sum, MIPStarRE.QPBT.applyOperatorToState_finset_sum, MIPStarRE.QPBT.applyOperatorToState_eq_of_mulVec_eq, MIPStarRE.QPBT.applyOperatorToState_postprocess_effect_eq_of_isConsistentOn, MIPStarRE.QPBT.MagicSquareRigidity.naimarkDilatedState_sub, MIPStarRE.QPBT.MagicSquareRigidity.applyOperatorToState_leftTensor_naimarkInflation, MIPStarRE.QPBT.MagicSquareRigidity.applyOperatorToState_rightTensor_naimarkInflation, MIPStarRE.QPBT.MagicSquareRigidity.norm_compressed_varObs_sub_eq, MIPStarRE.QPBT.MagicSquareRigidity.norm_varObs_sub_sq_le_two_mul_defect, MIPStarRE.QPBT.MagicSquareRigidity.norm_leak_varObs_sq_le_A, MIPStarRE.QPBT.MagicSquareRigidity.norm_leak_varObs_sq_le_B}
	\leanok
	\uses{def:ms-variable-agreement, def:consistent-measurement, def:povm-distance, thm:naimark, lem:povm-to-obs}
	None of the following is a named statement of the source article; each isolates one step of the proofs of Lemma~\ref{lem:ms-variable-agreement-zero} and Lemma~\ref{lem:ms-varobs-cross-player}. Let $\mathcal S$ be a strategy for $\mathfrak{G}^{\mathrm{MS}}$ and $j$ a cell. (i) $\Delta_j(\mathcal S) \ge 0$, and $\Delta_j(\mathcal S)$ equals the two-term sum displayed in Definition~\ref{def:ms-variable-agreement}. (ii) The action of an operator on a state is additive over finite sums of operators, and two operators whose matrix-vector products on a state agree act alike on that state; consequently a measurement consistent on a state in the sense of Definition~\ref{def:consistent-measurement} has each of its postprocessed effects acting alike on the two tensor factors of that state. (iii) On the dilated state of Theorem~\ref{thm:naimark}, the dilated state map is additive over differences, an inflated operator placed on either factor acts as the original operator does on the original state, and hence the ground compressions of the two players' dilated $\mathsf{Variable}_j$ reflections differ on the dilated state exactly as the original binary observables differ on $\ket{\psi}$. (iv) That difference has squared norm at most $2\,\Delta_j(\mathcal S)$, by Lemma~\ref{lem:povm-to-obs} with the unit-modulus weights $1$ and $-1$. (v) For a strategy of value at least $1-\eps$, each player's leakage term has squared norm at most $144\,\eps$.
\end{lemma}
\begin{proof}
	\leanok
	\uses{def:ms-variable-agreement, lem:povm-to-obs, thm:naimark}
	Item~(i) is immediate from the definition. In item~(ii) the first two assertions are the linearity of the map $M \mapsto M\ket{\varphi}$ and the definition of that map; the third follows by writing a postprocessed effect as the finite sum of the effects in its fibre. In item~(iii) the first two assertions hold coordinatewise on the dilated space, and the third is obtained by inserting the ground projection on the opposite factor and cancelling it. Item~(iv) is Lemma~\ref{lem:povm-to-obs} for the two-outcome alphabet. Item~(v) combines the leakage estimate for the dilated observable of a cell with the bound on the cell-consistency mass of a strategy of value $1-\eps$, every cell being incident to some constraint.
\end{proof}

\begin{corollary}[Rigidity for symmetric consistent strategies]\label{cor:ms-rigidity-symmetric-consistent}
	\lean{MIPStarRE.QPBT.exists_ms_rigidity_of_symmetric_consistent}
	\leanok
	\uses{thm:ms-rigidity, def:ms-variable-agreement, lem:ms-variable-agreement-zero, def:consistent-strategy, def:spcc, def:symmetric-game}
	Let $\eps \ge 0$ and let $\mathcal S = (\ket{\psi}, M)$ be a symmetric consistent strategy for $\mathfrak{G}^{\mathrm{MS}}$ in the sense of Definition~\ref{def:consistent-strategy}, so that $\mathcal S$ is projective and every measurement $\{M^x_a\}$ is consistent on $\ket{\psi}$; every SPCC strategy (Definition~\ref{def:spcc}) is of this form. Suppose that $\mathcal S$ succeeds in $\mathfrak{G}^{\mathrm{MS}}$ with probability at least $1-\eps$. Then the conclusion of Theorem~\ref{thm:ms-rigidity} holds with $\eta = \sqrt\eps$.
	This recovers the source's error scale for the completed binary effects of
	Definition~\ref{def:ms-variable-agreement}; comparison with the
	prescribed-answer effects remains a separate construction.
\end{corollary}
\begin{proof}
	\leanok
	\uses{thm:ms-rigidity, lem:ms-variable-agreement-zero}
	By Lemma~\ref{lem:ms-variable-agreement-zero} the agreement distances $\Delta_1(\mathcal S)$ and $\Delta_5(\mathcal S)$ vanish, so Theorem~\ref{thm:ms-rigidity} applies with $\delta = 0$ and $\eta = \sqrt\eps + \sqrt{0} = \sqrt\eps$. Projectivity is part of Definition~\ref{def:consistent-strategy} and is not otherwise used.
\end{proof}

% Formalization-support nodes for the Lean proof of Theorem~\ref{thm:ms-rigidity}, one
% for each module of the assembly. None of the statements below is a named result of the
% source article; together they record the public declarations that the Lean proof of
% the theorem rests on, and the proof of the theorem lists them in its \uses.

\begin{lemma}[Formalization support for the spectral algebra of the controlled swaps]\label{lem:ms-rigidity-support-swap-algebra}
	\lean{MIPStarRE.QPBT.MagicSquareRigidity.reflectionEffect_zero_mul_reflection, MIPStarRE.QPBT.MagicSquareRigidity.reflectionEffect_one_mul_reflection, MIPStarRE.QPBT.MagicSquareRigidity.reflectionEffect_zero_mul_X_sub_X_mul_one, MIPStarRE.QPBT.MagicSquareRigidity.X_mul_reflectionEffect_one_mul_X_sub_zero, MIPStarRE.QPBT.kroneckerMap_sub_left, MIPStarRE.QPBT.kroneckerMap_sub_right, MIPStarRE.QPBT.heteroKron_sub_left, MIPStarRE.QPBT.heteroKron_sub_right}
	\leanok
	\uses{def:tensor-product-strategy}
	None of the following is a named statement of the source article; each isolates one step of the Lean proof of Theorem~\ref{thm:ms-rigidity}. Let $X$ and $Z$ be $\pm 1$-valued observables on a finite-dimensional space and let $P^Z_0 = \tfrac12(I + Z)$ and $P^Z_1 = \tfrac12(I - Z)$ be the spectral effects of $Z$. Then $P^Z_0 Z = P^Z_0$ and $P^Z_1 Z = -P^Z_1$, and
	\[
		P^Z_0 X - X P^Z_1 = \tfrac12 (XZ + ZX)\;, \qquad
		X P^Z_1 X - P^Z_0 = -\tfrac12\, X (XZ + ZX)\;,
	\]
	so that both differences vanish exactly when $X$ and $Z$ anticommute. The tensor placement of Definition~\ref{def:tensor-product-strategy} respects differences in either factor, also for rectangular matrices, which is the form needed when a factor is the matrix of an isometry between distinct index sets.
\end{lemma}
\begin{proof}
	\leanok
	\uses{def:tensor-product-strategy}
	The first two identities follow from $Z^2 = I$; the two displayed ones by expanding the spectral effects and using $X^2 = I$. The tensor identities hold entrywise.
\end{proof}

\begin{lemma}[Formalization support for the large-error regime]\label{lem:ms-rigidity-support-constants}
	\lean{MIPStarRE.QPBT.MagicSquareRigidity.registerFiberEmbeddingMap, MIPStarRE.QPBT.MagicSquareRigidity.registerFiberEmbedding, MIPStarRE.QPBT.MagicSquareRigidity.registerFiberEmbedding_apply, MIPStarRE.QPBT.MagicSquareRigidity.conjTranspose_mul_le_one_neg, MIPStarRE.QPBT.MagicSquareRigidity.opFamilyDistSq_uniform_unit_binary_le, MIPStarRE.QPBT.MagicSquareRigidity.opDistSq_uniform_unit_le}
	\leanok
	\uses{def:povm-distance, def:EPR}
	None of the following is a named statement of the source article; each isolates one step of the Lean proof of Theorem~\ref{thm:ms-rigidity}. Once $\sqrt\eps + \sqrt\delta \ge 1$ the conclusion of Theorem~\ref{thm:ms-rigidity} is carried by estimates available for any witness: two unit vectors are at distance at most $2$; the squared state-dependent distance of Definition~\ref{def:povm-distance} between two contractions on a unit vector is at most $4$, so a two-outcome family has squared distance at most $8$ under the one-point distribution. A witness exists for every strategy, namely the embedding of a local space as the fibre of the extracted register over one fixed basis vector, which is an isometry.
\end{lemma}
\begin{proof}
	\leanok
	\uses{def:povm-distance}
	Each estimate is the triangle inequality together with $\|Mv\| \le \|v\|$ for a contraction $M$, summed over the two outcomes. The fibre embedding preserves norms coordinatewise.
\end{proof}

\begin{lemma}[Formalization support for the ideal target of the rigidity conclusion]\label{lem:ms-rigidity-support-ideal-target}
	\lean{MIPStarRE.QPBT.tauObservable_zero, MIPStarRE.QPBT.sum_pauliProj_eq_one, MIPStarRE.QPBT.posSemidef_pauliProj, MIPStarRE.QPBT.posSemidef_sum_ite_pauliProj, MIPStarRE.QPBT.sum_ite_pauliProj_nonneg, MIPStarRE.QPBT.sum_binary_marginal_pauliProj, MIPStarRE.QPBT.pauliMarginalMeasurement, MIPStarRE.QPBT.pauliMarginalMeasurement_effect, MIPStarRE.QPBT.conjTranspose_mul_le_one_sum_ite_pauliProj, MIPStarRE.QPBT.norm_reindexState_prodShuffle_vecTensor_eprState}
	\leanok
	\uses{def:generalized-pauli, def:EPR}
	None of the following is a named statement of the source article; each isolates one step of the Lean proof of Theorem~\ref{thm:ms-rigidity}. The Pauli basis projectors of one kind on a two-qubit register are positive semidefinite and sum to the identity, so any marginal of them is a positive contraction and the two-outcome marginal family is a measurement; and the ideal state $\ket{\mathrm{EPR}_2}^{\ot 2} \ot \ket{\mathrm{aux}}$ is a unit vector whenever $\ket{\mathrm{aux}}$ is.
\end{lemma}
\begin{proof}
	\leanok
	\uses{def:generalized-pauli, lem:pauli-observable-expansion}
	Completeness is the expansion of Lemma~\ref{lem:pauli-observable-expansion} evaluated at the zero label, where the generalized Pauli observable of Definition~\ref{def:generalized-pauli} is the identity. Positivity is the rank-one factorization of each projector on the corresponding normalized Pauli eigenvector. The norm of the ideal state is the product of the norms of its factors.
\end{proof}

\begin{lemma}[Formalization support for the two-qubit controlled-swap isometry]\label{lem:ms-rigidity-support-two-qubit-swap}
	\lean{MIPStarRE.QPBT.MagicSquareRigidity.sum_pi_fin_two, MIPStarRE.QPBT.MagicSquareRigidity.sum_norm_sq_binarySwapComponent, MIPStarRE.QPBT.MagicSquareRigidity.sum_norm_sq_binarySwapComponent_eq, MIPStarRE.QPBT.MagicSquareRigidity.twoBinarySwapMap, MIPStarRE.QPBT.MagicSquareRigidity.twoBinarySwapMap_apply, MIPStarRE.QPBT.MagicSquareRigidity.twoBinarySwapMap_norm, MIPStarRE.QPBT.MagicSquareRigidity.twoBinarySwapIsometry, MIPStarRE.QPBT.MagicSquareRigidity.twoBinarySwapIsometry_apply, MIPStarRE.QPBT.MagicSquareRigidity.msLocalCellObsA, MIPStarRE.QPBT.MagicSquareRigidity.isBinaryObservable_msLocalCellObsA, MIPStarRE.QPBT.MagicSquareRigidity.msCellObsA_eq_heteroKron, MIPStarRE.QPBT.MagicSquareRigidity.msVarObsA_eq_heteroKron, MIPStarRE.QPBT.MagicSquareRigidity.msVarObsB_eq_heteroKron, MIPStarRE.QPBT.MagicSquareRigidity.msAliceTwoQubitSwapIsometry, MIPStarRE.QPBT.MagicSquareRigidity.msBobTwoQubitSwapIsometry}
	\leanok
	\uses{def:ms-game, thm:naimark}
	None of the following is a named statement of the source article; each isolates one step of the Lean proof of Theorem~\ref{thm:ms-rigidity}. Let $(X_1, Z_1)$ and $(X_2, Z_2)$ be two pairs of $\pm 1$-valued observables on one local space. The map sending a vector $\psi$ to the family indexed by the register labels $e \in \F_2^{\{0,1\}}$ whose entry at $e$ is
	\[
		X_2^{e(1)} P^{Z_2}_{e(1)} X_1^{e(0)} P^{Z_1}_{e(0)} \psi
	\]
	is an isometry into the two-qubit register; no commutation between the two pairs is needed for that, and the first pair, applied innermost, carries the register coordinate $0$. Applied to the questionwise Naimark dilation of a strategy for $\mathfrak{G}^{\mathrm{MS}}$ it gives the two players' embeddings, Alice's from her variable reflections at the cells $0$ and $4$ together with her constraint reflections at the cells $1$ and $3$, Bob's from his variable reflections at those four cells; each of the three families of reflections is the tensor placement of a local reflection on the corresponding player's factor.
\end{lemma}
\begin{proof}
	\leanok
	\uses{thm:naimark}
	The squared norm of the image is the sum over the register labels of the squared norms of its entries, which the completeness of the spectral effects of each pair collapses to $\|\psi\|^2$, one pair at a time.
\end{proof}

\begin{lemma}[Formalization support for the second logical pair]\label{lem:ms-rigidity-support-second-pair}
	\lean{MIPStarRE.QPBT.MagicSquareRigidity.msCellObsA_second_pair_anticommute, MIPStarRE.QPBT.MagicSquareRigidity.msVarObsB_second_pair_anticommute, MIPStarRE.QPBT.MagicSquareRigidity.msCellObsA_close_msVarObsB_second_x, MIPStarRE.QPBT.MagicSquareRigidity.msCellObsA_close_msVarObsB_second_z, MIPStarRE.QPBT.MagicSquareRigidity.msVarObsB_comm_of_shared_constraint, MIPStarRE.QPBT.MagicSquareRigidity.msVarObsA_comm_msCellObsA_of_shared_constraint}
	\leanok
	\uses{def:ms-game, thm:naimark, def:tensor-product-value}
	None of the following is a named statement of the source article; each isolates one step of the Lean proof of Theorem~\ref{thm:ms-rigidity}. Let $\mathcal S$ be a strategy for $\mathfrak{G}^{\mathrm{MS}}$ of value at least $1 - \eps$, pass to its questionwise Naimark dilation (Theorem~\ref{thm:naimark}), and read all relations on the dilated state. Alice's constraint reflections at the cells $1$ and $3$, read from the first and from the second row, anticommute up to $624\sqrt\eps$, and so do Bob's variable reflections at those two cells. Each of Alice's two constraint reflections agrees with Bob's variable reflection at the same cell up to $12\sqrt\eps$. Two of Bob's variable reflections at cells of one constraint commute up to $96\sqrt\eps$, and one of Alice's variable reflections at a cell of agreement defect at most $\delta$ commutes with one of her constraint reflections at a cell of a shared constraint up to $4\sqrt{864\eps + 6\delta} + 96\sqrt\eps$.
\end{lemma}
\begin{proof}
	\leanok
	\uses{thm:naimark, lem:ms-varobs-cross-player}
	The anticommutation is the solution-group computation carried along the closed path $(1,3) \to (2,6) \to (5,7) \to (3,1)$ of cell pairs, which traverses the exceptional third column once and therefore returns the reversed product with the opposite sign; the path gives the scale $504\sqrt\eps$, to which substituting the two endpoints adds $24\sqrt\eps$ each. The cross-player agreements are the constraint--variable subtests at the cells $1$ and $3$. The commutation relations come from the joint measurability of the reflections read from one constraint of one player, the other player's reflections being brought to that player by the same subtests and, in the mixed case, by Lemma~\ref{lem:ms-varobs-cross-player}.
\end{proof}

\begin{lemma}[Formalization support for the components of a bipartite pair of controlled swaps]\label{lem:ms-rigidity-support-swap-pair}
	\lean{MIPStarRE.QPBT.MagicSquareRigidity.heteroKron_eq_left_mul_right, MIPStarRE.QPBT.MagicSquareRigidity.heteroKron_left_comm_right, MIPStarRE.QPBT.MagicSquareRigidity.isProj_reflectionEffect, MIPStarRE.QPBT.MagicSquareRigidity.contraction_reflectionEffect, MIPStarRE.QPBT.MagicSquareRigidity.contraction_swapFactor, MIPStarRE.QPBT.MagicSquareRigidity.contraction_heteroKron, MIPStarRE.QPBT.MagicSquareRigidity.applyOperatorToState_sub_vec, MIPStarRE.QPBT.MagicSquareRigidity.norm_mul_sub_le, MIPStarRE.QPBT.MagicSquareRigidity.norm_one_add_le, MIPStarRE.QPBT.MagicSquareRigidity.norm_one_sub_le, MIPStarRE.QPBT.MagicSquareRigidity.heteroKron_reflection_sub_zero, MIPStarRE.QPBT.MagicSquareRigidity.heteroKron_reflection_sub_one, MIPStarRE.QPBT.MagicSquareRigidity.heteroKron_reflection_zero_one, MIPStarRE.QPBT.MagicSquareRigidity.heteroKron_reflection_one_zero, MIPStarRE.QPBT.MagicSquareRigidity.isBinaryObservable_heteroKron_right, MIPStarRE.QPBT.MagicSquareRigidity.norm_swapComponent_01_le, MIPStarRE.QPBT.MagicSquareRigidity.norm_swapComponent_10_le, MIPStarRE.QPBT.MagicSquareRigidity.norm_apply_half_smul, MIPStarRE.QPBT.MagicSquareRigidity.heteroKron_split_right, MIPStarRE.QPBT.MagicSquareRigidity.norm_swapComponent_diag_sub_le}
	\leanok
	\uses{def:povm-distance}
	None of the following is a named statement of the source article; each isolates one step of the Lean proof of Theorem~\ref{thm:ms-rigidity}. Let $(X_{\mathrm A}, Z_{\mathrm A})$ and $(X_{\mathrm B}, Z_{\mathrm B})$ be pairs of $\pm 1$-valued observables on the two players' spaces, let $A_b = X_{\mathrm A}^b P^{Z_{\mathrm A}}_b$ and $B_c = X_{\mathrm B}^c P^{Z_{\mathrm B}}_c$ be the residual factors of the two controlled swaps, and let $v$ be any vector on the bipartite space. If $b \ne c$ then $\|(A_b \ot B_c) v\|$ is at most half of $\|(Z_{\mathrm A} \ot I - I \ot Z_{\mathrm B}) v\|$, and the two diagonal components $(A_0 \ot B_0)v$ and $(A_1 \ot B_1)v$ differ by at most the sum of the two cross-player distances of the pair on $v$ and half the norm of $(X_{\mathrm A} Z_{\mathrm A} + Z_{\mathrm A} X_{\mathrm A}) \ot I$ on $v$. Neither estimate uses Bob's anticommutator, and neither assumes $v$ normalized.
\end{lemma}
\begin{proof}
	\leanok
	\uses{def:povm-distance}
	Both estimates factor the component through the spectral effects, which are positive contractions, and trade Alice's reflection for Bob's at the cost of the cross-player distance. The diagonal comparison in addition conjugates the negative spectral effect by the shift reflection, which returns the positive one up to half the anticommutator, by the spectral identities above.
\end{proof}

\begin{lemma}[Formalization support for the transport of the second logical pair]\label{lem:ms-rigidity-support-swap-transport}
	\lean{MIPStarRE.QPBT.MagicSquareRigidity.NormBoundedBy, MIPStarRE.QPBT.MagicSquareRigidity.NormBoundedBy.of_contraction, MIPStarRE.QPBT.MagicSquareRigidity.NormBoundedBy.mono, MIPStarRE.QPBT.MagicSquareRigidity.NormBoundedBy.add, MIPStarRE.QPBT.MagicSquareRigidity.NormBoundedBy.sub, MIPStarRE.QPBT.MagicSquareRigidity.NormBoundedBy.mul, MIPStarRE.QPBT.MagicSquareRigidity.normBoundedBy_leftTensor, MIPStarRE.QPBT.MagicSquareRigidity.normBoundedBy_rightTensor, MIPStarRE.QPBT.MagicSquareRigidity.norm_leftTransport_le, MIPStarRE.QPBT.MagicSquareRigidity.norm_rightTransport_le, MIPStarRE.QPBT.MagicSquareRigidity.norm_mul_le_commutator_add, MIPStarRE.QPBT.MagicSquareRigidity.commutator_heteroKron_eq, MIPStarRE.QPBT.MagicSquareRigidity.norm_heteroKron_left_le, MIPStarRE.QPBT.MagicSquareRigidity.norm_heteroKron_right_le, MIPStarRE.QPBT.MagicSquareRigidity.norm_comm_swapFactor_left_le, MIPStarRE.QPBT.MagicSquareRigidity.norm_comm_swapFactor_right_le, MIPStarRE.QPBT.MagicSquareRigidity.norm_defect_transported_le}
	\leanok
	\uses{def:povm-distance}
	None of the following is a named statement of the source article; each isolates one step of the Lean proof of Theorem~\ref{thm:ms-rigidity}. Write $v_b = \bigl(X_1^b P^{Z_1}_b \ot Y_1^b P^{W_1}_b\bigr)\psi$ for the residual of the two-qubit controlled swap after its first logical pair. Every hypothesis of the second logical pair that holds on $\psi$ holds on $v_b$, at the cost of the commutation defects between the two pairs and of the cross-player agreement of the first pair. Two mechanisms carry this: a commutator evaluated on $Z\psi$ is controlled by the same commutator on $\psi$ once Alice's reflection is traded for Bob's, which commutes with every operator on Alice's factor; and an operator may be moved across a residual factor at the cost of its commutator with that factor.
\end{lemma}
\begin{proof}
	\leanok
	\uses{def:povm-distance}
	Both mechanisms are the triangle inequality applied to $MN = NM + [M, N]$, together with the fact that the reflections and the residual factors are contractions, so that they do not increase the norm of the vector they are applied to.
\end{proof}

\begin{lemma}[Formalization support for the joint state estimate]\label{lem:ms-rigidity-support-joint-state}
	\lean{MIPStarRE.QPBT.MagicSquareRigidity.applyOperatorToState_coord, MIPStarRE.QPBT.MagicSquareRigidity.isometryMatrix_of_family, MIPStarRE.QPBT.MagicSquareRigidity.isometryTensor_family_apply, MIPStarRE.QPBT.MagicSquareRigidity.reindexState_prodShuffle_vecTensor_eprState_apply, MIPStarRE.QPBT.MagicSquareRigidity.card_two_qubit_register, MIPStarRE.QPBT.MagicSquareRigidity.norm_sq_eq_sum_components, MIPStarRE.QPBT.MagicSquareRigidity.norm_le_card_mul_of_components, MIPStarRE.QPBT.MagicSquareRigidity.swapFactor, MIPStarRE.QPBT.MagicSquareRigidity.swapFactor_zero, MIPStarRE.QPBT.MagicSquareRigidity.swapFactor_one, MIPStarRE.QPBT.MagicSquareRigidity.contraction_swapFactor', MIPStarRE.QPBT.MagicSquareRigidity.norm_anticommutator_comm_first_le, MIPStarRE.QPBT.MagicSquareRigidity.norm_secondPair_cross_le, MIPStarRE.QPBT.MagicSquareRigidity.norm_secondPair_anticommutator_le, MIPStarRE.QPBT.MagicSquareRigidity.pi_fin_two_ext, MIPStarRE.QPBT.MagicSquareRigidity.sqrt_card_two_qubit_register, MIPStarRE.QPBT.MagicSquareRigidity.exists_residual_of_two_pairs, MIPStarRE.QPBT.MagicSquareRigidity.vecTensor_smul_right, MIPStarRE.QPBT.MagicSquareRigidity.reindexState_smul, MIPStarRE.QPBT.MagicSquareRigidity.exists_unit_residual}
	\leanok
	\uses{def:EPR, def:povm-distance}
	None of the following is a named statement of the source article; each isolates one step of the Lean proof of Theorem~\ref{thm:ms-rigidity}. Suppose each player carries two pairs of $\pm 1$-valued observables on a bipartite state, and suppose that both of Alice's pairs anticommute, that all four pairs of corresponding reflections agree across the players, and that the two pairs of each player commute, each up to $\eta$ on that state. Then the tensor of the two two-qubit controlled-swap embeddings carries the state to within an explicit multiple of $\eta$ of $\ket{\mathrm{EPR}_2}^{\ot 2} \ot w$ for some residual bipartite vector $w$, and $w$ may be replaced by a unit vector at twice that cost.
\end{lemma}
\begin{proof}
	\leanok
	\uses{def:EPR}
	The image of the state is the family of its components at the sixteen pairs of register labels. Each component factors as a component of the second pair applied to a component of the first, so the two estimates above apply twice, once on the state and once on the residual of the first pair, whose second-pair hypotheses the transport above supplies. Every component off the diagonal is small and every component on it is close to the one at the zero label; summing the sixteen labels gives the estimate, and normalizing the residual costs its distance from a unit vector, which the estimate itself bounds.
\end{proof}

\begin{lemma}[Formalization support for the two-qubit Pauli marginals]\label{lem:ms-rigidity-support-intertwine}
	\lean{MIPStarRE.QPBT.MagicSquareRigidity.swapFactor_mul_obs, MIPStarRE.QPBT.MagicSquareRigidity.sum_swapFactor_conjTranspose_mul, MIPStarRE.QPBT.MagicSquareRigidity.sum_leftTensor_swapFactor, MIPStarRE.QPBT.MagicSquareRigidity.sum_rightTensor_swapFactor, MIPStarRE.QPBT.MagicSquareRigidity.twoQubitPauliObs, MIPStarRE.QPBT.MagicSquareRigidity.binTrace_zmod_two, MIPStarRE.QPBT.MagicSquareRigidity.sum_ite_pauliProj_eq_reflectionEffect, MIPStarRE.QPBT.MagicSquareRigidity.applyOperatorToState_leftTensor_coord, MIPStarRE.QPBT.MagicSquareRigidity.applyOperatorToState_register_coord, MIPStarRE.QPBT.MagicSquareRigidity.twoQubitPauliObs_Z_apply, MIPStarRE.QPBT.MagicSquareRigidity.twoQubitPauliObs_X_apply, MIPStarRE.QPBT.MagicSquareRigidity.sum_twoSwapFactor_conjTranspose_mul, MIPStarRE.QPBT.MagicSquareRigidity.sum_rightTensor_twoSwapFactor}
	\leanok
	\uses{def:generalized-pauli}
	None of the following is a named statement of the source article; each isolates one step of the Lean proof of Theorem~\ref{thm:ms-rigidity}. The residual factors of a controlled swap absorb the phase reflection with the sign of the register label and form a complete family, on one and on two qubits alike, and this remains so after tensoring with the other player's factor. The marginal over the second register coordinate of the two-qubit Pauli basis projectors of one kind is the spectral effect of the corresponding two-qubit Pauli observable at the label concentrated on the first coordinate.
\end{lemma}
\begin{proof}
	\leanok
	\uses{def:generalized-pauli, lem:pauli-observable-expansion}
	The absorption identities are the spectral identities above. Completeness is the sum over the register labels of the products of the residual factors with their adjoints, in which the spectral effects of each pair sum to the identity. The marginal identity is the two-term expansion of Lemma~\ref{lem:pauli-observable-expansion} in the binary trace pairing, the marginal summing the projectors over the label of the second coordinate.
\end{proof}

\begin{lemma}[Formalization support for the transport of the first logical pair]\label{lem:ms-rigidity-support-marginals}
	\lean{MIPStarRE.QPBT.MagicSquareRigidity.heteroKron_left_mul_apply, MIPStarRE.QPBT.MagicSquareRigidity.twoQubitPauliObs_Z_mul_apply, MIPStarRE.QPBT.MagicSquareRigidity.add_single_zero_add_single_zero, MIPStarRE.QPBT.MagicSquareRigidity.twoQubitPauliObs_X_mul_apply, MIPStarRE.QPBT.MagicSquareRigidity.conjTranspose_mul_of_blocks, MIPStarRE.QPBT.MagicSquareRigidity.twoSwapMatrix, MIPStarRE.QPBT.MagicSquareRigidity.twoSwapMatrix_apply, MIPStarRE.QPBT.MagicSquareRigidity.twoSwap_matrix_intertwines_Z, MIPStarRE.QPBT.MagicSquareRigidity.twoSwap_shift_defect_apply, MIPStarRE.QPBT.MagicSquareRigidity.swapFactor_shift_defect_zero, MIPStarRE.QPBT.MagicSquareRigidity.swapFactor_shift_defect_one, MIPStarRE.QPBT.MagicSquareRigidity.sum_twoSwapFactor_gram, MIPStarRE.QPBT.MagicSquareRigidity.twoSwap_shift_defect_conjTranspose_mul, MIPStarRE.QPBT.MagicSquareRigidity.reflectionEffect_defect_eq, MIPStarRE.QPBT.MagicSquareRigidity.reflectionEffect_defect_conjTranspose_mul, MIPStarRE.QPBT.MagicSquareRigidity.leftTensor_isometryTensor_defect_eq, MIPStarRE.QPBT.MagicSquareRigidity.rightTensor_isometryTensor_defect_eq, MIPStarRE.QPBT.MagicSquareRigidity.norm_sq_toEuclideanLin_kron_left, MIPStarRE.QPBT.MagicSquareRigidity.norm_sq_toEuclideanLin_kron_right, MIPStarRE.QPBT.MagicSquareRigidity.norm_sq_leftTensor_isometryTensor_defect, MIPStarRE.QPBT.MagicSquareRigidity.norm_sq_rightTensor_isometryTensor_defect, MIPStarRE.QPBT.MagicSquareRigidity.reflectionEffect_heteroKron_left, MIPStarRE.QPBT.MagicSquareRigidity.reflectionEffect_heteroKron_right, MIPStarRE.QPBT.MagicSquareRigidity.dilatedEffect_var_A_eq, MIPStarRE.QPBT.MagicSquareRigidity.dilatedEffect_var_B_eq, MIPStarRE.QPBT.MagicSquareRigidity.isometryMatrix_msAliceTwoQubitSwapIsometry, MIPStarRE.QPBT.MagicSquareRigidity.isometryMatrix_msBobTwoQubitSwapIsometry, MIPStarRE.QPBT.MagicSquareRigidity.msJointAnticommutatorA, MIPStarRE.QPBT.MagicSquareRigidity.msJointAnticommutatorB, MIPStarRE.QPBT.MagicSquareRigidity.norm_msJointAnticommutatorA_le, MIPStarRE.QPBT.MagicSquareRigidity.norm_msJointAnticommutatorB_le, MIPStarRE.QPBT.MagicSquareRigidity.norm_ms_effect_defect_A_X, MIPStarRE.QPBT.MagicSquareRigidity.norm_ms_effect_defect_A_Z, MIPStarRE.QPBT.MagicSquareRigidity.norm_ms_effect_defect_B_X, MIPStarRE.QPBT.MagicSquareRigidity.norm_ms_effect_defect_B_Z}
	\leanok
	\uses{def:generalized-pauli, def:ms-game}
	None of the following is a named statement of the source article; each isolates one step of the Lean proof of Theorem~\ref{thm:ms-rigidity}. Let $V$ be the matrix of the two-qubit controlled swap of two pairs $(X_1, Z_1)$, $(X_2, Z_2)$, and let $\tau^X$, $\tau^Z$ be the two-qubit Pauli observables on the first register coordinate. Then $\tau^Z V = V Z_1$ exactly, and the shift transport defect $D = \tau^X V - V X_1$ satisfies $D^\dagger D = \tfrac12\, G^\dagger G$ with $G = X_1 Z_1 + Z_1 X_1$ the anticommutator of the first pair. The spectral effects of the two observables inherit these relations with a further factor $\tfrac14$, and tensoring with the other player's embedding preserves them. On the dilated Magic Square strategy this carries each player's $\mathsf{Variable}_1$ and $\mathsf{Variable}_5$ effects to the corresponding Pauli marginals, exactly for the phase and up to the norm of that player's joint anticommutator for the shift.
\end{lemma}
\begin{proof}
	\leanok
	\uses{def:generalized-pauli}
	The residual factors of the second controlled swap form a complete family and therefore drop out of $D^\dagger D$, which leaves the one-qubit computation of the spectral identities above. The factor $\tfrac14$ for the spectral effects is the square of the factor $\tfrac12$ relating an effect to its observable. The other player's embedding is an isometry, so it preserves the squared norms, and the tensor placement respects differences by the identities above.
\end{proof}

\begin{lemma}[Formalization support for the assembly of the small-error regime]\label{lem:ms-rigidity-support-assembly}
	\lean{MIPStarRE.QPBT.MagicSquareRigidity.msRigidityDefect, MIPStarRE.QPBT.MagicSquareRigidity.msRigidityDefect_nonneg, MIPStarRE.QPBT.MagicSquareRigidity.sqrt_consistency_le_msRigidityDefect, MIPStarRE.QPBT.MagicSquareRigidity.sqrt_value_le_msRigidityDefect, MIPStarRE.QPBT.MagicSquareRigidity.comm_le_msRigidityDefect, MIPStarRE.QPBT.MagicSquareRigidity.msRigidityDefect_le, MIPStarRE.QPBT.MagicSquareRigidity.heteroKron_left_commutator, MIPStarRE.QPBT.MagicSquareRigidity.heteroKron_left_anticommutator, MIPStarRE.QPBT.MagicSquareRigidity.heteroKron_right_commutator, MIPStarRE.QPBT.MagicSquareRigidity.ms_dilated_state_residual, MIPStarRE.QPBT.MagicSquareRigidity.ms_dilated_state_estimate}
	\leanok
	\uses{def:ms-game, def:ms-variable-agreement, thm:naimark, def:EPR}
	None of the following is a named statement of the source article; each isolates one step of the Lean proof of Theorem~\ref{thm:ms-rigidity}. Put $\Delta(\eps, \delta) = 4\sqrt{864\eps + 6\delta} + 624\sqrt\eps$. Let $\mathcal S$ be a strategy for $\mathfrak{G}^{\mathrm{MS}}$ of value at least $1 - \eps$ with $\Delta_1(\mathcal S), \Delta_5(\mathcal S) \le \delta$. On the questionwise Naimark dilation of $\mathcal S$ every hypothesis of the joint state estimate above holds at the scale $\Delta(\eps, \delta)$, and $\Delta(\eps, \delta) \le 744(\sqrt\eps + \sqrt\delta)$. Consequently the tensor of the two players' embeddings carries the dilated state to within $172608\,(\sqrt\eps + \sqrt\delta)$ of $\ket{\mathrm{EPR}_2}^{\ot 2} \ot \ket{\mathrm{aux}}$ for a unit vector $\ket{\mathrm{aux}}$.
\end{lemma}
\begin{proof}
	\leanok
	\uses{thm:naimark, lem:ms-varobs-cross-player, def:ms-variable-agreement}
	The anticommutation of each pair is the solution-group computation, the cross-player agreement at the cells of $\mathsf{Variable}_1$ and $\mathsf{Variable}_5$ is Lemma~\ref{lem:ms-varobs-cross-player} and at the two remaining cells it is the constraint--variable subtests, and the commutation of the two pairs of one player is the second-pair node above; each is bounded by $\Delta(\eps, \delta)$ using $\sqrt{x + y} \le \sqrt x + \sqrt y$ and the monotonicity of the square root. The final estimate is the joint state estimate at $116\,\Delta(\eps, \delta)$, doubled by the normalization of the residual, and $232 \cdot 744 = 172608$.
\end{proof}

\begin{lemma}[Formalization support for the regime split and the transfer to the given strategy]\label{lem:ms-rigidity-support-transfer}
	\lean{MIPStarRE.QPBT.ms_rigidity_coarse_bounds, MIPStarRE.QPBT.exists_ms_rigidity_of_one_le_sqrt_add, MIPStarRE.QPBT.norm_applyOperatorToState_sub_le_of_close, MIPStarRE.QPBT.sum_sq_le_two_mul, MIPStarRE.QPBT.ms_dilated_operator_distance_A, MIPStarRE.QPBT.ms_dilated_operator_distance_B, MIPStarRE.QPBT.ms_dilated_anticommutator_A, MIPStarRE.QPBT.ms_dilated_anticommutator_B}
	\leanok
	\uses{def:ms-game, def:povm-distance, thm:naimark}
	None of the following is a named statement of the source article; each isolates one step of the Lean proof of Theorem~\ref{thm:ms-rigidity}. For $C \ge 8$ and $\sqrt\eps + \sqrt\delta \ge 1$ the seven conclusions of Theorem~\ref{thm:ms-rigidity} hold at the scale $C(\sqrt\eps + \sqrt\delta)$ for the witness built from the fibre embedding. In the complementary regime each of the seven conclusions passes from the questionwise Naimark dilation to the given strategy: a bound $\kappa$ on the norm, on the extracted state, of the difference between a transported effect and its ideal Pauli marginal, together with a bound $\eta$ on the distance between the extracted state and the ideal one, bounds the corresponding squared distance of Definition~\ref{def:povm-distance} by $2(\kappa + 2\eta)^2$; and a bound $\kappa$ on the norm, on the extracted state, of that player's joint anticommutator bounds his or her anticommutator distance by $(\kappa + 2\eta)^2$.
\end{lemma}
\begin{proof}
	\leanok
	\uses{def:povm-distance, thm:naimark}
	The first assertion is the coarse estimates above. For the second, an inflated operator acts on the dilated state as the original operator acts on the given state, so a norm bound on the two outcome vectors bounds each summand of the state-dependent distance; summing the two outcomes of the binary alphabet contributes the factor $2$, and moving the comparison from the extracted state to the ideal one contributes $2\eta$ by the triangle inequality.
\end{proof}

The converse construction shows that any pair of anticommuting binary observables yields a perfect strategy for the Magic Square game.

\begin{theorem}\label{thm:ms-from-ac}
	\lean{MIPStarRE.QPBT.obsOf, MIPStarRE.QPBT.exists_ms_perfect_strategy_of_anticommuting}
	\leanok
	\uses{def:ms-game, def:consistent-measurement, def:spcc, def:EPR, def:tensor-product-value}
	Let $A = \{A_b\}_{b \in \F_2}$ and $B = \{B_b\}_{b \in \F_2}$ be two-outcome projective measurements acting on $(\C^q)^{\ot n}$ that are consistent on $\ket{\mathrm{EPR}_q}^{\ot n}$ (Definition~\ref{def:consistent-measurement}), and let $\mathcal O_A = A_0 - A_1$ and $\mathcal O_B = B_0 - B_1$ be the corresponding $\pm 1$-valued observables. Suppose that $\mathcal O_A \mathcal O_B = - \mathcal O_B \mathcal O_A$. Then there exists a symmetric strategy $\mathcal S = (\psi, M)$ for the Magic Square game with the following properties:
	\begin{enumerate}
		\item $\mathcal S$ is an SPCC (symmetric, projective, consistent, commuting) strategy of value $1$;
		\item the state is $\ket{\psi} = \ket{\mathrm{EPR}_q}^{\ot n} \ot \ket{\mathrm{EPR}_2}$;
		\item for $b \in \{0,1\}$, $M^{\mathsf{Variable}_1}_b = A_b \ot I$ and $M^{\mathsf{Variable}_5}_b = B_b \ot I$.
	\end{enumerate}
\end{theorem}
\begin{proof}
	\leanok
	\uses{def:consistent-measurement}
	The strategy is based on the canonical two-qubit operator solution of the Magic Square, with the state $\ket{\psi} = \ket{\mathrm{EPR}_q}^{\ot n} \ot \ket{\mathrm{EPR}_2}$ and the cells of the grid assigned the $\pm 1$-valued observables
	\[
		\begin{array}{ccc}
			\mathcal O_A \ot I & I \ot \sigma^X & \mathcal O_A \ot \sigma^X \\[2pt]
			I \ot \sigma^Z & \mathcal O_B \ot I & \mathcal O_B \ot \sigma^Z \\[2pt]
			\mathcal O_A \ot \sigma^Z & \mathcal O_B \ot \sigma^X & \mathcal O_A \mathcal O_B \ot \sigma^Z \sigma^X
		\end{array}
	\]
	on $(\C^q)^{\ot n} \ot \C^2$: on question $\mathsf{Variable}_j$ a player measures the two-outcome projective measurement of the observable in cell $j$, and on question $\mathsf{Constraint}_i$ a player simultaneously performs the three measurements of the corresponding row or column; item~3 holds by construction. Anticommutation of $\mathcal O_A$ and $\mathcal O_B$ gives
	\begin{equation}\label{eq:am-i-hermitian}
		\mathcal O_A \mathcal O_B \ot \sigma^Z \sigma^X
		= - \mathcal O_A \mathcal O_B \ot \sigma^X \sigma^Z
		= \mathcal O_B \mathcal O_A \ot \sigma^X \sigma^Z
		= (\mathcal O_A \mathcal O_B \ot \sigma^Z \sigma^X)^\dagger\;,
	\end{equation}
	so the bottom-right cell is Hermitian, and it squares to the identity; hence every cell is a genuine $\pm 1$-valued observable. The three observables in each row and each column commute pairwise --- for the third row and third column this again uses Equation~\eqref{eq:am-i-hermitian} --- so the $\mathsf{Constraint}$ measurements are well defined, and the strategy is symmetric, projective, and commuting. A two-outcome projective measurement is consistent on a state if and only if its observable $\mathcal O$ satisfies $\mathcal O \ot I \ket{\psi} = I \ot \mathcal O \ket{\psi}$; since $A$ and $B$ are consistent on $\ket{\mathrm{EPR}_q}^{\ot n}$ by hypothesis and $\sigma^X$, $\sigma^Z$ are consistent on $\ket{\mathrm{EPR}_2}$, each cell observable is a product of consistent, commuting factors and is therefore consistent on $\ket{\psi}$, and consistency of the $\mathsf{Constraint}$ measurements follows from the cellwise consistencies together with commutativity. Finally, the observables in each row and each of the first two columns multiply to $I$, while those of the third column multiply to $-I$; for the measurement $\{S_0, S_1\}$ obtained by measuring a $\mathsf{Constraint}_i$ triple and outputting the sum of the three bits, the observable $S_0 - S_1$ equals the product of the three cell observables, hence $\pm I$, which forces the three answer bits to sum to $0$ (respectively, $1$ for $\mathsf{Constraint}_6$) with certainty. Together with consistency this shows that every check of the Magic Square game is passed with probability $1$, so the strategy has value $1$.
\end{proof}

\section{The Pauli basis test}
% Paper label sec:qld-game kept for cross-reference fidelity.
\label{sec:qld-game}

The Pauli basis test is the quantum low individual degree test of Natarajan and Vidick, in the slightly modified form introduced by Natarajan and Wright~\cite{Natarajan2018NEEXP}. Informally, the test certifies that the players share the state $\ket{\mathrm{EPR}_q}^{\ot M}$, where $M = 2^m$. On a question of type $(\mathsf{Pauli}, W)$ with $W \in \{X,Z\}$, a player is expected to measure the generalized Pauli basis POVM $\{\tau^W_h\}_{h \in \F_q^M}$ of Chapter~\ref{chap:qpbt-algebra} and report the outcome $h$. Most subtests probe limited information about $h$: given a point $u_W \in \F_q^m$, a player is expected to report the evaluation $g_h(u_W)$ of the low-degree encoding $g_h$ of $h$, and for a fixed basis $W$ these probes are checked against each other by the classical low individual degree game of Section~\ref{sec:ld-verifier}. Consistency \emph{between} the two bases is checked through two-outcome probes: given additionally $r_W \in \F_q$, a player reports $\tr(g_h(u_W)\, r_W) \in \F_2$. The two binary measurements obtained in this way for $W = X$ and $W = Z$ either commute or anticommute, according to the value $\gamma \in \F_2$ of Equation~\eqref{eq:gamma-value} below. When $\gamma = 0$ the verifier checks commutation directly by asking for both bits simultaneously; when $\gamma = 1$ the verifier checks anticommutation by playing the Magic Square game, using Theorem~\ref{thm:ms-from-ac}.

\begin{definition}[Admissible parameters]\label{def:admissible}
	\lean{MIPStarRE.QPBT.AdmissibleParams, MIPStarRE.QPBT.AdmissibleParams.one_le_m}
	\leanok
	\uses{def:admissible-size}
	% The source introduces q, m, d in \N with \N the positive integers
	% (references/qpbt-paper/04_preliminaries.tex:6 and
	% references/qpbt-paper/08_classical_and_quantum_low_degree_tests.tex:903-905),
	% so d >= 1 is implicit in its notion of admissibility (lines 958-961 there);
	% we state it explicitly. It is load-bearing: at d = 0 the soundness function
	% delta_qld of thm:pauli vanishes identically while every strategy satisfies
	% the success hypothesis at eps = 1, and the low-degree check of
	% def:pauli-win-predicate invokes the classical game at ldparams = (q,m,d,1),
	% whose domain (Definition def:ld-game) requires d >= 1. The hypothesis
	% propagates to lem:pauli-completeness, thm:pauli, and cor:pauli-binary
	% through their quantification over admissible tuples.
	The tuple $\mathsf{qldparams} = (q,m,d)$ is \emph{admissible} if $d \ge 1$, $q$ is an admissible field size (Definition~\ref{def:admissible-size}), and $m$ divides $q$. Since an admissible field size satisfies $q \ge 2$, the divisibility forces $m \ge 1$, so an admissible tuple has all entries positive.
\end{definition}

\begin{definition}[Question distribution of the Pauli basis test]\label{def:pauli-question-distribution}
	\lean{MIPStarRE.QPBT.PauliType, MIPStarRE.QPBT.pauliCLLevel, MIPStarRE.QPBT.pauliCL}
	\lean{MIPStarRE.QPBT.pauliEdges}
	\lean{MIPStarRE.QPBT.pauliQuestionDistribution}
	\leanok
	\uses{def:admissible, def:ld-question-distribution, lem:ld-point-level, lem:ld-aline-level, lem:ld-dline-level, def:ms-game, def:cl-func, def:cl-dist, def:graph-distribution}
	% Admissibility of the full tuple is required here, not only the divisibility m | q:
	% the source parametrizes the Pauli basis test by the tuple qldparams introduced
	% together with Definition def:admissible (references/qpbt-paper/
	% 08_classical_and_quantum_low_degree_tests.tex:958-1075), and both the fixed binary
	% identification of F_q used by the map chi and the trace tr in eq:gamma-value
	% require q to be a power of 2.
	The game $\mathfrak{G}^{\mathrm{Pauli}}$ (also $\mathfrak{G}^{\mathrm{Pauli}}_{\mathsf{qldparams}}$) is parametrized by an admissible tuple $\mathsf{qldparams} = (q,m,d)$ (Definition~\ref{def:admissible}): admissibility makes $q$ a power of $2$, so that the field $\F_q$, its fixed identification with binary strings (Chapter~\ref{chap:qpbt-algebra}), and the trace $\tr\colon \F_q \to \F_2$ used by the win predicate below are defined, and it provides the divisibility $m \mid q$ required by the index map $\chi$ of Definition~\ref{def:ld-question-distribution}. The question type set is
	\begin{equation}
		\label{eq:pauli-type}
		\mathcal T^{\mathrm{pauli}} = \bigl( \{\mathsf{Point}, \mathsf{ALine}, \mathsf{DLine}, \mathsf{Pauli}, \mathsf{Pair}\} \times \{X,Z\} \bigr) \cup \mathcal T^{\mathrm{ms}} \cup \{\mathsf{Pair}\}\;,
	\end{equation}
	where $\mathcal T^{\mathrm{ms}}$ is the type set of the Magic Square game (Definition~\ref{def:ms-game}).

	The ambient space is $V^{\mathrm{pauli}} = (\F_q^m)^2 \times \F_q \times \F_q^m \times (\F_q)^2$, decomposed into the register subspaces $V_X, V_Z$ (each $m$-dimensional over $\F_q$), $V_I$ ($1$-dimensional), $V_V$ ($m$-dimensional), and $V_{R_X}, V_{R_Z}$ ($1$-dimensional each); elements of $V^{\mathrm{pauli}}$ are written $(u_X, u_Z, s, v, r_X, r_Z)$. For each type $t \in \mathcal T^{\mathrm{pauli}}$ define a CL function $L_t\colon V^{\mathrm{pauli}} \to V^{\mathrm{pauli}}$ as follows, where $L_{\mathsf{Point}}, L_{\mathsf{ALine}}, L_{\mathsf{DLine}}$ are the CL functions of Definition~\ref{def:ld-question-distribution}:
	\begin{enumerate}
		\item for $W \in \{X,Z\}$, $L_{(\mathsf{Point},W)}(u_X,u_Z,s,v,r_X,r_Z) = L_{\mathsf{Point}}(u_W,s,v)$, a $1$-level CL function, where the range $V_W \oplus V_I \oplus V_V$ of $L_{\mathsf{Point}}$ is embedded in $V^{\mathrm{pauli}}$ in the natural way (the same convention applies in items~2 and~3);
		\item for $W \in \{X,Z\}$, $L_{(\mathsf{ALine},W)}(u_X,u_Z,s,v,r_X,r_Z) = L_{\mathsf{ALine}}(u_W,s,v)$, a $2$-level CL function;
		\item for $W \in \{X,Z\}$, $L_{(\mathsf{DLine},W)}(u_X,u_Z,s,v,r_X,r_Z) = L_{\mathsf{DLine}}(u_W,s,v)$, a $3$-level CL function;
		\item for $t \in \mathcal T^{\mathrm{ms}} \cup \{\mathsf{Pair}, (\mathsf{Pair},X), (\mathsf{Pair},Z)\}$, $L_t(u_X,u_Z,s,v,r_X,r_Z) = (u_X,u_Z,0,0,r_X,r_Z)$, the $1$-level CL function projecting onto $V_X \oplus V_Z \oplus V_{R_X} \oplus V_{R_Z}$;
		\item for $W \in \{X,Z\}$, $L_{(\mathsf{Pauli},W)} = 0$, the identically zero ($0$-level) CL function.
	\end{enumerate}
	The maps $L_{(\mathsf{Point},W)}$, $L_{(\mathsf{ALine},W)}$, $L_{(\mathsf{DLine},W)}$ act as the zero function on $V_{\overline W} \oplus V_{R_X} \oplus V_{R_Z}$, where $\overline W$ denotes the basis other than $W$.

	The \emph{type graph} $G^{\mathrm{pauli}}$ has vertex set $\mathcal T^{\mathrm{pauli}}$ and the following edges: all $18$ edges of the Magic Square type graph $G^{\mathrm{ms}}$ (Definition~\ref{def:ms-game}); the edges
	\begin{align*}
		&(\mathsf{ALine},X) \,\text{---}\, (\mathsf{Point},X)\;, &&(\mathsf{DLine},X) \,\text{---}\, (\mathsf{Point},X)\;, &&(\mathsf{Point},X) \,\text{---}\, (\mathsf{Pauli},X)\;,\\
		&(\mathsf{ALine},Z) \,\text{---}\, (\mathsf{Point},Z)\;, &&(\mathsf{DLine},Z) \,\text{---}\, (\mathsf{Point},Z)\;, &&(\mathsf{Point},Z) \,\text{---}\, (\mathsf{Pauli},Z)\;,\\
		&(\mathsf{Point},X) \,\text{---}\, \mathsf{Variable}_1\;, &&(\mathsf{Point},Z) \,\text{---}\, \mathsf{Variable}_5\;, &&\\
		&(\mathsf{Point},X) \,\text{---}\, (\mathsf{Pair},X)\;, &&(\mathsf{Pair},X) \,\text{---}\, \mathsf{Pair}\;, &&\\
		&(\mathsf{Point},Z) \,\text{---}\, (\mathsf{Pair},Z)\;, &&(\mathsf{Pair},Z) \,\text{---}\, \mathsf{Pair}\;; &&
	\end{align*}
	and, in addition, a self-loop at every vertex.

	% The type pair is drawn from the graph distribution of Definition
	% def:graph-distribution, matching the source's formal definition
	% (references/qpbt-paper/07_types.tex:84-93 together with
	% references/qpbt-paper/08_classical_and_quantum_low_degree_tests.tex:1115-1120).
	% Sampling an undirected edge uniformly and then ordering its
	% endpoints uniformly would be a different distribution here: it gives each self-loop
	% twice the relative mass of an oriented non-loop pair, and G^pauli has a self-loop
	% at every vertex.
	The question distribution of $\mathfrak{G}^{\mathrm{Pauli}}$ is the typed CL distribution on $\mathcal T^{\mathrm{pauli}} \times V^{\mathrm{pauli}} \times \mathcal T^{\mathrm{pauli}} \times V^{\mathrm{pauli}}$ obtained by sampling the type pair $(t_{\mathrm A}, t_{\mathrm B})$ from the graph distribution $\mu_{G^{\mathrm{pauli}}}$ of Definition~\ref{def:graph-distribution}, that is, uniformly at random among all ordered pairs of types whose underlying unordered pair $\{t_{\mathrm A}, t_{\mathrm B}\}$ is an edge of $G^{\mathrm{pauli}}$, self-loops included; sampling $z \in V^{\mathrm{pauli}}$ uniformly at random; and setting $x_w = L_{t_w}(z)$ for $w \in \{\mathrm A,\mathrm B\}$. Explicitly, if $z = (u_X,u_Z,s,v,r_X,r_Z)$, then a player whose type is $(\mathsf{Point},W)$ receives $u_W$; a player whose type is $(\mathsf{ALine},W)$ receives $(u_0, s)$ with $u_0 = L^{\mathrm{Ln}}_{e_i}(u_W)$ and $i = \chi(s)$; a player whose type is $(\mathsf{DLine},W)$ receives $(u_0, s, v')$ with $i = \chi(s)$, $v' = \pi_{i-1}(v)$, and $u_0 = L^{\mathrm{Ln}}_{v'}(u_W)$; a player whose type lies in $\mathcal T^{\mathrm{ms}} \cup \{\mathsf{Pair}, (\mathsf{Pair},X), (\mathsf{Pair},Z)\}$ receives $(u_X, u_Z, r_X, r_Z)$; and a player whose type is $(\mathsf{Pauli},W)$ receives $0$.
\end{definition}

\begin{lemma}[Conditionally linear Pauli question maps]\label{lem:pauli-question-cl-components}
	\lean{MIPStarRE.QPBT.isCondLinear_pauliCL}
	\leanok
	\uses{def:pauli-question-distribution, def:cl-func, lem:ld-point-level, lem:ld-aline-level, lem:ld-dline-level}
	For every Pauli question type $t \in \mathcal T^{\mathrm{pauli}}$, the map $L_t$ of Definition~\ref{def:pauli-question-distribution} is conditionally linear. Its levels are $1$ for the point maps, $2$ for the axis-parallel line maps, $3$ for the diagonal line maps, $1$ for the shared projection maps, and $0$ for the Pauli maps.
\end{lemma}
\begin{proof}
	\leanok
	\uses{lem:ld-point-level, lem:ld-aline-level, lem:ld-dline-level, lem:cl-reindex, lem:cl-one-level-linear}
	Transport the point, axis-parallel line, and diagonal line representations to the selected basis register. The shared projection onto the point and auxiliary registers is linear, while the Pauli maps are zero. These representations have levels $1$, $2$, $3$, $1$, and $0$, respectively.
\end{proof}

\begin{lemma}[Common level of the Pauli question maps]\label{lem:pauli-question-cl-family}
	\lean{MIPStarRE.QPBT.isTypedCondLinearFamily_pauliCL}
	\leanok
	\uses{def:pauli-question-distribution, def:typed-cl-functions}
	For every Pauli question type $t \in \mathcal T^{\mathrm{pauli}}$, the map $L_t$ of Definition~\ref{def:pauli-question-distribution} is a $3$-level CL function. Thus $L = \{L_t\}_{t \in \mathcal T^{\mathrm{pauli}}}$ is a finite typed family of $3$-level CL functions.
\end{lemma}
\begin{proof}
	\leanok
	\uses{lem:pauli-question-cl-components, lem:cl-level-mono}
	The component maps have levels at most $3$ by Lemma~\ref{lem:pauli-question-cl-components}. Lemma~\ref{lem:cl-level-mono} raises each representation to the common level $3$.
\end{proof}

\begin{lemma}[Pauli question-distribution identity]\label{lem:pauli-question-typed-equality}
	\lean{MIPStarRE.QPBT.pauliQuestionDistribution_eq_typedCL}
	\leanok
	\uses{def:pauli-question-distribution, def:typed-cl-distributions, def:graph-distribution}
	For the Pauli type graph $G^{\mathrm{pauli}}$ and the family $L$ of Definition~\ref{def:pauli-question-distribution}, the Pauli question distribution equals $\mu^{G^{\mathrm{pauli}}}_{L,L}$.
\end{lemma}
\begin{proof}
	\leanok
	\uses{def:typed-cl-distributions, def:graph-distribution, def:cl-dist, lem:distribution-map-comp, lem:uniform-bind-on-finset-map}
	Both distributions first use the uniform law on ordered pairs whose unordered pair is an edge of $G^{\mathrm{pauli}}$, and then push forward a common uniform seed in $V^{\mathrm{pauli}}$ by the same pair of maps. Writing the ordered-edge law as the push-forward of the uniform law on the ordered-edge set, Lemma~\ref{lem:uniform-bind-on-finset-map} presents each side as the push-forward of the uniform law on the product of the ordered-edge set with $V^{\mathrm{pauli}}$ along the same map.
\end{proof}

\begin{lemma}[The Pauli question distribution is a typed CL distribution]\label{lem:pauli-question-typed-cl}
	\lean{MIPStarRE.QPBT.pauliQuestionDistribution_isTypedCL}
	\leanok
	\uses{def:pauli-question-distribution, def:typed-cl-distributions}
	Fix an admissible tuple $\mathsf{qldparams} = (q,m,d)$ (Definition~\ref{def:admissible}) and let $L = \{L_t\}_{t \in \mathcal T^{\mathrm{pauli}}}$ be the family of Definition~\ref{def:pauli-question-distribution}. The family $L$ is a finite $\mathcal T^{\mathrm{pauli}}$-typed family of $3$-level conditionally linear functions, and the question distribution of $\mathfrak{G}^{\mathrm{Pauli}}_{\mathsf{qldparams}}$ equals the $(\mathcal T^{\mathrm{pauli}}, G^{\mathrm{pauli}})$-typed conditionally linear distribution $\mu^{G^{\mathrm{pauli}}}_{L,\,L}$ of Definition~\ref{def:typed-cl-distributions}.
\end{lemma}
\begin{proof}
	\leanok
	\uses{lem:pauli-question-cl-family, lem:pauli-question-typed-equality}
	Combine the common-level family assertion of Lemma~\ref{lem:pauli-question-cl-family} with the distribution identity of Lemma~\ref{lem:pauli-question-typed-equality}.
\end{proof}

\begin{remark}[Formalization status of Lemma~\ref{lem:pauli-question-typed-cl}]\label{rem:pauli-question-typed-cl-status}
	Both halves of Lemma~\ref{lem:pauli-question-typed-cl} are proved. The common-level assertion is Lemma~\ref{lem:pauli-question-cl-family}, and the distribution identity is Lemma~\ref{lem:pauli-question-typed-equality}, whose proof identifies the two common-seed laws through Lemma~\ref{lem:uniform-bind-on-finset-map}. The Pauli question distribution therefore carries the typed conditionally linear structure without further hypotheses, exactly as the low-degree question distribution does through Lemma~\ref{lem:ld-question-typed-cl}.
\end{remark}

\begin{definition}[Win predicate of the Pauli basis test]\label{def:pauli-win-predicate}
	\lean{MIPStarRE.QPBT.PauliAnswer, MIPStarRE.QPBT.validPauliAnswer}
	\lean{MIPStarRE.QPBT.pauliWinPredicate, MIPStarRE.QPBT.pauliBasisTest}
	\leanok
	\uses{def:pauli-question-distribution, def:ld-win-predicate, def:ms-game, def:indicator-vector, def:admissible, def:game, def:low-degree-encoding, def:subfield-trace}
	Fix an admissible tuple $\mathsf{qldparams} = (q,m,d)$ (Definition~\ref{def:admissible}) and let $M = 2^m$. Answers in the game $\mathfrak{G}^{\mathrm{Pauli}}_{\mathsf{qldparams}}$ take the following forms, depending on the question type; answers not of the prescribed form are rejected.
	% The paper's format table writes the line contents loosely as (u_W, s) and (u_W, s, v);
	% by the sampling procedure the actual contents are the canonical pairs (u_0, s) and
	% triples (u_0, s, v') with v' = pi_{i-1}(v). We state the canonical form.
	\begin{itemize}
		\item $(\mathsf{Point},W)$: content $y \in \F_q^m$; answer an element of $\F_q$.
		\item $(\mathsf{ALine},W)$: content $(u_0, s) \in \F_q^m \times \F_q$; answer a univariate polynomial $f\colon \F_q \to \F_q$ of degree at most $d$.
		\item $(\mathsf{DLine},W)$: content $(u_0, s, v') \in \F_q^m \times \F_q \times \F_q^m$; answer a univariate polynomial $f\colon \F_q \to \F_q$ of degree at most $md$.
		\item $\mathsf{Pair}$: content $\omega = (u_X, u_Z, r_X, r_Z) \in (\F_q^m)^2 \times \F_q^2$; answer $(\beta_X, \beta_Z) \in \F_2^2$.
		\item $(\mathsf{Pair},W)$: content $\omega$ as above; answer an element of $\F_2$.
		\item $\mathsf{Constraint}_i$: content $\omega$ as above; answer $(\alpha_{v_1}, \alpha_{v_2}, \alpha_{v_3}) \in \F_2^3$, indexed by the variable indices $(v_1,v_2,v_3)$ of $\mathsf{Constraint}_i$.
		\item $\mathsf{Variable}_j$: content $\omega$ as above; answer an element of $\F_2$.
		\item $(\mathsf{Pauli},W)$: content $0$; answer an element of $\F_q^M$.
	\end{itemize}
	The win predicate $D^{\mathrm{pauli}} = D^{\mathrm{pauli}}_{\mathsf{qldparams}}$ is the function of the tuple $(t_{\mathrm A}, x_{\mathrm A}, t_{\mathrm B}, x_{\mathrm B}, a_{\mathrm A}, a_{\mathrm B})$ determined by the following clauses, evaluated for $w \in \{\mathrm A, \mathrm B\}$ and $W \in \{X,Z\}$; in all cases where no action is indicated, the predicate accepts. In the last four clauses, the content of the player whose type lies in $\mathcal T^{\mathrm{ms}} \cup \{\mathsf{Pair}, (\mathsf{Pair},X), (\mathsf{Pair},Z)\}$ is a tuple $(u_X, u_Z, r_X, r_Z)$, from which the predicate first computes
	\begin{equation}
		\label{eq:gamma-value}
		\gamma = \tr\bigl( (\mathrm{ind}_m(u_X) \cdot r_X) \cdot (\mathrm{ind}_m(u_Z) \cdot r_Z) \bigr) \in \F_2\;,
	\end{equation}
	where $\mathrm{ind}_m(\cdot) \in \F_q^M$ is the indicator vector of the low-degree encoding (Chapter~\ref{chap:qpbt-algebra}), the inner dot denotes the dot product on $\F_q^M$, and $\tr\colon \F_q \to \F_2$ is the finite-field trace.
	\begin{enumerate}
		\item (\textbf{Consistency check}) If $t_{\mathrm A} = t_{\mathrm B}$, accept if and only if $a_{\mathrm A} = a_{\mathrm B}$.
		\item (\textbf{Low-degree check}) Let $\mathsf{ldparams} = (q,m,d,1)$; admissibility of $\mathsf{qldparams}$ gives $d \ge 1$ and $m \mid q$, so this is a parameter tuple as in Definition~\ref{def:ld-game}.
		\begin{enumerate}
			\item If $t_w = (\mathsf{Point},W)$ and $t_{\overline w} = (\mathsf{ALine},W)$, accept if $D^{\mathrm{ld}}_{\mathsf{ldparams}}$ accepts $(\mathsf{Point}, x_w, \mathsf{ALine}, x_{\overline w}, a_w, a_{\overline w})$.
			\item If $t_w = (\mathsf{Point},W)$ and $t_{\overline w} = (\mathsf{DLine},W)$, accept if $D^{\mathrm{ld}}_{\mathsf{ldparams}}$ accepts $(\mathsf{Point}, x_w, \mathsf{DLine}, x_{\overline w}, a_w, a_{\overline w})$.
		\end{enumerate}
		\item (\textbf{Consistency check}) If $t_w = (\mathsf{Point},W)$ and $t_{\overline w} = (\mathsf{Pauli},W)$, accept if $g_{a_{\overline w}}(x_w) = a_w$, where $g_{a_{\overline w}}$ is the low-degree encoding of $a_{\overline w} \in \F_q^M$ (Chapter~\ref{chap:qpbt-algebra}).
		\item (\textbf{Commutation check}) If $t_w = (\mathsf{Pair},W)$ and $t_{\overline w} = \mathsf{Pair}$, accept if $a_w = \beta_W$ or $\gamma \neq 0$.
		\item (\textbf{Consistency check}) If $t_w = (\mathsf{Point},W)$ and $t_{\overline w} = (\mathsf{Pair},W)$, accept if $\tr(a_w r_W) = a_{\overline w}$ or $\gamma \neq 0$.
		\item (\textbf{Magic Square check}) If $t_w = \mathsf{Constraint}_i$ and $t_{\overline w} = \mathsf{Variable}_j$, accept if $\gamma = 0$, or if $a_w$ satisfies the equation of $\mathsf{Constraint}_i$ and $\alpha_j = a_{\overline w}$.
		\item (\textbf{Consistency check}) If $t_w = (\mathsf{Point},W)$ and $t_{\overline w} = \mathsf{Variable}_j$, accept if $\gamma = 0$, or if $j = 1$, $W = X$, and $\tr(a_w r_X) = a_{\overline w}$, or if $j = 5$, $W = Z$, and $\tr(a_w r_Z) = a_{\overline w}$.
	\end{enumerate}
	The game $\mathfrak{G}^{\mathrm{Pauli}}_{\mathsf{qldparams}}$ is the typed game with type set $\mathcal T^{\mathrm{pauli}}$, the question distribution of Definition~\ref{def:pauli-question-distribution}, and win predicate $D^{\mathrm{pauli}}_{\mathsf{qldparams}}$.
\end{definition}

Note the one-sided gating on $\gamma$: the commutation check and the $\mathsf{Pair}$ consistency check accept automatically when $\gamma \neq 0$, while the Magic Square check and the $\mathsf{Variable}$ consistency check accept automatically when $\gamma = 0$.

\begin{lemma}[Completeness of the Pauli basis test]\label{lem:pauli-completeness}
	\lean{MIPStarRE.QPBT.pauliBasisTestSymm, MIPStarRE.QPBT.exists_spcc_value_one}
	\lean{MIPStarRE.QPBT.honestStrategy, MIPStarRE.QPBT.honestStrategy_isSPCC}
	\leanok
	\uses{def:pauli-question-distribution, def:pauli-win-predicate, def:admissible, def:spcc, def:tensor-product-value}
	For every admissible tuple $\mathsf{qldparams} = (q,m,d)$ (Definition~\ref{def:admissible}), the Pauli basis test $\mathfrak{G}^{\mathrm{Pauli}}_{\mathsf{qldparams}}$ has a value-$1$ SPCC strategy.
\end{lemma}
\begin{proof}
	\leanok
	\uses{thm:ms-from-ac, def:pauli-win-predicate, def:bracket, def:consistent-measurement, def:EPR, def:indicator-vector, lem:pauli-observable-expansion, lem:twisted-commutation, lem:low-degree-encoding-line-restriction}
	The strategy uses the state $\ket{\psi} = \ket{\mathrm{EPR}_q}^{\ot M} \ot \ket{\mathrm{EPR}_2}$ with $M = 2^m$. On a question $((\mathsf{Point},W), y)$ a player measures the generalized Pauli basis POVM $\{\tau^W_h\}_{h \in \F_q^M}$ and reports $g_h(y)$; on $((\mathsf{ALine},W), \ell)$ or $((\mathsf{DLine},W), \ell)$ they report the restriction of $g_h$ to the line $\ell$, expressed as a univariate polynomial through the canonical parameterization of $\ell$; and on $(\mathsf{Pauli},W)$ they report $h$ itself --- in the bracket notation of Definition~\ref{def:bracket}, these measurements are $\tau^W_{[g_\cdot(y) = a]} \ot I$, $\tau^W_{[(g_\cdot)|_\ell = f]} \ot I$, and $\tau^W_a \ot I$. For a content $\omega = (u_X, u_Z, r_X, r_Z)$, consider the two binary measurements $A_b = \tau^X_{[\tr(g_\cdot(u_X)\, r_X) = b]} \ot I$ and $B_b = \tau^Z_{[\tr(g_\cdot(u_Z)\, r_Z) = b]} \ot I$. Since $g_h(u) = h \cdot \mathrm{ind}_m(u)$ (Chapter~\ref{chap:qpbt-algebra}), the relation between the Pauli projections and observables gives $\mathcal O_A = A_0 - A_1 = \tau^X(\mathrm{ind}_m(u_X)\, r_X) \ot I$ and $\mathcal O_B = \tau^Z(\mathrm{ind}_m(u_Z)\, r_Z) \ot I$, and the twisted commutation relation yields
	\[
		\mathcal O_A \mathcal O_B = (-1)^\gamma\, \mathcal O_B \mathcal O_A
	\]
	with $\gamma$ as in Equation~\eqref{eq:gamma-value}: the two measurements commute when $\gamma = 0$ and anticommute when $\gamma = 1$. If $\gamma = 0$, define $E^{\mathsf{Pair},\omega}_{\beta_X,\beta_Z} = A_{\beta_X} B_{\beta_Z}$ and $E^{(\mathsf{Pair},X),\omega}_a = A_a$, $E^{(\mathsf{Pair},Z),\omega}_a = B_a$, and let the $\mathsf{Constraint}$ and $\mathsf{Variable}$ measurements be trivial (identity on the all-zeros outcome). If $\gamma = 1$, let the $\mathsf{Constraint}$ and $\mathsf{Variable}$ measurements be those produced by Theorem~\ref{thm:ms-from-ac} from the anticommuting consistent pair $(A, B)$, acting on the auxiliary $\ket{\mathrm{EPR}_2}$ factor; keep $E^{(\mathsf{Pair},X),\omega}_a = A_a$ and $E^{(\mathsf{Pair},Z),\omega}_a = B_a$ for the $(\mathsf{Pair},W)$ questions, and let the joint $\mathsf{Pair}$ measurement be trivial. The strategy is symmetric and projective by construction; every measurement is either a post-processed Pauli basis measurement, a measurement supplied by Theorem~\ref{thm:ms-from-ac}, or (in the $\gamma = 0$ case) a product of two commuting consistent measurements, so the strategy is consistent, and commutativity along every edge of $G^{\mathrm{pauli}}$ is checked case by case --- for instance, $(\mathsf{Point},X)$ and $\mathsf{Variable}_1$ commute because both are $X$-basis measurements. Consistency and projectivity give the consistency checks with certainty, and the low-degree checks pass because the answers agree with the strategy determined by the polynomial $g_h$ in the classical game. For the remaining checks: when $\gamma = 0$ the commutation check passes by construction, and when $\gamma = 1$ the Magic Square check passes by Theorem~\ref{thm:ms-from-ac}; the consistency checks 5 and 7 pass with certainty because $\tau^W_{[\tr(g_\cdot(y)\, r_W) = \cdot]} \ot I$ coincides with $E^{(\mathsf{Pair},W),\omega}$, respectively with the embedded $\mathsf{Variable}_1$ and $\mathsf{Variable}_5$ measurements of Theorem~\ref{thm:ms-from-ac} (item~3 there). Hence the strategy has value $1$.
\end{proof}

The proof above differs from the source in one inessential respect. When
$\gamma\neq 0$ the source makes the $(\mathsf{Pair},W)$ measurements trivial
along with the joint $\mathsf{Pair}$ measurement, whereas the strategy stated
here --- and the formal witness --- keeps $A$ and $B$ on those questions. Both
choices are admissible: the commutation and $\mathsf{Pair}$-consistency checks
accept automatically when $\gamma\neq 0$, and $A$ and $B$ remain consistent
coarse-grainings of a single Pauli basis. Only the existential witness changes,
so the SPCC and value-one conclusions of the lemma are unaffected.

We now state the main soundness property of the Pauli basis test; the statement is an adaptation of the self-testing theorem of~\cite{Natarajan2018NEEXP} (Theorem~6.4 there). Unlike Theorem~\ref{lem:ld-soundness}, no projectivity assumption is made on the strategy.

\begin{theorem}[Soundness of the Pauli basis test]\label{thm:pauli}
	\lean{MIPStarRE.QPBT.pauli_soundness}
	\leanok
	\uses{def:pauli-question-distribution, def:pauli-win-predicate, def:admissible, def:tensor-product-value, def:EPR, lem:pauli-observable-expansion, def:povm-distance}
	There exists a function
	\[
		\delta_{\mathrm{qld}}(\eps, m, d, q) = a\,(md)^a \bigl( \eps^b + q^{-b} + 2^{-bmd} \bigr)
	\]
	for universal constants $a \ge 1$ and $0 < b < 1$ such that the following holds. For all admissible parameter tuples $\mathsf{qldparams} = (q,m,d)$ and for all strategies $\mathcal S = (\ket\psi, A, B)$ for the game $\mathfrak{G}^{\mathrm{Pauli}}_{\mathsf{qldparams}}$ that succeed with probability at least $1 - \eps$, there exist local isometries $\phi_{\mathrm A}\colon \hilb_{\mathrm A} \to \hilb_{\mathrm A'} \ot \hilb_{\mathrm A''}$ and $\phi_{\mathrm B}\colon \hilb_{\mathrm B} \to \hilb_{\mathrm B'} \ot \hilb_{\mathrm B''}$, where $\ket\psi \in \hilb_{\mathrm A} \ot \hilb_{\mathrm B}$ and $\hilb_{\mathrm A''}, \hilb_{\mathrm B''} \cong (\C^q)^{\ot M}$ with $M = 2^m$, and a state $\ket{\mathrm{aux}} \in \hilb_{\mathrm A'} \ot \hilb_{\mathrm B'}$, such that:
	\begin{enumerate}
		\item $\bigl\| \phi_{\mathrm A} \ot \phi_{\mathrm B} \ket\psi - \ket{\mathrm{aux}} \ot \ket{\mathrm{EPR}_q}^{\ot M} \bigr\| \le \delta_{\mathrm{qld}}(\eps, m, d, q)$;
		% The paper here has the typo "phi_alice \ot phi_B"; both maps are the isometries above.
		\item letting $\tilde A^x_a = \phi_{\mathrm A} A^x_a \phi_{\mathrm A}^\dagger$ and $\tilde B^y_b = \phi_{\mathrm B} B^y_b \phi_{\mathrm B}^\dagger$, for $W \in \{X, Z\}$,
		\begin{align*}
			\tilde A^{(\mathsf{Pauli},W)}_u \ot I_{\mathrm{B'B''}} &\approx_{\delta_{\mathrm{qld}}} (\tau^W_u)_{\mathrm A''} \ot I_{\mathrm{A'B'B''}}\;,\\
			I_{\mathrm{A'A''}} \ot \tilde B^{(\mathsf{Pauli},W)}_u &\approx_{\delta_{\mathrm{qld}}} I_{\mathrm{A'A''B'}} \ot (\tau^W_u)_{\mathrm B''}\;,
		\end{align*}
		where the $\approx_{\delta_{\mathrm{qld}}}$ statements hold with respect to the state $\ket{\mathrm{aux}}_{\mathrm{A'B'}} \ot \ket{\mathrm{EPR}_q}^{\ot M}_{\mathrm{A''B''}}$ and the answer summation is over $u \in \F_q^M$.
	\end{enumerate}
\end{theorem}
\begin{proof}
	\leanok
	\uses{thm:pauli-arbitrary-strategy-raw-isometry-support}
	Apply Theorem~\ref{thm:pauli-arbitrary-strategy-raw-isometry-support}. It
	supplies the universal constants, local isometries, auxiliary state, and all
	three displayed estimates for every $\eps\geq0$, with the operator estimates
	formed from the prescribed-answer effects themselves.
\end{proof}

The argument adapts the analysis of the quantum
low-degree test of Natarajan and Vidick to low individual degree. It is
developed in Chapters~\ref{chap:qpbt-observables}
and~\ref{chap:qpbt-combining} and concluded in
Chapter~\ref{chap:qpbt-extraction}. The established argument uses directly
indexed low-degree soundness and transfers the resulting completed Pauli
measurements to the prescribed-answer effects. The seed-indexed derivation
printed in the source remains subject to the divisibility and identification
obligations documented in \cite{gap:qpbt_ld-dimension-divisibility} and is
recorded separately in Remark~\ref{rem:pauli-printed-seed-argument}.
For the Magic Square step, the soundness argument uses the original-state
anticommutation bounds of
Lemma~\ref{lem:qld-ms-anticommutator-original-state}, obtained from winning
probability alone. The additional variable-agreement hypotheses of
Theorem~\ref{thm:ms-rigidity} are not inputs to that argument.

The strategies appearing in Lemma~\ref{lem:pauli-completeness} and in the conclusion of Theorem~\ref{thm:pauli} are phrased in terms of qudit Pauli measurements and maximally entangled states of $q$-dimensional qudits. For the intended application it is more convenient to work with qubits: since the field sizes $q$ used here are powers of $2$, the conclusions can be restated in terms of qubit EPR pairs and qubit Pauli measurements.

\begin{corollary}[Qubit form of the soundness of the Pauli basis test]\label{cor:pauli-binary}
	\lean{MIPStarRE.QPBT.pauli_soundness_qubit}
	\leanok
	\uses{def:pauli-question-distribution, def:pauli-win-predicate, def:admissible, def:tensor-product-value, def:EPR, def:povm-distance, def:binary-representation}
	There exists a function
	\[
		\delta_{\mathrm{qld}}(\eps, m, d, q) = a\,(md)^a \bigl( \eps^b + q^{-b} + 2^{-bmd} \bigr)
	\]
	% The paper's displayed function here writes the argument order (eps, q, m, d); the
	% formula is the same, and its own item 1 reverts to (eps, m, d, q). See the remark below.
	for universal constants $a \ge 1$ and $0 < b < 1$ such that the following holds. For all admissible parameter tuples $\mathsf{qldparams} = (q,m,d)$ and for all strategies $\mathcal S = (\ket\psi, A, B)$ for the game $\mathfrak{G}^{\mathrm{Pauli}}_{\mathsf{qldparams}}$ that succeed with probability at least $1 - \eps$, there exist local isometries $\phi_{\mathrm A}\colon \hilb_{\mathrm A} \to \hilb_{\mathrm A'} \ot \hilb_{\mathrm A''}$ and $\phi_{\mathrm B}\colon \hilb_{\mathrm B} \to \hilb_{\mathrm B'} \ot \hilb_{\mathrm B''}$, where $\ket\psi \in \hilb_{\mathrm A} \ot \hilb_{\mathrm B}$ and $\hilb_{\mathrm A''}, \hilb_{\mathrm B''} \cong (\C^2)^{\ot M \log q}$ with $M = 2^m$, and a state $\ket{\mathrm{aux}} \in \hilb_{\mathrm A'} \ot \hilb_{\mathrm B'}$, such that:
	\begin{enumerate}
		\item $\bigl\| \phi_{\mathrm A} \ot \phi_{\mathrm B} \ket\psi - \ket{\mathrm{aux}} \ot \ket{\mathrm{EPR}_2}^{\ot M \log q} \bigr\| \le \delta_{\mathrm{qld}}(\eps, m, d, q)$;
		\item letting $\tilde A^x_a = \phi_{\mathrm A} A^x_a \phi_{\mathrm A}^\dagger$ and $\tilde B^y_b = \phi_{\mathrm B} B^y_b \phi_{\mathrm B}^\dagger$, for $W \in \{X, Z\}$,
		\begin{align*}
			\tilde A^{(\mathsf{Pauli},W)}_u \ot I_{\mathrm{B'B''}} &\approx_{\delta_{\mathrm{qld}}} (\sigma^W_u)_{\mathrm A''} \ot I_{\mathrm{A'B'B''}}\;,\\
			I_{\mathrm{A'A''}} \ot \tilde B^{(\mathsf{Pauli},W)}_u &\approx_{\delta_{\mathrm{qld}}} I_{\mathrm{A'A''B'}} \ot (\sigma^W_u)_{\mathrm B''}\;,
		\end{align*}
		where the $\approx_{\delta_{\mathrm{qld}}}$ statements hold with respect to the state $\ket{\mathrm{aux}}_{\mathrm{A'B'}} \ot \ket{\mathrm{EPR}_2}^{\ot M \log q}_{\mathrm{A''B''}}$ and the answer summation is over $u \in \F_q^M$. Here $\sigma^W_u$ denotes the tensor product of single-qubit Pauli measurements
		\[
			\bigotimes_{i=1}^{M} \bigotimes_{j=1}^{\log q} \sigma^W_{u_{ij}}\;,
		\]
		where $\kappa(u_i) = (u_{ij})_{1 \le j \le \log q}$ under the bijection $\kappa\colon \F_q \to \F_2^{\log q}$ of Chapter~\ref{chap:qpbt-algebra}.
	\end{enumerate}
\end{corollary}
\begin{proof}
	\leanok
	\uses{thm:pauli, lem:pauli-binary}
	By Lemma~\ref{lem:pauli-binary}, since $q$ is a power of $2$ there are local isometries that map $\ket{\mathrm{EPR}_q}^{\ot M}$ to $\ket{\mathrm{EPR}_2}^{\ot M \log q}$ and map the generalized Pauli measurements $\tau^W_u$ to the qubit tensor products $\bigotimes_{i=1}^M \bigotimes_{j=1}^{\log q} \sigma^W_{u_{ij}}$, where $\kappa(u_i) = (u_{i1},\ldots,u_{i \log q})$. Composing these isometries with those produced by Theorem~\ref{thm:pauli} gives the statement.
\end{proof}

\begin{remark}[Notational drift in the soundness functions]\label{rem:delta-qld-argument-order}
	The paper writes the soundness function of the Pauli basis test with two different argument orders: $\delta_{\mathrm{qld}}(\eps, m, d, q)$ in Theorem~\ref{thm:pauli} and Lemma~\ref{lem:delta-bound}, and $\delta_{\mathrm{qld}}(\eps, q, m, d)$ in the displayed function of Corollary~\ref{cor:pauli-binary}, whose own item~1 reverts to $\delta_{\mathrm{qld}}(\eps, m, d, q)$; the formula is the same throughout, and here the order $(\eps, m, d, q)$ is used everywhere. Note also two deliberate asymmetries with the classical soundness function of Theorem~\ref{lem:ld-soundness}: there the polynomial prefactor is $a(dmk)^a$ and the exponent satisfies $0 < b \le 1$, while in Theorems~\ref{thm:pauli} and Corollary~\ref{cor:pauli-binary} the prefactor is $a(md)^a$ and the exponent satisfies the strict inequality $0 < b < 1$.
\end{remark}

\section{Canonical parameters of the Pauli basis test}
% Paper label sec:qld-complexity kept for cross-reference fidelity (the complexity
% statements of that subsection are outside the scope of this development).
\label{sec:qld-complexity}

We specify a canonical setting of the parameter tuple $\mathsf{qldparams}$ as a function of an integer $R$, and bound the resulting soundness function. Here $\log$ denotes the base-two logarithm, and $\ln$ denotes the natural logarithm.

\begin{definition}[Canonical parameters of the Pauli basis test]\label{def:introparams}
	\lean{MIPStarRE.QPBT.introParamsC, MIPStarRE.QPBT.introParamsTuple, MIPStarRE.QPBT.AdmissibleParams.ofTuple, MIPStarRE.QPBT.introParams}
	\leanok
	\uses{thm:pauli}
	Let $a, b$ be the universal constants of Theorem~\ref{thm:pauli}, and let $c$ be the smallest \emph{even} integer that is at least $(b+a)/b$. For every integer $R \ge 4$, define the tuple $\mathsf{introparams}(R) = (q, m, d)$ by
	\begin{itemize}
		\item $q = 2^k$ for $k = c \lceil \log \log R \rceil + 1$,
		\item $m = 2^j$, where $j$ is the unique integer with $2^j \le c \lceil \log R \rceil + 1 < 2^{j+1}$,
		\item $d = 1$.
	\end{itemize}
	The letter $k$ here denotes the exponent of the field size and is unrelated to the simultaneity parameter $k$ of the classical game. With this choice the Pauli basis test certifies the presence of an $M = 2^m$-qudit maximally entangled state of local dimension $q$, equivalently an $(M \log q)$-qubit EPR state, and the choice $d = 1$ runs the low individual degree test multilinearly.
\end{definition}

\begin{lemma}[Soundness bound for the canonical parameters]\label{lem:delta-bound}
	\lean{MIPStarRE.QPBT.introParamsTuple_isAdmissible, MIPStarRE.QPBT.le_two_pow_introParams_m, MIPStarRE.QPBT.exists_deltaQld_introParams_bound}
	\leanok
	\uses{def:introparams, def:admissible, thm:pauli}
	For all integers $R \ge 4$, the parameter tuple $\mathsf{introparams}(R) = (q,m,d)$ is admissible (Definition~\ref{def:admissible}) and satisfies $2^m \ge R$. Furthermore, there exist universal constants $a' \ge 1$ and $0 < b' \le 1$ such that, for every $\eps \ge 0$, the function $\delta_{\mathrm{qld}}(\eps, m, d, q)$ of Theorem~\ref{thm:pauli} is at most
	\[
		a' \bigl( (\log R)^{a'} \cdot \eps^{b'} + (\log R)^{-b'} \bigr)\;.
	\]
\end{lemma}
\begin{proof}
	\leanok
	\uses{def:introparams, def:admissible, def:admissible-size, thm:pauli, lem:delta-bound-envelope-support}
	Since $c$ is even, the exponent $k = c \lceil \log\log R \rceil + 1$ is odd, so $q = 2^k$ is an admissible field size; $m$ divides $q$ because $2^j \le c \lceil \log R \rceil + 1 \le 2 \log^c R \le 2^k$, so $j \le k$ and both are powers of $2$; and $d = 1 \ge 1$. Hence the tuple is admissible. Since $m = 2^j \ge (c/2) \log R$, we get $2^m \ge R^{c/2} \ge R$, using $c \ge 2$.

	For the soundness bound, write $\delta_{\mathrm{qld}}(\eps,m,d,q) = a(md)^a \eps^b + a(md)^a q^{-b} + a(md)^a 2^{-bmd}$ and use $d = 1$, $\log R \le m \le 2c \log R$, and $q = 2^k \ge \log^c R$. Set
	\[
		C = \frac{\lceil a + b \rceil !}{(b \ln 2)^{\lceil a + b \rceil}}
		\qquad \text{and} \qquad
		a' = a\bigl( (2c)^a + C \bigr)\;,
	\]
	both depending only on the universal constants $a$ and $b$; since $a \ge 1$, $c \ge 2$ and $C \ge 0$, one has $a' \ge a \ge 1$. The first term is at most $a(2c)^a (\log R)^a \eps^b \le a' (\log R)^{a'} \eps^b$, using $\log R \ge 2 \ge 1$ and $a \le a'$. The second term is at most $a(2c)^a (\log R)^{a - cb} \le a(2c)^a (\log R)^{-b}$, since $cb - a \ge b$ by the choice of $c$. For the third term, split $m^a = m^{a+b} \cdot m^{-b}$ and bound the first factor against the exponential: for all real $s$, all $t > 0$ and all $x \ge 1$, the Taylor estimate $y^n / n! \le e^y$ of the exponential, taken at $y = (t \ln 2) x$ and $n = \lceil s \rceil_{+} = \max(\lceil s \rceil, 0)$, gives
	\[
		x^s\, 2^{-tx} \le \frac{\lceil s \rceil_{+}!}{(t \ln 2)^{\lceil s \rceil_{+}}}\;.
	\]
	With $s = a + b$, $t = b$ and $x = m$ the right-hand side is the constant $C$ above, so
	\[
		a(md)^a\, 2^{-bm}
		= a \bigl( m^{a+b}\, 2^{-bm} \bigr) m^{-b}
		\le a\, C\, m^{-b}
		\le a\, C\, (\log R)^{-b}\;,
	\]
	the last step using $m \ge \log R \ge 1$. Summing the three bounds, the statement follows with $a'$ as above and $b' = b$.
\end{proof}

\begin{lemma}[Formalization support for Lemma~\ref{lem:delta-bound}]
\label{lem:delta-bound-envelope-support}
\lean{MIPStarRE.LDT.rpow_mul_two_rpow_neg_le}
\leanok
For real $s$, $t > 0$ and $x \ge 1$, writing $\lceil s \rceil_{+} = \max(\lceil s \rceil, 0)$,
\[
  x^{s}\, 2^{-t x} \le \frac{\lceil s \rceil_{+}!}{(t \ln 2)^{\lceil s \rceil_{+}}} .
\]
This is the explicit envelope used for the third term in the proof of Lemma~\ref{lem:delta-bound}, with $x = m$, $s = a + b$ and $t = b$.
\end{lemma}
\begin{proof}
\leanok
Since $x \ge 1$, $x^{s} \le x^{\lceil s \rceil_{+}}$, and the Taylor lower bound $y^{n}/n! \le e^{y}$ at $y = t x \ln 2$ and $n = \lceil s \rceil_{+}$ gives $x^{n} (t \ln 2)^{n} \le n!\, 2^{t x}$.
\end{proof}

\begin{remark}[A comparison of exponents in the source proof of Lemma~\ref{lem:delta-bound}]\label{rem:delta-bound-exponent-comparison}
	The source proof bounds the third term of $\delta_{\mathrm{qld}}$ by asserting $2^{-bmd} \le q^{-b}$ ``because $k \le m$'', comparing the field exponent $k = c \lceil \log\log R \rceil + 1$ with $m = 2^j$. The inequality $k \le m$ does not follow from Definition~\ref{def:introparams}: if the universal constants of Theorem~\ref{thm:pauli} were $a = 1$ and $b = 1/2$, then $c = 4$, and $R = 5$ gives $k = 9$ and $m = 8$. The proof above instead bounds the third term directly, comparing $m^{a+b} 2^{-bm}$ with an explicit constant and using only $m \ge (c/2) \log R \ge \log R$, which yields the same conclusion after enlarging the universal constant $a'$.  The comparison and correction are documented in \cite{gap:qpbt_delta-bound-exponent-comparison}.
\end{remark}

\section{Formalization support for finite sampling}

The following finite-distribution identities are not named statements of the
source article.  They make explicit the push-forward calculations used in the
low-degree question and line-point laws above.

\begin{lemma}[Extensionality of finite distributions]\label{lem:distribution-ext-support}
	\lean{MIPStarRE.QPBT.Distribution.ext_of_support_of_weight}
	\leanok
	Two finite distributions with the same support and the same weight at every
	point are equal.
\end{lemma}
\begin{proof}
	\leanok
	The two laws have the same finite support and assign the same mass to every
	point, so they coincide.
\end{proof}

\begin{lemma}[Composition of push-forwards]\label{lem:distribution-map-comp}
	\lean{MIPStarRE.QPBT.Distribution.map_map}
	\leanok
	Let $\mu$ be a finite distribution on $A$, and let $e\colon A\to B$ and
	$f\colon B\to C$. Then
	\[
		f_\ast(e_\ast\mu)=(f\circ e)_\ast\mu.
	\]
\end{lemma}
\begin{proof}
	\leanok
	\uses{lem:distribution-ext-support}
	The supports agree by composition of finite images.  For each $c\in C$,
	partition the fiber of $f\circ e$ over $c$ by the intermediate value $e(a)$;
	the corresponding sums of weights agree.
\end{proof}

\begin{lemma}[Weights of the uniform law]\label{lem:uniform-distribution-weight}
	\lean{MIPStarRE.QPBT.uniformDistribution_weight_apply}
	\leanok
	If $A$ is a nonempty finite set, then its uniform distribution assigns weight
	$1/|A|$ to every $a\in A$.
\end{lemma}
\begin{proof}
	\leanok
	This is the defining weight of the uniform distribution on a nonempty finite
	set.
\end{proof}

\begin{lemma}[Fibre weights of a uniform push-forward]\label{lem:uniform-map-fibre-weight}
	\lean{MIPStarRE.QPBT.uniformDistribution_map_weight}
	\leanok
	\uses{lem:uniform-distribution-weight}
	Let $A$ be a nonempty finite set and let $e\colon A\to C$. Then for every
	$c\in C$,
	\[
		(e_\ast\operatorname{Unif}(A))(c)=\frac{|e^{-1}(c)|}{|A|}.
	\]
\end{lemma}
\begin{proof}
	\leanok
	\uses{lem:uniform-distribution-weight}
	The push-forward assigns to $c$ the total uniform mass of its fibre, and
	every point of $A$ carries mass $1/|A|$.
\end{proof}

\begin{lemma}[Uniform push-forward along a balanced map]\label{lem:uniform-map-equal-fibers}
	\lean{MIPStarRE.QPBT.uniformDistribution_map_of_card_fiber}
	\leanok
	Let $A$ and $B$ be nonempty finite sets and let $e\colon A\to B$. If there is
	an integer $c$ such that $|e^{-1}(b)|=c$ for every $b\in B$, then
	$e_\ast\operatorname{Unif}(A)=\operatorname{Unif}(B)$.
\end{lemma}
\begin{proof}
	\leanok
	\uses{lem:distribution-ext-support, lem:uniform-distribution-weight}
	Every fiber is nonempty, so $e$ is surjective and both supports are $B$.
	Counting the fibers gives $|A|=c|B|$; summing the weight $1/|A|$ over a
	$c$-element fiber gives $1/|B|$.
\end{proof}

\begin{lemma}[Uniformity under an equivalence]\label{lem:uniform-map-equivalence}
	\lean{MIPStarRE.QPBT.uniformDistribution_map_equiv}
	\leanok
	An equivalence $E\colon A\simeq B$ between nonempty finite sets satisfies
	$E_\ast\operatorname{Unif}(A)=\operatorname{Unif}(B)$.
\end{lemma}
\begin{proof}
	\leanok
	\uses{lem:uniform-map-equal-fibers}
	Every fiber of an equivalence has cardinality one, so the result follows from
	Lemma~\ref{lem:uniform-map-equal-fibers}.
\end{proof}

\begin{lemma}[First marginal of a uniform product]\label{lem:uniform-product-first-marginal}
	\lean{MIPStarRE.QPBT.uniformDistribution_map_fst}
	\leanok
	For nonempty finite sets $A$ and $B$, the first marginal of
	$\operatorname{Unif}(A\times B)$ is $\operatorname{Unif}(A)$.
\end{lemma}
\begin{proof}
	\leanok
	\uses{lem:uniform-map-equal-fibers}
	Every fiber of the first projection has cardinality $|B|$; apply
	Lemma~\ref{lem:uniform-map-equal-fibers}.
\end{proof}

\begin{lemma}[Second marginal of a uniform product]\label{lem:uniform-product-second-marginal}
	\lean{MIPStarRE.QPBT.uniformDistribution_map_snd}
	\leanok
	For nonempty finite sets $A$ and $B$, the second marginal of
	$\operatorname{Unif}(A\times B)$ is $\operatorname{Unif}(B)$.
\end{lemma}
\begin{proof}
	\leanok
	\uses{lem:uniform-map-equal-fibers}
	Every fiber of the second projection has cardinality $|A|$; apply
	Lemma~\ref{lem:uniform-map-equal-fibers}.
\end{proof}

\begin{lemma}[Dependent uniform sampling from a finite subset]
	\label{lem:uniform-bind-on-finset-map}
	\lean{MIPStarRE.QPBT.bind_uniformOnFinset_map}
	\leanok
	Let $S$ be a finite subset of $A$, let $B$ be a nonempty finite set, let
	$g\colon A\to B\to C$, and let $f\colon\Sigma\to A$ be an injection from a
	nonempty finite set $\Sigma$ with image $S$.  Sampling $a$ uniformly from
	$S$, then $b$ uniformly from $B$, and returning $g(a,b)$ has the same law as
	pushing $\operatorname{Unif}(\Sigma\times B)$ forward along
	$(x,b)\mapsto g(f(x),b)$.
\end{lemma}
\begin{proof}
	\leanok
	\uses{lem:distribution-ext-support, lem:uniform-map-fibre-weight}
	Both supports are the image of $S\times B$ under $g$.  Since $f$ is a
	bijection from $\Sigma$ onto $S$, we have $|S|=|\Sigma|$, and sums over $S$
	may be reindexed by $\Sigma$.  At each $c\in C$, both weights are the number
	of pairs $(x,b)\in\Sigma\times B$ with $g(f(x),b)=c$, divided by
	$|\Sigma|\,|B|$.
\end{proof}

\begin{lemma}[Dependent uniform sampling as a product push-forward]\label{lem:uniform-bind-map}
	\lean{MIPStarRE.QPBT.bind_uniformDistribution_map}
	\leanok
	Let $A$ and $B$ be nonempty finite sets and let $g\colon A\to B\to C$.
	Sampling $a$ uniformly from $A$, then $b$ uniformly from $B$, and returning
	$g(a,b)$ has the same law as pushing $\operatorname{Unif}(A\times B)$ forward
	along $(a,b)\mapsto g(a,b)$.
\end{lemma}
\begin{proof}
	\leanok
	\uses{lem:uniform-bind-on-finset-map}
	This is the case $S=A$, $\Sigma=A$, and $f$ the identity of
	Lemma~\ref{lem:uniform-bind-on-finset-map}.
\end{proof}

\begin{lemma}[Uniform first marginal after an equivalence]
	\label{lem:uniform-first-marginal-after-equivalence}
	\lean{MIPStarRE.QPBT.uniformDistribution_map_fst_of_equiv}
	\leanok
	Let $A$, $B$, and $C$ be nonempty finite sets, let $E\colon A\simeq B\times C$,
	and suppose $f(a)$ is the first component of $E(a)$ for every $a\in A$. Then
	$f_\ast\operatorname{Unif}(A)=\operatorname{Unif}(B)$.
\end{lemma}
\begin{proof}
	\leanok
	\uses{lem:distribution-map-comp, lem:uniform-map-equivalence, lem:uniform-product-first-marginal}
	Push the uniform law along $E$, compose with the first projection, and apply
	uniformity under the equivalence followed by the first-marginal identity.
\end{proof}

\begin{lemma}[Exact fibers of the coordinate-index map]\label{lem:chi-seed-fibers}
	\lean{MIPStarRE.QPBT.LdParams.seedFiberCard_pos, MIPStarRE.QPBT.LdParams.seedFiberProduct_eq}
	\lean{MIPStarRE.QPBT.seedFiberEquiv, MIPStarRE.QPBT.seedFiberEquiv_fst}
	\lean{MIPStarRE.QPBT.seedOfIndexResidue, MIPStarRE.QPBT.seedFiberEquiv_seedOfIndexResidue, MIPStarRE.QPBT.chiIndex_seedOfIndexResidue}
	\leanok
	\uses{def:ld-game, def:ld-question-distribution, def:binary-representation}
	Write $\chi_0(s)=\chi(s)-1$ for the corresponding zero-based coordinate
	index. Every scalar seed has a unique
	representation as a pair
	\[
		(i,r)\in\{0,\ldots,m-1\}\times\{0,\ldots,q/m-1\},
	\]
	and the first component is $\chi_0(s)$. Conversely, every such pair determines
	a unique scalar seed. In particular, the residue set is nonempty and
	$q=m(q/m)$.
\end{lemma}
\begin{proof}
	\leanok
	\uses{def:binary-representation}
	Use the fixed bijection between $\F_q$ and $\{0,\ldots,q-1\}$, the equality
	$q=m(q/m)$, and quotient--remainder decomposition with divisor $q/m$.
\end{proof}

\begin{lemma}[Uniformity of the coordinate index]\label{lem:chi-index-uniform}
	\lean{MIPStarRE.QPBT.uniformDistribution_map_chiIndex}
	\leanok
	\uses{def:ld-question-distribution}
	For a uniformly random $s\in\F_q$, the zero-based index $\chi_0(s)$ is
	uniformly distributed on $\{0,\ldots,m-1\}$.
\end{lemma}
\begin{proof}
	\leanok
	\uses{lem:chi-seed-fibers, lem:uniform-first-marginal-after-equivalence}
	Under the equivalence of Lemma~\ref{lem:chi-seed-fibers}, $\chi_0$ is the first
	projection. Apply Lemma~\ref{lem:uniform-first-marginal-after-equivalence}.
\end{proof}

\begin{lemma}[Idempotence of the prefix projection]
	\label{lem:prefix-projection-idempotent}
	\lean{MIPStarRE.QPBT.prefixProjection_idempotent}
	\leanok
	\uses{def:ld-question-distribution}
	For every coordinate $i$ and direction $v$,
	$\pi_i(\pi_i(v))=\pi_i(v)$.
\end{lemma}
\begin{proof}
	\leanok
	A coordinate before $i$ is set to zero by both sides, while every coordinate
	at or after $i$ is retained by both sides.
\end{proof}

\begin{lemma}[Idempotence of the canonical representative]
	\label{lem:line-representative-idempotent}
	\lean{MIPStarRE.QPBT.lineRepMap_apply_self}
	\leanok
	\uses{def:line-representative}
	For all $u,v\in\F_q^m$,
	$L^{\mathrm{Ln}}_v(L^{\mathrm{Ln}}_v(u))=L^{\mathrm{Ln}}_v(u)$.
\end{lemma}
\begin{proof}
	\leanok
	The canonical representative map is a projection onto the chosen complement
	of the span of $v$, and every projection is idempotent.
\end{proof}

\begin{definition}[Seed-bearing line descriptions]\label{def:ld-line-description}
	\lean{MIPStarRE.QPBT.LineKind, MIPStarRE.QPBT.LineDesc}
	\lean{MIPStarRE.QPBT.aLineDescOf, MIPStarRE.QPBT.dLineDescOf}
	\leanok
	\uses{def:ld-question-distribution, def:line-representative, lem:line-representative-idempotent}
	An axis-line description records a canonical base point and a scalar seed,
	whose coordinate index selects the axis direction. A diagonal-line
	description records a canonical base point, a scalar seed, and a direction
	whose coordinates preceding the selected index vanish. The axis and diagonal
	CL outputs determine such descriptions by taking their canonical
	representatives.
\end{definition}

\begin{lemma}[Prefix property of a diagonal description]
	\label{lem:line-description-prefix}
	\lean{MIPStarRE.QPBT.LineDesc.diagonal_prefix_zero}
	\leanok
	\uses{def:ld-line-description}
	If a seed-bearing line description is diagonal, then every coordinate of its
	direction preceding the index selected by its seed is zero.
\end{lemma}
\begin{proof}
	\leanok
	This is the prefix-zero property carried by a diagonal line description.
\end{proof}

\section{Explicit verifier cardinality}

The quantitative estimates below use the exact size of the finite set of
question-pair types in the Pauli basis test.  This is a Lean-only numerical
specialization; it does not alter the verifier or the source soundness theorem.

\begin{lemma}[Cardinality of the Pauli verifier edge set]
\label{lem:pauli-edge-card-explicit}
\lean{MIPStarRE.QPBT.pauli_edge_card}
\leanok
\uses{def:pauli-question-distribution}
The type of directed Pauli verifier edges has cardinality
\[
  \lvert \mathcal E_{\mathrm{Pauli}}\rvert=86.
\]
\end{lemma}
\begin{proof}
\leanok
Expand the finite union defining the verifier edges and count its disjoint
question-pair classes.
\end{proof}
